\section{Programs with a fixed nef part and special termination}
\label{sec:programs}
\label{prog:section}

Fix a dimension \(n\), and assume \(\Ggood_d\) for every \(d<n\).
This section supplies two program constructions used in the induction.
An ordinary dlt adjoint that is already the sum of a nef rational line and
its divisorial negative part has a nef model, reached by contracting that
negative part while preserving the nef line. A suitably chosen scaling of
a pseudo-effective ordinary dlt adjoint starting on a smooth space has
special termination. Its proof uses \(\Ggood_d\) only on proper log
canonical strata and compares two descriptions of negative multiplicities
as the scaling parameter tends to zero. Along the way, cohomology transport
supplies both the generic scaling directions and the exceptional splitting
carried by the terminal models used in Section~\ref{sec:boundary}.

The restricted generalized model input is stated first for use in these
ordinary constructions. Its proof occupies the last three subsections:
polarization and descent, an already-projective relative construction,
and the dimension induction. That proof is independent of the assumptions
\(\Ggood_d\) used for ordinary special termination.

\subsection{Program inputs and descent of actual lines}

A normal compact Kähler space is \emph{globally \(\Q\)-factorial} if
every global Weil divisor has a positive Cartier multiple. It is
\emph{globally strongly \(\Q\)-factorial} if every global coherent
rank-one reflexive sheaf has an invertible reflexive power. The strong
property implies the first one by applying it to the sheaf associated to a
Weil divisor. Both properties concern global objects; the strong property
makes no assertion about reflexive sheaves defined only on arbitrary
analytic open subsets. For a rational line \(J\) we write
\(\{J\}=c_1(J)\), and use the same braces for the class of a rational
Cartier divisor.  A trace of a rational line under a bimeromorphic map
means its reflexive transform. Bott--Chern transforms will be used only
along the detected birational steps described below; on first Chern classes they
agree with reflexive transforms whenever the step is ordinary.
All resolutions in the compact arguments below are smooth compact
Kähler spaces, and their indicated maps to the spaces being resolved
are projective.  A common resolution is projective over each
indicated model.

\begin{definition}[A resolution adapted to log canonical strata]
\label{def:lc-strata-resolution}
An effective rational dlt pair \((T,B)\) on a normal compact Kähler space
satisfies the \emph{lc-strata resolution convention} if it admits a
projective log resolution \(r:W\to T\) with smooth compact Kähler source
such that the strict boundary and the exceptional support together have
distinct smooth components forming a simple normal crossing divisor.
Writing
\[
 K_W+B_W^{\mathrm{cr}}=r^*(K_T+B)
\]
with compatible canonical choices, every \(r\)-exceptional coefficient
of \(B_W^{\mathrm{cr}}\) is strictly less than one, and \(r\) is an
isomorphism at the general point of every log canonical center of
\((T,B)\).
\end{definition}

This convention asserts the existence of one such resolution. It imposes
no condition on later resolutions chosen for other purposes. The ordinary
dlt models supplied below for the applications in
Sections~\ref{sec:boundary}, \ref{sec:rank-one}, and~\ref{sec:signed} admit
one by choosing a defining dlt resolution that preserves the simple normal
crossing open set meeting the general points of all log canonical centers,
and resolving any remaining data away from that open set.

An \emph{ordinary negative step} for \(J=K_T+B\) is a projective
bimeromorphic divisorial contraction, or a diagram
\[
 T\xrightarrow{f} Z\xleftarrow{f^+}T^+
\]
of projective small bimeromorphic morphisms, with connected fibers and
normal base, such that \(-J\) is \(f\)-ample and, in the small case,
the trace \(J^+\) is \(f^+\)-ample.  The contracted curves on \(T\)
span one nonzero ray for degrees of global rational lines.  We require
this ray condition on the contracting side only.  All spaces occurring
in these steps are compact Kähler, and the working spaces are globally
strongly \(\Q\)-factorial.  The boundary is pushed forward in a
divisorial step and strictly transformed in a small step.

On a common smooth compact Kähler resolution, projective over \(T\)
and the next space with maps \(p\) and \(q\), the natural meromorphic comparisons of
canonical bundles give the actual rational-line comparison
\begin{equation}
 p^*J\sim_{\Q}q^*J^+ +F,\qquad F\geq0,\qquad q_*F=0.
 \label{prog:comparison}
\end{equation}
Here and below these adjoint identities use the local meromorphic
canonical identifications.  In particular they do not choose a global
meromorphic frame for an arbitrary line bundle.  Negativity proves
Equation~\eqref{prog:comparison}; moreover, its support contains the full
inverse image of the non-isomorphism locus in the contraction base.
Thus discrepancies increase strictly for a place whose center maps
into that locus, and do not decrease elsewhere.  This preserves klt
and dlt singularities.  These assertions, including strictness, are
proved in \cite[Lemma~7.3]{CB}.  Their proof is local over the
contraction base and has no dimension restriction.

For use on strata, we record why the full support assertion holds
for a detected small step. Choose a sufficiently divisible integer
\(r>0\) and the evaluation ideal
\[
 \operatorname{im}\bigl(f^*f_*\OO_T(rJ)\longrightarrow\OO_T(rJ)\bigr)
          =\mathcal I\otimes\OO_T(rJ).
\]
Outside the exceptional locus of \(f\), evaluation is an isomorphism.
On a positive-dimensional projective fiber, every section vanishes:
\(rJ\) has negative degree on every fiber curve, and curves cover
every positive-dimensional fiber component. Connectedness gives the
same vanishing on the whole nontrivial fiber. Hence the zero set of
\(\mathcal I\) is exactly the exceptional locus. On a common
resolution principalizing \(\mathcal I\), relative generation on
the positive side identifies \(rF\) in
Equation~\eqref{prog:comparison} with the divisor of
\(\mathcal I\OO_W\). Thus \(\Supp F\) is the full inverse image
of that locus. On a common resolution of a longer negative program,
discrepancy monotonicity makes the cumulative comparison dominate
this first-step divisor. These are support statements for the actual
comparison, stronger than effectiveness and exceptionality alone.

We next state the restricted model inputs used in this section. A nef
b-class is a Bott--Chern class that is globally nef on one fixed smooth
compact Kähler carrier, with its linear traces on the other models. The
generalized adjoint on \(T\) is
\(\alpha=\{K_T+B\}+\beta_T\). Generalized klt, abbreviated gklt,
means that all generalized log discrepancies, computed on that carrier,
are positive. The modified-bigness hypothesis below concerns the whole
boundary-plus-nef trace \(\{B\}+\beta_T\). We use the convention of
\cite[Definition~2.8]{HX}: a modified-big trace is the pushforward of
a big Bott--Chern class on a modification. The trace itself may be
a current; the generalized adjoint is required to define a
Bott--Chern class.
For a proper map $f:T\to S$, let $q_f$ be the quotient map to
$H^{1,1}_{\BC}(T,\R)/f^*H^{1,1}_{\BC}(S,\R)$. Relative
pseudo-effectivity means
$q_f(\alpha)\in\overline{q_f(\operatorname{Psef}(T))}$, where
$\operatorname{Psef}(T)$ is the absolute pseudo-effective cone.
Equivalently, for Kähler classes $\omega_T,\omega_S$, the class
$\alpha$ is pseudo-effective over $S$ if for every $\epsilon>0$
some $c_\epsilon\geq0$ makes
$\alpha+\epsilon\omega_T+c_\epsilon f^*\omega_S$ big.
This imposes no projectivity hypothesis on $f$. For the actual
relative line classes of an already projective map, this is the
usual relative pseudo-effective cone.

\begin{proposition}[Restricted Kähler model inputs]\label{prog:mori}
Let \((T,B+\boldsymbol\beta)\) be an effective compact Kähler gklt
pair of dimension \(d\), with \(T\) globally strongly
\(\Q\)-factorial. Suppose that its nef b-data are globally nef on a
fixed smooth carrier and that \(\{B\}+\beta_T\) is globally
modified big. Write \(\alpha=\{K_T+B\}+\beta_T\).
\begin{enumerate}
\item If \(\alpha\) is nef, there is a proper morphism with connected
fibers \(g:T\to Z\) to a normal compact Kähler space and a Kähler
class \(\omega_Z\) such that \(\alpha=g^*\omega_Z\).
\item For a proper morphism \(\pi:T\to V\) to a normal compact
Kähler space, if \(\alpha\) is pseudo-effective over \(V\), there
is a chosen good log terminal model over \(V\). It is nonextracting,
compact Kähler and globally strongly \(\Q\)-factorial. Its adjoint
is nef over \(V\), and on a common projective resolution its comparison
with \(\alpha\) is effective and exceptional over the new model,
with strictly positive coefficient at every prime contracted from \(T\).
The construction retains compatible forward Bott--Chern traces of the
specified data on a common smooth carrier.
\item A compact polytope of these data on one fixed carrier has
finitely many marked relative canonical models and relative weak log
canonical models with normal compact Kähler targets.
\end{enumerate}
The relative assertion applies to proper morphisms, without assuming
that \(\pi\) is projective. It is an assertion about a chosen model,
not termination of every generalized flip sequence.
\end{proposition}

\begin{proof}
These are the semiampleness, proper relative model and finite-model
assertions of Theorem~\ref{prog:restricted-package}, proved below
by dimension induction independently of \(\Ggood\).
Finite marked weak-model geography is recorded in
Lemma~\ref{prog:finite-geography}. Projective analytic model
constructions used in that induction are supplied by
Proposition~\ref{prog:projective-good-model}. None of these statements
promotes an undetected generalized-ray contraction to a projective map.
\end{proof}

We use separately the analytic cone theorem
\cite[Theorem~1.3]{HX}: an \(\alpha\)-negative analytic extremal ray
of an effective gklt adjoint has a rational curve generator \(C\) with
\(0<-\alpha\cdot C\leq2d\). Existence of the supporting contraction
in the applications below follows from the nef assertion above.
When an ordinary rational adjoint is negative on that ray,
Lemma~\ref{prog:detected-projectivity} and
Proposition~\ref{prog:detected-flip} give its projective contraction and ordinary
flip. We call these \emph{detected ordinary steps}.

The Bott--Chern transform through such a birational step has a concrete
description which does not require the source to be smooth. Write the
step as \(T\xrightarrow{f}Z\xleftarrow{f^+}T^+\), with
\(f^+=\id\) for a divisorial contraction, and let \(J\) be its
negative rational adjoint. For \(\gamma\in H^{1,1}_{\BC}(T,\R)\),
choose the unique real number \(c\) for which
\(\gamma-c\{J\}\) annihilates the contracted analytic ray.
Every contracted curve is on that ray. Lemma~\ref{prog:bc-descent}
therefore gives a unique class \(\delta\) on \(Z\) with
\(\gamma-c\{J\}=f^*\delta\), and we put
\[
 \gamma^+=(f^+)^*\delta+c\{J^+\}.
\]
The rational-singularity and vanishing hypotheses of that lemma hold:
slightly lowering the rational Cartier floor gives a klt adjoint still
antiample over \(Z\), and relative vanishing gives rational
singularities on the base. On a common resolution, the pullback
difference is \(c\) times the adjoint comparison in
Equation~\eqref{prog:comparison}; it is exceptional over the new
model, and over both models in a flip. This defines a linear transform.
For the Chern class of a rational line it agrees with its reflexive trace,
by the coherent line comparison \cite[Lemma~7.2]{CB} and exceptional
negativity. The already-projective relative program below only needs a
ray for degrees of global rational lines; this Bott--Chern construction
is used for its subsequent detected analytic-ray programs.

When such a class difference lies in an exceptional divisor span, we write
it as \(\{F\}\) for the representing real exceptional divisor \(F\).
This divisor is unique: an exceptional divisor with zero class has zero
degree on every contracted curve, and exceptional negativity applied to
both signs makes it zero.

We also use local projective analytic results.  For a projective
analytic morphism over a Stein neighborhood of a compactum with
Fujino's property (P), relative klt base point freeness says that, if
\(L\) is a relatively nef Cartier line and
\(aL-(K+\Gamma)\) is relatively ample for some positive integer
\(a\), then all sufficiently high powers of \(L\) are relatively
generated after shrinking
\cite[Theorems~6.2 and 6.5]{FujinoAnalyticBCHM}.
We use the finite negative-ray truncation in the relative cone theorem
in the same setting.  Finally, the multigraded adjoint finite-generation
theorem applies to a projective morphism from a smooth space, for simultaneous
effective SNC klt boundaries with a common relatively ample rational
summand \cite[Theorem~3.1]{DHP}.  Finitely many smaller Stein
neighborhoods suffice over a compact base.

\begin{lemma}[Descent without changing a Cartier exponent]
\label{prog:integral-descent}
Let \(f:T\to Z\) be a projective bimeromorphic contraction with normal
target.  Suppose an ordinary rational klt adjoint is \(f\)-antiample.
If a Cartier line \(L\) has degree zero on every contracted curve,
then \(f_*L\) is an invertible sheaf and evaluation is an isomorphism
\[
 f^*(f_*L)\simeq L.
\]
The same conclusion holds for an ordinary dlt adjoint whose floor is
rational Cartier and which is \(f\)-antiample.  In a flip, a line with
zero contracted degree therefore has the same Cartier exponent on
both sides, and the two lines are pullbacks of a line on the common
base.
\end{lemma}

\begin{proof}
For dlt input, decrease the floor by a sufficiently small positive
rational multiple.  The resulting pair is klt and remains
\(f\)-antiample.  For this klt pair the base point free hypothesis
holds for \(L\), because \(L\) has zero fiber degrees and the negative
adjoint is relatively ample.  Over a smaller base neighborhood every
sufficiently high power of \(L\) is therefore generated.  The
associated morphism is constant on each connected fiber: its
tautological line has degree zero on every fiber curve, and every
positive-dimensional projective image contains a curve.  Its graph
then factors through \(Z\); the graph projection is finite
bimeromorphic and \(Z\) is normal. Descend two consecutive generated
powers \(L^k,L^{k+1}\) to lines \(M_k,M_{k+1}\). The line
\(M_{k+1}\otimes M_k^{-1}\) pulls back to \(L\). Projection formula
identifies this quotient with
\(f_*L\), so these local descents and their evaluation maps agree on
overlaps.  Pull the descended line to the positive side of a flip.
This is the argument of \cite[Lemma~7.4]{CB}.
\end{proof}

\subsection{Relative klt programs and small models}

The local finite-generation theorem also supplies an ordinary klt program
over a compact base. We use it to obtain globally strongly
\(\Q\)-factorial small models of strata. The construction keeps track of
degrees of global rational lines, which are the degrees needed for exact
Cartier descent.

\begin{proposition}[Relative klt programs over a compact base]
\label{prog:relative-klt}
Let \(\pi:T\to V\) be projective bimeromorphic, where \(V\) is a
normal compact Kähler space and \(T\) is normal, compact Kähler, and
globally strongly \(\Q\)-factorial.  Let \((T,B)\) be an effective
rational klt pair with rational-line adjoint \(J=K_T+B\).
There is a finite sequence of ordinary \(J\)-negative steps over
\(V\), all projective over \(V\), whose last adjoint is curve-nef
over \(V\), meaning that it has nonnegative degree on every curve
contracted over \(V\). Every working space is compact Kähler and globally
strongly \(\Q\)-factorial.  On a common resolution the comparison
from the first adjoint to the last is effective and exceptional over
the last space.
\end{proposition}

\begin{proof}
We give the reduction to the local finite-generation input, including why
it yields a single finite program over the compact base. Let
\(\mathcal N_{\R}\) be the finite-dimensional space of degrees of global
rational lines on curves over \(V\), and put \(r=\dim\mathcal N_{\R}\).
If \(r=0\), then \(J\) is already curve-nef. Assume \(r>0\). Choose
rational lines \(H_1,\ldots,H_r\) whose degree classes form a basis and
such that both \(H_i\) and \(J+H_i\) are relatively ample. Such a basis
exists because the relatively ample classes whose sum with \(J\) is also
relatively ample form a nonempty open set containing sufficiently positive
classes. Define the rational simplex of adjoints
\[
 \Pi=\operatorname{conv}\{J,J+H_1,\ldots,J+H_r\}.
\]
We will choose a general direction \(H=\sum_i\theta_iH_i\) with
\(\theta_i>0\) and \(\sum_i\theta_i=1\). Its entire scaling segment lies
in this simplex:
\[
 J+tH=(1-t)J+\sum_i t\theta_i(J+H_i)\in\Pi
 \qquad(0\leq t\leq1),
\]
and it lies in the relative interior for \(0<t<1\).

We produce the klt representatives for these exact vertices on each
sufficiently small Stein neighborhood separately. A relative Cartier line
and its inverse have local effective representatives there: their proper
direct images are coherent of generic rank one, so Cartan generation
supplies nonzero local sections. This uses that \(\pi\) is bimeromorphic.
Represent each \(H_i\) by a divided general free effective divisor
\(G_i\sim_{\Q}H_i\), using a sufficiently divisible relatively generated
multiple. These finitely many general divisors preserve klt with \(B\).
Choose a relatively ample integral line \(A\), and local effective
representatives \(A^+\sim A\), \(A^-\sim -A\) by the preceding argument.
For one sufficiently small positive rational \(\epsilon\), all boundaries
\[
 \Delta_0=B+\epsilon(A^++A^-),\qquad
 \Delta_i=B+G_i+\epsilon(A^++A^-)
\]
are effective and klt. Their adjoints represent \(J\) and \(J+H_i\),
respectively, and \(\epsilon A^+\) is a common effective relatively ample
summand. In particular, every rational adjoint in \(\Pi\) is relatively big.

On this Stein member, choose one simultaneous projective log resolution
of the finite family. Give each exceptional prime a coefficient in
\((\max\{0,c_0,\ldots,c_r\},1)\), where the \(c_j\) are its crepant
coefficients for the vertex pairs. If
\(E\geq0\) is resolution-exceptional with \(-E\) resolution-ample, then
subtracting a sufficiently small multiple of \(E\) from the pullback of
\(\epsilon A^+\) leaves a relatively ample rational divisor. A sufficiently
small positive rational multiple of this divisor can be removed from every
boundary while leaving all coefficients effective and below one; the strict
exceptional margins ensure this along the exceptional primes. A divided
general free representative of this common summand preserves the simultaneous SNC
condition. The smooth multigraded theorem now applies to these vertices.
Exceptional corrections leave the local adjoint rings unchanged by
projection to a normal space; clearing the finitely many line
identifications simultaneously makes them multiplicative. This is the
reduction in \cite[Section~7.2]{CB}, using \cite[Theorem~3.1]{DHP}.

Here are the two consequences of finite generation that we need.
The normalized main relative Proj of any rational adjoint \(\Theta\in\Pi\)
extracts no prime divisor, and its trace is a relatively
ample rational line.  Indeed, resolve the degree-one base ideal after
a common Veronese and write
\[
 p^*(m\Theta)\sim q^*H_\Theta+G,\qquad G\geq0,
\]
with \(H_\Theta\) tautological and relatively ample. Generation says that
all sections in degree \(k\) vanish at least along \(kG\).
If a component of \(G\) were not \(q\)-exceptional, relative
generation of \(q_*\OO(G)\otimes H_\Theta^k\) would supply a section with
smaller vanishing there. If a \(p\)-exceptional prime \(Q\) were not
\(q\)-exceptional, the same argument with \(q_*\OO(Q)\otimes H_\Theta^k\)
would contradict \(p_*\OO(Q)=\OO_T\). This proves nonextraction.
There are also only finitely many \emph{marked} normalized main Proj
models, where the marking records their common bimeromorphic open
over \(V\).  To see this, choose finitely many homogeneous generators
locally.  Proj charts of each diagonal ring use homogeneous monomials
as denominators; the supports of those monomials in the finite set of
generators determine the degree-zero localizations and their gluing.
Only finitely many support patterns occur.  Normalization and taking
the main component preserve this finiteness.  A finite cover of the
compact base gives finitely many global marked models, since marked
local identifications agree on the common dense open and glue
uniquely.  These are \cite[Lemmas~7.5 and 7.6]{CB}.

For completeness, we construct the scaling program controlled by this
finite list. On a working model, let \(A_0\) be a relatively ample line.
The closed curve cone in the dual of its global degree space has compact
slice \(A_0\cdot z=1\): for every line \(D\), both \(kA_0+D\) and \(kA_0-D\)
are relatively ample for \(k\) large, which bounds every coordinate.
Local cone truncations on the finite Stein cover express the negative
part of this slice, after any positive ample truncation, using
finitely many actual curve classes.  A negative extremal ray can
therefore be separated by a rational nef support annihilating just
that ray.  A positive multiple of the support minus the klt adjoint
is relatively ample.  Relative base point freeness contracts exactly
the ray.  If the contraction is divisorial, its single exceptional
prime and Lemma~\ref{prog:integral-descent} prove the global strong
property downstairs.  If it is small, the relative Proj of the
adjoint is its small positive model.  For any global line \(D\) on
the negative model, killing its ray degree and applying the same
lemma gives an actual rational identity
\[
 D\sim_{\Q}cJ+f^*D_Z,\qquad c\in\Q.
\]
It transforms to the positive side, proving the global strong
property there.  This is the continuation construction of
\cite[Proposition~7.7]{CB}; it also preserves projectivity over \(V\)
and compact Kählerness.

The transforms of \(H_1,\ldots,H_r\) continue to span the degree spaces:
use pullback in a divisorial step and the preceding decomposition of
each line in a small step.  A line of zero global relative degrees
stays so by Lemma~\ref{prog:integral-descent}; curves in the
contraction bases can be lifted through projective morphisms.
There are only countably many possible finite sequences of curve rays,
and each determines its contractions and flips uniquely. We may therefore
choose \(H\) in the relative interior of
\(\operatorname{conv}\{H_1,\ldots,H_r\}\) so that it avoids all ties
between independent ray degrees in advance.
At a positive nef wall, write \(w=J+tH\) for the traces on the current
model, and choose a preceding parameter \(s>t\) for which \(J+sH\) is
relatively ample. On the zero face of the normalized compact curve slice,
\[
 J=-\frac{t}{s-t}(J+sH)
\]
in degrees. Thus \(J\) is uniformly negative on that face. Choose the
positive ample truncation so that its compact remainder \(K\) misses the
face. The finite rational polyhedral test and the general choice of \(H\)
make the face a single ray \(R\). The wall class has a positive minimum on
\(K\) and is positive on every other truncated curve generator. Hence a
sufficiently small open neighborhood of \(w\) inside the annihilator
\(R^\perp\) consists of nef classes with zero face exactly \(R\).

Suppose first that \(R^\perp\) has positive dimension. It is a rational
hyperplane because \(R\) is represented by an integral curve. Choose a
rational simplex of support classes in that neighborhood containing the
degree class of \(w\) in its relative interior, and represent each vertex
by a rational line \(N\). For each such \(N\), compactness gives an
integer \(a>0\) for which \(aN-J\) is relatively ample. Relative base
point freeness contracts exactly \(R\) and descends \(N\) to a relatively
ample rational line on the contraction base. These bases agree because
the contracted curves agree. The difference between \(w\) and the
corresponding convex combination of these rational lines has zero
relative degree. In the finite rational span of the lines involved, the
degree map has rational kernel, since rational lines have rational degrees
on integral curves. The difference is therefore a real combination of
rational lines of zero relative degree. Lemma~\ref{prog:integral-descent}
descends these lines, and lifting curves from the base shows that their
descents still have zero degree over \(V\). They do not affect relative
ampleness. Convexity now shows that \(w\) descends to a relatively ample
real line class over \(V\). If \(R^\perp=0\), the same zero-degree descent
applies to \(w\); the contraction base has no curves over \(V\), hence is
finite over \(V\), and relative ampleness is automatic.

After a flip, let \(w^+\) be the pullback of this wall class to the positive
side. For sufficiently small \(\delta>0\),
\[
 J^++(t-\delta)H^+
   =\left(1-\frac{\delta}{t}\right)w^+
      +\frac{\delta}{t}J^+
\]
is relatively ample over \(V\): the wall is pulled back from a relatively
ample class on the contraction base, and \(J^+\) is relatively ample over
that base. After a divisorial contraction, openness of the ample cone on
the base gives the same nonempty interval of relative ampleness.

Each working model is consequently one of the marked ample models
already counted. Indeed, for a fixed finite prefix and an interior
parameter \(0<t<1\), replace the direction and parameter by nearby
rational ones. They remain in \(\Pi\), and all strict signs in that prefix
and the final relative ampleness persist; the associated relative section
ring has this working model as its Proj. A repeated marked working model
would have the same trace of \(B\), because all steps are nonextracting,
and hence the same discrepancies. This contradicts the strict increase in
Equation~\eqref{prog:comparison} at an intervening nontrivial step. The
finite list therefore forces termination.
The last adjoint is curve-nef, since the continuation construction
applies whenever it is not.  Composing
Equation~\eqref{prog:comparison} proves the final comparison.
\end{proof}

\begin{corollary}[Small strong models]\label{prog:small-model}
Let \((S,\Delta)\) be an effective rational dlt pair on a normal
compact Kähler space, with rational-line adjoint. Suppose that an effective
rational Cartier divisor \(D\) satisfies
\(\Supp D=\Supp\lfloor\Delta\rfloor\). There is a projective small morphism
\(h:Y\to S\), with \(Y\) globally strongly \(\Q\)-factorial and
compact Kähler, which is crepant for the full adjoint and is an
isomorphism over the smooth locus of \(S\).  The transformed full
pair is dlt.  The same assertion for an ordinary klt pair allows
\(D=0\).
\end{corollary}

\begin{proof}
Choose a small rational \(\epsilon>0\) for which
\((S,\Delta-\epsilon D)\) is effective and klt.  On a projective log
resolution \(r:W\to S\), use the strict lowered boundary and give
every exceptional prime a coefficient strictly between its crepant
coefficient and \(1\), and at least zero.  The resulting klt adjoint
has the actual form
\[
 K_W+\Gamma\sim_{\Q}r^*(K_S+\Delta-\epsilon D)+F,
 \qquad F\geq0,
\]
where every \(r\)-exceptional prime has positive coefficient in
\(F\).  Apply Proposition~\ref{prog:relative-klt} over \(S\).
On its endpoint the trace \(F_Y\) is effective, exceptional over
\(S\), and relatively nef.  Negativity makes \(F_Y=0\).
Since all initially exceptional primes had positive coefficient,
\(Y\to S\) is small and the lowered adjoint is crepant.
Restoring \(\epsilon D\) gives the full crepant identity.

A projective small morphism to a smooth germ is an isomorphism:
transport a relatively ample line to the smooth germ, where it is
Cartier; smallness makes the original line its pullback, contradicting
relative ampleness on a nontrivial fiber.  Thus \(h\) is an isomorphism
over an SNC open meeting all lc centers of the dlt pair.  The full
crepant comparison then proves dlt on \(Y\).
\end{proof}

\subsection{Scaling toward zero and keeping a nef rational line}

The next construction runs every positive truncation of an ordinary
dlt scaling.  It also explains why a sufficiently large multiple of
a prescribed nef rational line forces every step to be trivial for
that line.  The integer that clears the line will be the same at
every step.

For the detected ordinary scalings constructed below, let
\(\alpha_0(\lambda)=\{H\}+\lambda\omega\) on the first space
\(T_0\), with initial parameter \(1\), and let \(\alpha_i(\lambda)\)
be its linear transform on a working model \(T_i\). Put
\(\lambda_{-1}=1\). If the next step occurs at threshold
\(\lambda_i\), its ray is annihilated by \(\alpha_i(\lambda_i)\)
and is negative for \(\alpha_i(0)\). The trace \(\alpha_i(\lambda)\)
is nef on \([\lambda_i,\lambda_{i-1}]\), which we call its
\emph{working interval}. If \(\alpha_i(0)\) is nef, we set
\(\lambda_i=0\) and stop. Working intervals are allowed to be
degenerate; strict decrease of positive thresholds will be arranged only
in the proof of special termination.

\begin{lemma}[Positive truncations and a fixed nef line]
\label{prog:nef-trivial}
Let \((T,B)\) be an effective rational dlt pair on a globally strongly
\(\Q\)-factorial compact Kähler \(d\)-fold, and put \(J=K_T+B\).
\begin{enumerate}
\item For a sufficiently large Kähler class \(\omega\), there is an
ordinary \(J\)-program with scaling of \(\omega\), starting with
\(\{J\}+\omega\) Kähler. Every positive truncation is finite.
If \(J\) is pseudo-effective, the program either reaches a nef adjoint
or has positive thresholds tending to zero. If \(J\) is not
pseudo-effective, the program is finite and ends with a Mori fiber
contraction.
\item Let \(P\) be a nef rational line with \(mP\) Cartier, and choose
\(b\in\Q\) with \(b>4dm\). The same alternatives hold for
\(H=J+bP\), according to pseudo-effectivity of \(H\), with
\(\omega\) enlarged if necessary. Every birational step is an ordinary
\(J\)-negative step on which \(P\) is trivial, and the final Mori
ray, if present, is also \(P\)-trivial. The same line \(mP\) descends
through every birational contraction and pulls back to its positive
side. Its traces remain analytically nef; if \(P\) is semiample, its
same generated multiple and associated morphism are preserved, also
through the final Mori contraction.
\end{enumerate}
\end{lemma}

\begin{proof}
Write \(H=J\) in the first case and \(H=J+bP\) in the second.
Fix \(\omega\) with \(\{H\}+\omega\) Kähler. For each
\(0<\eta<1\), choose a positive rational \(\epsilon\) so small that,
with \(F=\lfloor B\rfloor\),
\[
 (T,B-\epsilon F)\ \text{is klt},\qquad
 \eta\omega+\epsilon\{F\}\ \text{is Kähler}.
\]
For \(\lambda\in[\eta,1]\) use the presentations
\[
 \{H\}+\lambda\omega
  =\{K_T+B-\epsilon F\}
    +\bigl(b\{P\}+\lambda\omega+\epsilon\{F\}\bigr),
\]
omitting \(b\{P\}\) in the first case. Their b-data are globally
nef on the fixed initial carrier, and their boundary-plus-nef traces
are globally modified big. These are effective gklt data.

We first construct a step whenever the unscaled working trace is not
nef. Choose a positive shift strictly below the current wall and use the
gklt presentation of its current trace \(A\). The ray at the wall is
\(A\)-negative. On a compact slice of the analytic curve cone, the
cone theorem makes the \(A\)-negative part locally polyhedral after
a sufficiently small positive Kähler truncation. Separating the
chosen extremal ray from the other finitely many rays in that
truncation gives a Kähler class \(\omega_R\) such that
\(A+\omega_R\) is nef and its zero face is precisely this ray;
see the supporting construction in Proposition~\ref{prog:ordinary-step}.
Add the pullback of this new Kähler class to the fixed nef carrier.
The resulting adjoint still has globally nef data and modified-big
boundary-plus-nef trace, so the nef assertion of
Proposition~\ref{prog:mori} supplies its supporting contraction.
This uses a new supporting direction on the current space; it does
not require the trace of the original scaling direction to be Kähler.
We verify a rational ordinary detector before asserting projectivity
or constructing a flip.
In the second case nefness of \(P\) makes \(J\) negative on the ray;
the first case has this property directly. Lower the floor by a small
rational amount, retaining negativity of \(J_\epsilon\), so that
\(-J\cdot R\leq2(-J_\epsilon\cdot R)\). The ordinary klt cone
length bound gives a rational curve generator \(C\) with
\[
                         0<-J\cdot C\leq4d.
\]
In the second case, if \(P\cdot C>0\), integrality of \(mP\) gives
\(P\cdot C\geq1/m\), contradicting
\[
                         H\cdot C\geq-4d+b/m>0.
\]
Thus in that case the ray is \(P\)-trivial. The lowered ordinary adjoint is a
global rational detector. The detected projectivity and flip
constructions in Lemma~\ref{prog:detected-projectivity}
and Proposition~\ref{prog:detected-flip} supply the projective ordinary step.
The difference of \(J\) and a positive rational multiple of the
lowered adjoint has zero contracted degree; exact descent in
Lemma~\ref{prog:integral-descent} shows that restoring the floor
preserves the positive-side sign. This is the ordinary replacement
of Proposition~\ref{prog:ordinary-replacement}, and the full pair
remains dlt by Equation~\eqref{prog:comparison}.

In the second case, Lemma~\ref{prog:integral-descent} descends the actual Cartier line
\(mP\) without changing \(m\). Analytic nefness descends and pulls
back under these projective contractions \cite[Lemma~2.44]{HX}.
A generated multiple descends as a generated line by projection
formula and pulls back as one. Thus the same length estimate and the
same \(b\) apply at the next step. The corresponding assertion on
a final Mori ray follows from exactly the same estimate. For a
semiample \(P\), the map of its fixed generated multiple is constant
on each connected projective contracted fiber: otherwise its image
contains a curve of positive degree. The map therefore factors through
the normal contraction base, including in the fiber-type case.

The traces used in this construction are linear by the degree-zero
descent description preceding Lemma~\ref{prog:integral-descent}.
At a wall the scaled class is pulled back from its base. The comparison
for an earlier parameter is a nonnegative multiple of the ordinary
adjoint comparison. Consequently the gklt presentations on any fixed
positive segment remain gklt along its program, and every working
model is a weak model for a member of that segment. The finite marked
weak-model assertion in Proposition~\ref{prog:mori} applies.
A marked model cannot recur: its boundary trace would be the same,
whereas a nontrivial intervening ordinary step strictly increases a
discrepancy. There are therefore only finitely many steps with
threshold at least \(\eta\).

For decreasing cutoffs, any already-constructed finite prefix remains
valid. The floor perturbations cancel in the displayed classes and
continue to cancel under their linear transforms; their nonpositive
comparisons preserve the new gklt data. Starting at its last nef wall
therefore extends the prefix to the smaller cutoff. This gives one
compatible scaling program. A positive limiting threshold would lie
in a finite positive truncation, so an infinite program has limit zero.
Pseudo-effectivity excludes a negative Mori endpoint. If \(H\) is
not pseudo-effective, its pseudo-effective threshold in the initial
Kähler direction is positive. Every nef working trace, together with
its effective exceptional comparison, makes the initial scaled class
pseudo-effective, so the scaling thresholds are bounded below by
that positive number. The program is therefore finite. Its endpoint
cannot be nef, since that comparison would make \(H\)
pseudo-effective; it is the stated Mori fiber contraction.
\end{proof}

We call these programs \emph{detected ordinary scalings}. Their steps
are ordinary negative steps of one analytic extremal ray, with the
linear Bott--Chern transform described above.
The next lemma connects their nef working traces
to the analytic negative part, including its limit at parameter zero.

\begin{lemma}[Negative parts along a detected ordinary scaling]
\label{prog:negative-along-scaling}
In a scaling from Lemma~\ref{prog:nef-trivial}, write \(H=J\) in its first
case and \(H=J+bP\) in its second, and use the notation
\(\alpha_0(\lambda)=\{H\}+\lambda\omega\) above. For a parameter
\(\lambda\) in the working interval of \(T_i\), a common smooth
resolution \(p:W\to T_0\), \(q:W\to T_i\) gives
\begin{equation}
 p^*\alpha_0(\lambda)=q^*\alpha_i(\lambda)+\{E_i(\lambda)\},
 \qquad E_i(\lambda)\geq0,\qquad q_*E_i(\lambda)=0.
 \label{prog:working-negative}
\end{equation}
Here \(E_i(\lambda)\) is an actual real exceptional divisor, and
\[
 N(p^*\alpha_0(\lambda))=E_i(\lambda).
\]
For the limit statement, let \(H\) be any rational line on a normal compact
Kähler space \(T\) whose pullback to a smooth resolution is
pseudo-effective, and let \(A\) be a Kähler class on \(T\). On any fixed
smooth resolution \(r:U\to T\) and for every prime \(Q\subset U\),
\[
 \lim_{\lambda\downarrow0}\nu_Q(r^*(\{H\}+\lambda A))
   =\nu_Q(r^*\{H\}).
\]
Thus the same convergence holds for every divisorial place, after choosing
a resolution on which it appears.
\end{lemma}

\begin{proof}
For \(\lambda\) in the working interval of \(T_i\), every earlier
threshold is at least \(\lambda\). Each earlier step is therefore negative
or trivial for this parameter. Composing its comparisons gives
Equation~\eqref{prog:working-negative}. Since \(\alpha_i(\lambda)\) is
nef, Lemma~\ref{bir:exceptional} identifies the exceptional divisor with
the negative part.

On the fixed resolution \(U\), choose a Kähler class \(\Omega\) and
\(c>0\) such that \(c\Omega-r^*A\) is Kähler. Monotonicity of each
minimal multiplicity under addition of a nef class gives
\[
 \nu_Q(r^*\{H\}+\lambda c\Omega)
 \leq \nu_Q(r^*(\{H\}+\lambda A))
 \leq \nu_Q(r^*\{H\}).
\]
The left side tends to the right side by the definition using small Kähler
perturbations. This proves the convergence on \(U\), and the same argument
applies on a resolution representing any chosen divisorial place.
\end{proof}

\begin{proposition}[Contracting a known negative part]
\label{prog:contract-negative}
Let \((T,B)\) be an effective rational dlt pair on a globally
strongly \(\Q\)-factorial compact Kähler space, and put \(J=K_T+B\).
Suppose there is an actual rational-line identity
\[
 J\sim_{\Q}P+E,
\]
where \(P\) is analytically nef, \(E\) is an effective rational
Cartier divisor, and, on one smooth compact Kähler resolution
\(\pi:W\to T\) projective over \(T\),
\[
 N(\pi^*\{J\})=\pi^*E.
\]
There is a finite ordinary \(J\)-negative program to a nef model
\(T'\) which contracts precisely the prime components of \(E\)
among the primes of \(T\).  Every step is \(P\)-trivial, with a
fixed Cartier exponent as in Lemma~\ref{prog:nef-trivial}, and
\[
 J_{T'}\sim_{\Q}P_{T'}.
\]
If \(P\) is semiample, then so is this actual last adjoint.
\end{proposition}

\begin{proof}
Lemma~\ref{bir:pullup} makes the stated equality valid on any
higher resolution.  Choose \(b\) as in
Lemma~\ref{prog:nef-trivial}.  Adding \(bP\) leaves the same
negative part.  Indeed, on a smooth resolution write
\(\alpha=\pi^*\{J\}\), \(p=\pi^*\{P\}\), and \(e=\pi^*E=N(\alpha)\).
Subadditivity and homogeneity of negative multiplicities give
\begin{align*}
 N(\alpha+bp)&\leq e,\\
 (1+b)e=N((1+b)\alpha)
 &\leq N(\alpha+bp)+bN(\{e\})
 \leq N(\alpha+bp)+be.
\end{align*}
Thus \(N(\alpha+bp)=e\).

Run the scaling for \(H=J+bP\).
Let \(Q\) be the strict transform on a fixed resolution of a prime
component of \(E\).  Its multiplicity in \(N(\pi^*\{H\})\) is
positive. The limit in Lemma~\ref{prog:negative-along-scaling} makes its multiplicity
positive for every sufficiently small positive parameter.  If
that prime still survived on the corresponding working model,
Equation~\eqref{prog:working-negative} would make its multiplicity zero,
because the comparison there is exceptional over that model.
It must therefore be contracted after finitely many steps.  The
same conclusion holds if a nef endpoint is reached before taking
the limit.  There are only finitely many components of \(E\), so
after a finite prefix all have disappeared.

Throughout the program \(P\) descends and remains nef.  The
actual identity \(J_i\sim_{\Q}P_i+E_i\), with \(E_i\) the
codimension-one pushforward of \(E\), follows from the reflexive
trace convention. Once \(E_i=0\), we have \(J_i\sim_{\Q}P_i\); since
\(P_i\) is nef, the program stops. Conversely, a divisorial negative step can
contract only a prime of \(E_i\): its ray is \(P_i\)-trivial,
so the effective divisor \(E_i\) has negative degree on every
contracted curve, and those curves are contained in its support.
This proves the exact assertion about contracted primes.
\end{proof}

The next lemma records the cohomology carried by a detected ordinary scaling
that begins on a smooth space. Its surjectivity will let us choose one
generic initial scaling direction for all later models. Its exceptional
splitting will also carry the span of line classes to the terminal models
used in Section~\ref{sec:boundary}.

\begin{lemma}[Cohomology through ordinary steps]
\label{prog:cohomology-transport}
Let a finite prefix of one of the detected ordinary scalings constructed
in Lemma~\ref{prog:nef-trivial} start on a smooth compact Kähler
space, and let \(V\) be any working space.
For every smooth compact Kähler resolution \(p:U\to V\)
projective over \(V\),
\begin{align}
 H^2(U,\R)
 &=p^*H^2(V,\R)\ \oplus\
   \bigoplus_{Q\ {\rm exceptional\ for}\ p}\R\{Q\}, \label{prog:H2}\\
 H^{1,1}(U,\R)
 &=p^*H^{1,1}_{\BC}(V,\R)\ \oplus\
   \bigoplus_{Q\ {\rm exceptional\ for}\ p}\R\{Q\}.
 \label{prog:H11}
\end{align}
The linear transform from the first space is surjective onto
\(H^2(V,\R)\) and \(H^{1,1}_{\BC}(V,\R)\), respectively. For a small correspondence between
working spaces, the difference of the pulled-back classes and
their transforms on a common resolution lies in the span of the
primes exceptional over both spaces.
Under the displayed decompositions, the Chern class of a global
holomorphic line on \(U\) is the sum of the pullback of a rational
holomorphic line class on \(V\) and an exceptional rational divisor
class.  In particular, passage to the exceptional quotient preserves
the real span of line Chern classes.
\end{lemma}

\begin{proof}
For a modification of a smooth compact Kähler manifold, the
degree-two modification theorem gives these decompositions, and
the exceptional classes have type \((1,1)\).  Suppose they hold
at one step.  Its contraction base and both working spaces have
rational singularities: slight lowering of the rational Cartier
floor makes the ordinary pair klt, and the detected contraction has
a base with rational singularities.  These contractions are
bimeromorphic, so a common resolution and Leray give
\(R^if_*\OO=0\) for \(i>0\); this verifies the additional
hypothesis in Lemma~\ref{prog:bc-descent}.

The ray of this detected contraction is one ray for Bott--Chern degrees.
This also suffices for degree-two classes.  Indeed, pull a class
\(\gamma\in H^2(V,\R)\) to a smooth resolution \(p:U\to V\)
and take its Hodge decomposition. By the decomposition already proved for
this working model in Equation~\eqref{prog:H11}, its \((1,1)\)-part is \(p^*\beta+\{E\}\) for a
Bott--Chern class \(\beta\) and an exceptional real divisor
\(E\).  The other Hodge components have zero degree on curves.
On every \(p\)-contracted curve the degree of \(E\) is therefore
zero, so both signs of negativity make \(E=0\).  Lifting any
curve of \(V\) through the projective resolution now shows that
\(\gamma\) and \(\beta\) have the same curve degrees.  Thus
subtracting a multiple of \(\{J\}\) from either a degree-two or
a Bott--Chern class can annihilate every curve of the contracted
fiber.  By
Lemma~\ref{prog:bc-descent}, a degree-two class, or a Bott--Chern class,
with these zero degrees is pulled back from the base.
Subtract a multiple of \(\{J\}\) to make any given class
annihilate that ray, descend the remainder, and use the trace
\(\{J^+\}\) to define its transform.  On a common resolution the
difference is a multiple of the exceptional adjoint comparison in
Equation~\eqref{prog:comparison}.  For a flip the two exceptional prime
sets on that resolution are the same.  For a divisorial
contraction the target resolution has just the one additional
exceptional prime of the step.  The old decomposition therefore
spans the new one and proves surjectivity of the transform.

The sums are direct.  If the pullback of a class from \(V\) is
an exceptional divisor class, that divisor has zero degree on
all curves over \(V\).  Applying exceptional negativity to both
signs makes the divisor zero; injectivity of pullback gives that
the class is zero.  The pullback injectivity in degree two and in
Bott--Chern cohomology is also
Lemma~\ref{prog:bc-descent}, applied to the resolution.  Passing to
higher smooth resolutions proves the assertion for every \(p\).

For the line assertion, let \(\mathscr H\) be a holomorphic line on
\(U\), and put \(\mathscr Q=(p_*\mathscr H)^{**}\).  Global strong
\(\Q\)-factoriality gives an invertible sheaf
\(\mathscr M=\mathscr Q^{[m]}\) for some \(m>0\).  The coherent
evaluation comparison of \cite[Lemma~7.2]{CB} gives
\[
 \mathscr H^{\otimes m}\simeq p^*\mathscr M\otimes\OO_U(E),
 \qquad E\ \text{an integral \(p\)-exceptional divisor}.
\]
Indeed \((p_*\mathscr H^{\otimes m})^{**}=\mathscr M\), since the
two sheaves agree at the general points of all divisors of \(V\).
Local generators of \(p_*\mathscr H^{\otimes m}\) in a frame
of \(\mathscr M\) compare their evaluations with the pulled-back
coefficients by meromorphic quotients.  Those quotients agree on the
isomorphism locus and hence glue; their zeros and poles are exceptional.
Taking Chern classes proves the assertion without choosing a meromorphic
section of \(\mathscr H\).
\end{proof}

Before returning to special termination, we record the consequence used in
the boundary argument: a lower-dimensional manifold of Kodaira dimension
zero has a terminal torsion model carrying the preceding exceptional
splitting.

\begin{corollary}[Terminal models in Kodaira dimension zero]
\label{prog:terminal-kappa-zero}
Let \(X\) be a smooth compact Kähler \(d\)-fold with \(d<n\)
and \(\kappa(X,K_X)=0\).  A finite ordinary \(K_X\)-negative
program reaches a globally strongly \(\Q\)-factorial compact
Kähler terminal space \(X_{\min}\) whose actual canonical
rational line is torsion.  For every smooth compact Kähler
resolution projective over \(X_{\min}\), the decompositions in
Equations~\eqref{prog:H2} and~\eqref{prog:H11} hold.
\end{corollary}

\begin{proof}
A pluricanonical section makes \(K_X\) pseudo-effective.
By \(\Ggood_d\), on a smooth modification
\(\mu:\widetilde X\to X\) there is an actual identity
\[
 \mu^*K_X\sim_{\Q}P+E,\qquad P\ \text{semiample},\quad
 E=N(\mu^*\{K_X\}).
\]
Lemma~\ref{bir:sections} identifies the divisible section spaces
with those of \(P\).  Hence \(\kappa(P)=0\); a semiample line of
Iitaka dimension zero has a trivial positive multiple.
Proposition~\ref{bir:finite-descent}, applied with \(e=\mu\) and \(F=0\),
gives an effective rational divisor \(E_X=N(K_X)\) with
\(K_X\sim_{\Q}E_X\).
Lemma~\ref{bir:pullup} then gives
\(N(\mu^*\{K_X\})=\mu^*E_X\).
Apply Proposition~\ref{prog:contract-negative} with \(P=0\).
It reaches a model with torsion canonical line.  It is terminal:
exceptional places over the initial smooth space have log
discrepancy greater than \(1\), and a divisor contracted from
that space starts with log discrepancy \(1\) and acquires a
strict increase in Equation~\eqref{prog:comparison}.  Finally apply
Lemma~\ref{prog:cohomology-transport}.
\end{proof}

\subsection{A comparison model for the restricted scaling}

We now prepare the lower-dimensional argument that excludes an
infinite sequence of small transformations on a log canonical
stratum.  We first construct a semiample small model when the
negative multiplicities have centers away from the floor.  We then
make one model nef on an entire interval of perturbations.  This
interval is the precise lower-dimensional conclusion needed below.

\begin{lemma}[A semiample small comparison model]
\label{prog:stratum-model}
Let \((S,\Delta)\) be an effective rational dlt pair on a normal
compact Kähler space of dimension \(d<n\).  Suppose
\(J=K_S+\Delta\) is a pseudo-effective rational line, and that an
effective rational Cartier divisor \(D\) satisfies
\(\Supp D=\Supp\lfloor\Delta\rfloor\). Assume for every divisorial
place \(Q\) over \(S\) that
\begin{equation}
 \nu_Q(J)=0
 \quad\text{if the center of \(Q\) is a prime of \(S\), or meets
 \(\Supp\lfloor\Delta\rfloor\).}
 \label{prog:away-floor}
\end{equation}
There is a globally strongly \(\Q\)-factorial compact Kähler dlt
pair \((M,\Delta_M)\), small bimeromorphic with \((S,\Delta)\),
whose adjoint \(J_M\) is an actual semiample rational line.
Over a neighborhood \(U\) of the floor in \(S\), the comparison is
a projective small crepant morphism \(M_U\to U\).

For every \(\gamma\in H^{1,1}_{\BC}(S,\R)\) there is a transformed
class \(\gamma_M\in H^{1,1}_{\BC}(M,\R)\) which is its pullback over
\(U\). On a common smooth resolution \(p:W\to S,\ q:W\to M\),
\[
 q^*\gamma_M-p^*\gamma=\{E_\gamma\},
\]
where \(E_\gamma\) is a real divisor supported on primes exceptional over
both spaces, and \(\Supp E_\gamma\) is disjoint from \(p^{-1}(U)\).
\end{lemma}

\begin{proof}
Resolve the dlt pair, giving new exceptional primes coefficient
one.  The resolved adjoint is the pullback of \(J\) plus an
effective exceptional divisor.  Apply \(\Ggood_d\) to this
smooth pair.  Lemma~\ref{bir:exceptional} subtracts the
exceptional resolution error from its negative part.  On a
smooth higher model \(r:\widetilde S\to S\) we obtain the actual
identity
\begin{equation}
 r^*J\sim_{\Q}P+E,\qquad P\ \text{semiample},\qquad
 E=N(r^*\{J\}).
 \label{prog:stratum-decomposition}
\end{equation}
Every component of \(E\) is exceptional over \(S\) and has
center disjoint from the floor, by Equation~\eqref{prog:away-floor}.
Choose \(m\) divisible enough that \(mJ,mP,mE\) are integral and
\(mP\) is generated.  Lemma~\ref{bir:sections} shows that
multiplication by the section of \(mE\) identifies the complete
section spaces.  Over a neighborhood of the floor, \(E\) is
absent upstairs.  The line \(mr^*J\) is generated there, and
these sections descend to \(S\) by normality.  Thus \(mJ\)
is generated on a neighborhood \(U\) of the entire floor.

Let \(\rho_\Gamma:\Gamma\to S\) be the normalized graph of the map defined by
the complete system \(|mJ|\), and write
\(\varphi:\Gamma\to\PP(H^0(S,mJ)^*)\) for its morphism to projective
space. The map \(\rho_\Gamma\) is projective and is an isomorphism over \(U\).
The rational line
\[
 P_\Gamma=\frac1m\varphi^*\OO(1)
\]
is semiample. The generated moving system \(|mP|\) induces a morphism
\(\psi:\widetilde S\to\Gamma\) whose pullback of \(P_\Gamma\) is \(P\).
Choose a small positive rational \(\epsilon\)
for which \((S,\Delta-\epsilon D)\) is effective and klt.
On a projective log resolution \(W_0\to\Gamma\to S\), give
every exceptional prime a nonnegative coefficient strictly
above its crepant coefficient for the lowered pair and strictly
below \(1\).  Its klt adjoint equals the pullback of
\(J-\epsilon D\) plus an effective error positive on every
prime exceptional over \(S\).

Run Proposition~\ref{prog:relative-klt} over \(\Gamma\), and write
\(g:Y\to\Gamma\) for the endpoint morphism and \(h=\rho_\Gamma g:Y\to S\).
The surviving
error \(F_Y\) is effective and has positive coefficient on
every prime of \(Y\) exceptional over \(S\).  Over \(U\), where
\(\Gamma=S\), it is also exceptional and relatively nef, since
the lowered adjoint is relatively nef and its other term is
pulled back from \(S\).  Local exceptional negativity makes it
zero there.  Hence \(h\) is small over \(U\), and it is an
isomorphism over the smooth part of \(U\) by the argument in
Corollary~\ref{prog:small-model}.  Restore the floor.  The full
pair on \(Y\) is dlt near it, by the crepant small comparison
and the SNC open, and is klt elsewhere.  No exceptional prime
has center meeting the floor, so \(h^*D\) is its strict
transform and
\[
 J_Y\sim_{\Q}h^*J+F_Y.
\]

Put \(P_Y=g^*P_\Gamma\). On a common smooth model
\(a:\widehat S\to\widetilde S\), \(b:\widehat S\to Y\) over \(S\),
the maps \(\psi a\) and \(g b\) agree on the common dense open and hence
everywhere by separatedness of \(\Gamma\). Thus \(a^*P\sim_{\Q}b^*P_Y\).
Define \(E_Y^0=b_*a^*E\). This is an effective rational divisor,
exceptional over \(S\) and supported away from the floor. Pushing
Equation~\eqref{prog:stratum-decomposition} to \(Y\) gives
\[
 h^*J\sim_{\Q}P_Y+E_Y^0.
\]
The global strong property on \(Y\) makes \(E_Y^0\) rational Cartier.
The divisor \(a^*E-b^*E_Y^0\) is \(b\)-exceptional and
rationally linearly trivial. Exceptional negativity applied to both signs therefore
gives \(a^*E=b^*E_Y^0\).

Define \(E_Y=E_Y^0+F_Y\). Adding the restored-boundary error gives
\begin{equation}
 J_Y\sim_{\Q}P_Y+E_Y,\qquad E_Y\geq0,
 \label{prog:graph-adjoint}
\end{equation}
where \(E_Y\) is exceptional over \(S\), is supported away from the floor,
and contains every prime exceptional for \(h\). It is rational Cartier by
the global strong property. On \(\widehat S\), Lemma~\ref{bir:pullup}
gives
\[
 N(a^*r^*\{J\})=a^*E=b^*E_Y^0.
\]
Exceptional translation for \(b^*F_Y\) over \(S\) then gives
\(N(b^*\{J_Y\})=b^*E_Y\). The hypotheses of
Proposition~\ref{prog:contract-negative} are now satisfied.
Each contracted curve lies in \(\Supp E_Y\), since its
\(P_Y\)-degree is zero and its \(E_Y\)-degree is negative.
Nontrivial connected projective fibers are covered by curves,
so this contraction and its positive replacement are unchanged
over a neighborhood of the floor.  The same assertion persists
with the pushed-forward effective divisor at each step.
All primes exceptional over \(S\) are contracted and no prime
coming from \(S\) is contracted.  The endpoint \(M\) is
therefore small with \(S\), has the asserted morphism over
\(U\), and has \(J_M\sim_{\Q}P_M\) semiample.

Pull \(\gamma\) from \(S\) to \(Y\). Although \(Y\) may be singular,
its subsequent program consists of the detected ordinary analytic-ray
steps in Proposition~\ref{prog:contract-negative}. At step \(i\),
choose \(c_i\in\R\) so that
\(\gamma_i-c_i\{J_i\}\) has degree zero on its contracted ray.
All curves in a fiber have class on that same analytic ray. The
working spaces have rational singularities after slightly lowering
the Cartier floor, and relative vanishing gives rational singularities
on the contraction base. Lemma~\ref{prog:bc-descent} therefore
descends this class uniquely to the base. Pull it to the positive side
and add \(c_i\{J_{i+1}\}\), defining \(\gamma_{i+1}\).
The pullback difference is \(c_i\) times the ordinary adjoint
comparison, hence is an actual real exceptional divisor class.
No smooth-start cohomology splitting is needed for this construction.

Telescope these comparisons on a common resolution. The original
map \(Y\to S\) was a morphism, and \(\gamma\) was pulled back
through it. Since \(S\) and \(M\) match in codimension one, the
resulting divisor is exceptional over both. Each step is an
isomorphism over the chosen neighborhood \(U\) of the floor.
There the class remains the pullback from \(S\); the corresponding
comparison divisor has zero class and is exceptional. Negativity
applied to both signs makes it zero there. Thus its support is
disjoint from the inverse image of \(U\), as claimed.
\end{proof}

\begin{lemma}[One nef model for an interval]\label{prog:nef-interval}
Under the hypotheses of Lemma~\ref{prog:stratum-model}, let \(A\) be any Kähler
class on \(S\).  There are a normal globally strongly
\(\Q\)-factorial compact Kähler space \(V\), small bimeromorphic
with \(S\), a number \(\delta>0\), and transformed classes
\(\{J_V\},A_V\) such that
\begin{equation}
 \{J_V\}+tA_V\ \text{is nef for every }0\leq t\leq\delta.
 \label{prog:whole-nef-interval}
\end{equation}
On a common smooth resolution of \(S\) and \(V\), the
differences of pullbacks of \(J,A\) and these transforms lie
in the common exceptional span.
\end{lemma}

\begin{proof}
Let \(P_M=J_M\) denote the semiample rational line from
Lemma~\ref{prog:stratum-model}, and use its transform \(A_M\).
On a common resolution \(p:W\to S,\ q:W\to M\), let \(F_A\) be the
real divisor in the common exceptional span defined by
\[
 q^*A_M-p^*A=\{F_A\}.
\]
It is effective: \(-F_A\) is \(q\)-nef, because \(p^*A\) is
nef and the other term is pulled back from \(M\); apply
exceptional negativity. Its support is disjoint from \(p^{-1}(U)\).
In particular \(A_M\) is pseudo-effective, and its negative
multiplicity at every prime of \(M\) is zero.  Indeed,
\(q^*A_M=p^*A+\{F_A\}\), and
Lemma~\ref{bir:exceptional} applied over \(S\) gives
\(N(q^*A_M)=F_A\).

Let \(D_M\) be the rational Cartier trace of \(D\).
Choose \(\tau>0\) so small that \(A+\tau\{D\}\) is Kähler.
For a fixed sufficiently small \(t>0\), consider on \(M\) the
boundary \(\Delta_M-t\tau D_M\) and the nef b-class represented on
\(W\) by \(t p^*(A+\tau\{D\})\). Its trace on \(M\) is
\(t(A_M+\tau\{D_M\})\), so its generalized adjoint is
\begin{equation}
 \{P_M\}+tA_M.
 \label{prog:interval-gadjoint}
\end{equation}
These are effective gklt data.  Over \(U\) the b-class
descends to \(M\), so its generalized discrepancy
contribution is zero there, and subtracting \(t\tau D_M\)
raises every zero log discrepancy of the dlt pair.
Away from the floor the ordinary pair is klt; on one fixed
log resolution its finitely many subunit coefficients stay
subunit for small \(t\).  This proves gklt.  The generalized
boundary is big.  On a common resolution the pullback of the
trace of \(A+\tau\{D\}\) minus \(p^*(A+\tau\{D\})\) is effective exceptional
by the same negativity argument as for \(F_A\).  The latter
pullback is nef and big, so the trace is big; the remaining
boundary is effective.

Choose an integer \(m>0\) such that \(mP_M\) is generated, and
let \(g:M\to Z\) be the Stein factor of its morphism. Then \(Z\)
is normal projective and \(mP_M=g^*H\) for an ample generated Cartier line
\(H\) on \(Z\). Set \(L=\{P_M\}+tA_M\). The preceding
preparation makes \(L\) an effective gklt adjoint with globally nef
b-data and globally modified-big boundary-plus-nef trace. It is
pseudo-effective because \(P_M\) is nef and \(A_M\) is modified
nef. Apply the proper relative assertion of
Proposition~\ref{prog:mori} over \(Z\). This produces a chosen
nonextracting good log terminal model \(\phi:M\dashrightarrow V\)
and a morphism \(g_V:V\to Z\), with \(L_V\) nef over \(Z\).
The hypothesis here is properness: the semiample morphism \(M\to Z\)
need not be projective. Put
\[
             mP_V=g_V^*H,\qquad A_V=(L_V-\{P_V\})/t.
\]
The relative construction preserves the actual line pulled back from
\(Z\), and its forward Bott--Chern transport preserves the displayed
linear relation; these are the transforms of \(P_M,A_M\).

The map \(\phi\) is small. On a common projective resolution
\(a:W'\to M\), \(c:W'\to V\), its comparison is
\[
             a^*L=c^*L_V+\{E\},\qquad E\geq0,
\]
where \(E\) is exceptional over \(V\) and has positive coefficient
at the strict transform of each prime contracted from \(M\).
The forward trace \(L_V\) is pseudo-effective, so exceptional
translation gives
\[
                  N(a^*L)=N(c^*L_V)+E\geq E.
\]
But \(L\) is modified nef on \(M\); the coefficient of
\(N(a^*L)\) at a strict prime of \(M\) is zero. Thus no prime of
\(M\) is contracted. Nonextraction then proves smallness.

Fix \(b>2dm\). We claim that \(L_V+b\{P_V\}\) is absolutely
nef. Otherwise the cone theorem for the effective gklt adjoint
\(L_V\) gives an \(L_V\)-negative rational extremal generator
\(C\) such that
\[
 (L_V+b\{P_V\})\cdot C<0,\qquad 0<-L_V\cdot C\leq2d.
\]
Indeed the cone summand on which \(L_V\) is nonnegative also has
nonnegative \(P_V\)-degree, so a negative ray summand must account
for any failure of nefness. If \(P_V\cdot C=0\), ampleness of
\(H\) makes \(C\) vertical over \(Z\), contradicting relative
nefness of \(L_V\). Otherwise Cartier integrality gives
\(P_V\cdot C\geq1/m\), contradicting
\[
 (L_V+b\{P_V\})\cdot C\geq-2d+b/m>0.
\]
This proves the claim. Both \(\{P_V\}\) and
\((1+b)\{P_V\}+tA_V\) are nef. By convexity,
\[
 \{P_V\}+sA_V\ \text{is nef for }0\leq s\leq\frac{t}{1+b}.
\]
Thus \(\delta=t/(1+b)\) works. The actual line \(P_V\) has the
same generated multiple as \(P_M\), since both are pulled back from
\(H/m\). Its class is the transform of \(J\).
The relative comparison and linear transport give exceptional-divisor
comparisons for \(P_M,A_M\). Since \(S,M,V\) agree in codimension
one, their final comparison divisors are exceptional over both \(S\)
and \(V\), as required.
\end{proof}

\subsection{Special termination}

\begin{theorem}[Special termination for a chosen dlt scaling]
\label{prog:special-termination}
Assume \(\Ggood_d\) for every \(d<n\).  Let \((X,B)\) be an
effective rational dlt pair on a smooth compact Kähler
\(n\)-fold, and suppose \(L=K_X+B\) is pseudo-effective.
There is a Kähler scaling direction \(\omega\), with
\(\{L\}+\omega\) Kähler, and a scaling as in
Lemma~\ref{prog:nef-trivial}, such that after finitely many
steps the exceptional, flipping, and flipped loci are disjoint
from the reduced floors of the working pairs.
\end{theorem}

\begin{proof}
We first choose a scaling with strict positive walls and Kähler interior
classes. We then reduce the transformations on each log canonical stratum
to small diagrams away from its smaller strata. Finally, the single nef
interval supplied above rules out infinitely many such diagrams.

\smallskip\noindent\emph{Choice of scaling.}
For every possible finite prefix of the detected ordinary truncation scalings from \(X\),
Lemma~\ref{prog:cohomology-transport} makes the transform of
\(H^{1,1}(X,\R)\) onto the working Bott--Chern space
surjective.  There are countably many such sequences: curve
rays belong to countably many integral homology classes, and
a ray determines its contraction and its flip uniquely.
On each working space, two distinct curve rays \(C,C'\) with nonzero
\(L\)-degrees give a proper hyperplane of initial directions defined by
\[
 (L\cdot C)(\omega\cdot C')-(L\cdot C')(\omega\cdot C)=0,
\]
where the degrees use the traces on that working space. Whenever both rays
define finite walls, this equation says that their wall parameters agree.
Avoid all these
hyperplanes, and also avoid \(\omega\cdot C=0\) whenever
\(L\cdot C=0\).  A general Kähler direction, enlarged to
make the initial scaled class Kähler, has these properties.

At a wall the scaled class is pulled back from the contraction
base.  Immediately after a flip, curves in an opposite-side
fiber have positive \(L\)-degree. Another negative ray at
the same wall would satisfy the excluded wall equation with that
opposite-side ray. Immediately after a
divisorial contraction, lift any putative new wall curve
through the projective morphism. Its lifted ray is distinct
from the contracted ray and satisfies the same wall equation; if its
\(L\)-degree were zero, its \(\omega\)-degree would also be
zero.  Both possibilities were excluded.  Hence consecutive
positive thresholds strictly decrease, and each working
interval has nonempty interior.  Its interior nef class is
big: the class on \(X\) is a positive Kähler perturbation of
the pseudo-effective class \(L\), and bigness is preserved
under the birational comparison.  No rational curve can have
zero interior degree.  An \(L\)-negative or \(L\)-positive
such curve would violate nefness at one of the two nearby
parameters; the zero-\(L\) case was excluded.  The gklt
presentation of a positive truncation and
\cite[Theorem~4.3]{HLX} therefore make the interior class
Kähler.  By Lemma~\ref{prog:nef-trivial}, infinitely many
thresholds would tend to zero.

\smallskip\noindent\emph{Reduction to small transformations on a stratum.}
Suppose the sequence is infinite.  The homological divisor
count \cite[Lemma~7.8]{CB} discards all divisorial ambient
steps after a finite prefix.  We next reduce the remaining
flips to small transformations on a stratum.
For a working ambient pair \((X_i,B_i)\), ordinary dlt
adjunction on the normalization \(S\) of an lc center gives
an effective dlt pair \((S,\Delta_S)\) and the
actual rational-line residue identity
\[
 K_S+\Delta_S\sim_{\Q}(K_{X_i}+B_i)|_S.
\]
Its lc centers are the proper nested lc centers.  The
coefficients of these differents form a DCC set depending
only on the original coefficients and the chain length.  At
one restriction they have the form
\[
 1-\frac1r+\frac{\sum_j k_jb_j}{r},
 \qquad 0\leq\sum_j k_jb_j\leq1,
\]
which preserves DCC under iteration.  These assertions,
including the actual residue comparisons, are
\cite[Lemma~7.16]{CB}.  Restricting
Equation~\eqref{prog:comparison} along a strict adjunction chain
gives an effective comparison on strata, strictly positive
for a place whose center maps into the ambient exceptional
locus \cite[Lemma~7.17]{CB}.  Both lemmas are
dimension-free.

Identify a surviving lc center with its strict transform on the common
isomorphism open. There are finitely many lc centers at each stage, and
their surviving list stabilizes. A new zero-discrepancy place was
already a zero-discrepancy place, and strictness keeps the
general point of its center outside the surgery. Starting with the smallest
centers, induct on their dimensions. Fix a surviving center and suppose the
loci already avoid all its proper subcenters. We will prove that the loci
eventually avoid this center as well. The induced diagrams on its normalized
strata are then isomorphisms near their non-klt loci.  On
each stratum, let \(D_S\) be the sum of the restrictions of the ambient
floor components not used in its generic adjunction chain. Strong
\(\Q\)-factoriality of the ambient space makes \(D_S\) rational Cartier,
and the nested-center description gives
\(\Supp D_S=\Supp\lfloor\Delta_S\rfloor\). Slightly lowering this
effective divisor makes the stratum pair klt.

After a further tail the induced transformations on this
stratum are small and their boundary transforms agree.
Indeed, a prime extracted in a restricted transformation
has discrepancy strictly less than \(1\) before extraction:
its new boundary is effective and strict comparison
increases its discrepancy.  Its center is away from the
unchanged non-klt locus.  At the start of this tail there are
only finitely many such places.  On a fixed compact log
resolution, the positive SNC weights \(1-b_j\) have a
positive minimum, and the Jacobian inequality bounds every
place below the strict cutoff \(1\); include also the
finitely many nonexceptional positive-boundary primes.
This is \cite[Lemma~7.9]{CB}, with its lc assertion away
from the non-klt locus.  Discrepancies do not decrease, so
all later extractions belong to this same finite list.
Repeated extraction of one place would give a strictly
decreasing sequence of its boundary coefficients, contrary
to the adjunction DCC property.

Once extractions stop, \cite[Lemma~7.8]{CB} makes
divisorial contractions finite.  Its dimension-free count
is the dimension of the span of prime divisor cycles in
\(H_{2d-2}\) of the compact \(d\)-dimensional stratum.
Restriction to the common isomorphism open uses
Borel--Moore homology: removing codimension at least two
from the target has no effect in this degree, while a lost
prime has nonzero class by its positive Kähler volume.
The count strictly drops.  The coefficients on the
finitely many matched boundary primes then stabilize by
monotonicity and DCC.  Neither morphism of a restricted
diagram can still contract a divisor to the common
intermediate space, since strict comparison would change
the discrepancy of that matched divisor.  The diagrams are
therefore small, with identical boundary transforms.

We justify the passage from an isomorphic restricted diagram to
absence of ambient surgery near the stratum, also in higher
codimension. Its identified boundaries give the same adjunction data,
not only isomorphic underlying spaces. Choose a common projective log
resolution preserving the general points of the surviving lc centers,
and restrict its adjoint comparisons along the strict adjunction
chain described above. The generic point of each surviving center
lies outside the surgery: a zero-discrepancy place remains such a
place, whereas a center contained in the surgery would acquire
strictly larger discrepancy. Thus the strict transform \(S_W\) of
the normalized stratum is not contained in the comparison support.
For a zero-dimensional stratum this already excludes any intersection
with the surgery, so suppose its dimension is positive.
Its restriction is consequently a well-defined effective divisor.
Strict adjunction identifies its class with the difference of the
two pulled-back stratum adjoints. When the restricted diagram is an
isomorphism with the same boundary, that difference is zero.
An effective nonzero divisor on the compact Kähler resolution of
\(S_W\) has positive mass against a Kähler power, so the restricted
divisor is zero.

If an ambient step nevertheless met the stratum, take the first such
step. The full-support assertion following
Equation~\eqref{prog:comparison}, and surjectivity of \(S_W\) onto
the stratum, force its comparison support to meet \(S_W\).
Noncontainment makes this a nonzero effective restriction, and the
cumulative comparison dominates it. This contradicts the vanishing
just proved. Earlier steps are isomorphisms on a neighborhood of the
stratum by the choice of this first step. The ambient diagram is
therefore an isomorphism near the whole stratum. This verifies the
support conclusion used in the induction on proper subcenters;
for a higher-codimension stratum it is noncontainment in the support,
not merely being different from a divisor component, that is needed.

\smallskip\noindent\emph{Excluding the remaining walls.}
It remains to show that only finitely many of these small
diagrams are nontrivial.  Write \(S\) for the first stratum
in this last tail and \(J=K_S+\Delta\) for its adjoint.
Let its restricted original scaling be
\(\{J\}+\lambda H\).  Choose an interior parameter
\(\lambda_0>0\) on this working space, above the remaining
walls, and set
\[
 A=\{J\}+\lambda_0H.
\]
It is Kähler, as the restriction of the interior ambient
Kähler class.  The exact change of parameter is
\begin{equation}
 \{J\}+tA=(1+t)(\{J\}+\lambda H),\qquad
 \lambda=\frac{\lambda_0t}{1+t},\qquad
 t=\frac{\lambda}{\lambda_0-\lambda}.
 \label{prog:reparameterization}
\end{equation}
It applies for \(0\leq\lambda<\lambda_0\).  The new walls
\(t_i\) strictly decrease toward zero, and the traces
\(\{J_i\}+tA_i\) are nef on their successive restricted
working intervals.

On a common resolution of a nontrivial small stratum step, let \(D_i\)
be the real exceptional divisor defined by
\[
 p_i^*\{J_i\}-q_i^*\{J_{i+1}\}=\{D_i\}.
\]
The strict adjoint comparison makes \(D_i\) an effective
nonzero divisor exceptional over both strata.  Equality of
the scaled pullbacks at its wall gives the exact
Bott--Chern equality
\begin{equation}
 p_i^*(\{J_i\}+tA_i)-q_i^*(\{J_{i+1}\}+tA_{i+1})
 =\left(1-\frac{t}{t_i}\right)\{D_i\}.
 \label{prog:restricted-comparison}
\end{equation}
For \(t\) in a later working interval all preceding factors
on the right are nonnegative.  Telescoping on a common
resolution expresses the pullback of \(\{J\}+tA\) as the
pullback of its nef working trace plus an effective divisor
exceptional over that working stratum.
Lemma~\ref{bir:exceptional} identifies the latter divisor
with the negative part.

Push the resulting positive currents to one fixed smooth
resolution of \(S\) and let \(t\downarrow0\).  Closedness of
the pseudo-effective cone proves that \(J\) is
pseudo-effective.  The negative multiplicities converge
valuation by valuation by the limit statement in
Lemma~\ref{prog:negative-along-scaling}. The comparison divisors
have no strict prime of \(S\), since the maps are small.
Their coefficients also vanish at every place whose center
meets the floor: each finite composition is unchanged on a
neighborhood of the floor by the induction on smaller
centers.  Passing to the limit proves
Equation~\eqref{prog:away-floor}.

The stratum has dimension less than \(n\), and the divisor \(D_S\)
constructed above is rational Cartier with support equal to its floor.
Lemmas~\ref{prog:stratum-model} and
\ref{prog:nef-interval} give one small model \(V\) on
which all transformed classes of \(\{J\}+tA\) are nef
for \(0\leq t\leq\delta\).  On a common smooth
resolution \(p:W\to S,\ q:W\to V\), define the real divisor \(F(t)\)
in the common exceptional span by
\begin{equation}
 p^*(\{J\}+tA)-q^*(\{J_V\}+tA_V)=\{F(t)\}.
 \label{prog:affine-comparison}
\end{equation}
Uniqueness of the representing exceptional divisor makes \(F(t)\)
coefficientwise affine in \(t\). It is effective on
\([0,\delta]\): \(-F(t)\) is \(p\)-nef because the
class pulled back from \(V\) is nef, so exceptional
negativity applies.  It is exceptional over \(V\) as
well.  Lemma~\ref{bir:exceptional} gives
\[
 N\bigl(p^*(\{J\}+tA)\bigr)=F(t)
 \qquad(0\leq t\leq\delta).
\]
Every divisorial negative multiplicity is thus affine on
this whole interval.  This remains true on higher
resolutions by Lemma~\ref{bir:pullup}.

Choose a nontrivial stratum wall \(t_i\in(0,\delta)\).
On a common resolution of \(S,V\) and the adjacent
strata, their nef working models compute the same
negative part on their respective open intervals.  Their
two affine expressions agree at \(t_i\).  But
Equation~\eqref{prog:restricted-comparison} says that, at a prime
with coefficient \(d_i>0\) in \(D_i\), their difference
is
\[
 \left(1-\frac{t}{t_i}\right)d_i.
\]
Their slopes differ by \(-d_i/t_i\).  A single affine
function on \([0,\delta]\) cannot agree with both on
two open intervals.  This contradiction excludes every
such wall.  Infinitely many walls would tend to zero, so
the nontrivial diagrams on this stratum are finite.

There are finitely many surviving strata.  Induction on
their dimensions makes all sufficiently late ambient
flipping and flipped loci disjoint from the floor.  This
proves the theorem.
\end{proof}

\begin{corollary}[A signed adjoint supported on the floor]
\label{prog:signed-nef}
Assume \(\Ggood_d\) for every \(d<n\).  In the situation of
Theorem~\ref{prog:special-termination}, suppose moreover
that \(L\) is rationally linearly equivalent, as an actual
rational line, to a signed rational divisor supported on
\(\lfloor B\rfloor\).  The scaling can be chosen to reach
a nef ordinary dlt model after finitely many steps.
\end{corollary}

\begin{proof}
The signed representation persists under codimension-one
pushforward.  A curve disjoint from the transformed floor
has degree zero for each component in this representation,
and hence for the transformed adjoint.  Every negative
operation must therefore meet the floor.  After the finite
prefix supplied by Theorem~\ref{prog:special-termination}
no further negative operation is possible.  The
continuation in Lemma~\ref{prog:nef-trivial} makes the
last adjoint nef.
\end{proof}

\subsection{Polarizations, ordinary replacements, and descent}
\label{prog:foundations}

The program constructions below use the analytic cone
\(\overline{\mathrm{NA}}(T)\), dual to the nef cone in
\(H^{1,1}_{\BC}(T,\R)\) for spaces with rational singularities
\cite[Proposition~2.4]{DHP}. Degrees of global line bundles form a
possibly smaller numerical space. We pass from the analytic cone to
projective geometry by producing an actual relatively ample line.
The ensuing descent statements apply on singular working models as
well as on smooth initial spaces.

\begin{lemma}[Degree-zero Bott--Chern descent]
\label{prog:bc-descent}
Let \(f:T\to Z\) be a proper surjective morphism with connected fibers
between normal compact complex spaces with rational singularities.
Suppose either
\begin{enumerate}
\item \(f\) is bimeromorphic and both spaces belong to Fujiki's class
\(\mathcal C\); or
\item \(f\) is projective and an effective rational boundary
\(\Delta\) makes \((T,\Delta)\) klt with
\(-(K_T+\Delta)\) relatively nef and big.
\end{enumerate}
Pullback is injective on both \(H^2(-,\R)\) and
\(H^{1,1}_{\BC}(-,\R)\). In either group its image consists exactly
of the classes with degree zero on every curve contracted by \(f\).
\end{lemma}

\begin{proof}
These are, respectively, the two cases of
\cite[Lemma~2.6(1),(2)]{DH2024}. The first case is birational descent
between spaces with rational singularities and does not require a
projectivity criterion. The second case uses relative vanishing for
the ordinary rational klt pair and the projective relative curve
space. In particular, it can be applied to a projective log-Fano
contraction before invoking a canonical bundle formula.
\end{proof}

In applications to an ordinary birational negative contraction,
the rationality hypotheses are automatic. The source is locally klt.
For a dlt pair, first decrease its rational Cartier floor slightly;
relative antiampleness persists. Relative vanishing gives
\(R^if_*\OO_T=0\) for \(i>0\)
\cite[Theorem~5.2]{FujinoAnalyticBCHM}. Composing a resolution of \(T\)
with \(f\), the Leray spectral sequence proves that \(Z\) has rational
singularities. The same reasoning works when the klt boundary is
chosen separately on finitely many Stein base neighborhoods. For
fiber-type contractions in case~(ii), rationality of the base follows
from \cite[Lemma~8.8(i)]{DHP}; thus that application also precedes
any canonical bundle formula.

\begin{lemma}[Projectivity from a rational line detecting the ray]
\label{prog:detected-projectivity}
Let $f:T\to Z$ be a proper contraction between normal compact
K\"ahler spaces, with $T$ having rational singularities. Suppose
$\gamma$ is K\"ahler on $Z$ and
\[
 \alpha=f^*\gamma,\qquad
 \overline{\mathrm{NA}}(T)\cap\alpha^\perp=R
\]
is a nonzero ray. If a global rational line bundle $D$ has
$D\cdot R<0$, then $-D$ is relatively ample and $f$ is projective.
\end{lemma}
\begin{proof}
Fix a K\"ahler class $\omega$ on $T$ and normalize $\overline{\mathrm{NA}}(T)$ by
$\omega\cdot z=1$. The resulting slice $S$ is compact. Its unique
point on $R$ has negative $D$-degree, so $d=c_1(D)$ is negative on
a neighborhood of that point in $S$. On the complementary compact
set, $\alpha$ has a strictly positive minimum, while $d$ is bounded.
For all sufficiently small $s>0$, the class $\alpha-sd$ is therefore
strictly positive on $S$ and is K\"ahler by cone duality.

Choose an integer $m>0$ making $mD$ integral. Then
\[
 c_1(-mD)+f^*((m/s)\gamma)=(m/s)(\alpha-sd)
\]
is K\"ahler. On a neighborhood in the base with a potential for
$\gamma$, absorb that potential in a smooth metric on the same global
line bundle $-mD$. Its curvature is positive on the fibres. The analytic
positive-line-bundle criterion and the fibrewise relative-ampleness
criterion make $-mD$ relatively ample. The line bundle was global
throughout, so no gluing of unrelated local polarizations is required.
\end{proof}


An analytic contraction supplied by a supporting class therefore becomes
projective as soon as one rational line has negative degree on its ray.
For a small contraction, that same line will define the flip globally.
The local ordinary adjoints used to prove finite generation need not
glue; their section algebras will be Veronese subalgebras of one algebra.

\begin{lemma}[Divisor representatives over a Stein base]
\label{prog:stein-representatives}
Let $\pi:Y\to U$ be a projective surjective morphism with connected fibres
between normal irreducible complex spaces, where $U$ is Stein, and let
$\mathcal L$ be a holomorphic line bundle on $Y$.
If $\pi$ is bimeromorphic, $\mathcal L$ has a holomorphic section
which is not identically zero, and hence an effective Cartier divisor
representative. More generally, after shrinking $U$ around any specified
point, $\mathcal L$ has a nonzero meromorphic section.
\end{lemma}

\begin{proof}
Suppose first that $\pi$ is bimeromorphic. Grauert's proper direct-image
theorem makes $\mathcal F=\pi_*\mathcal L$ coherent. On the nonempty
open subset $U^\circ$ where $\pi$ is an isomorphism, it is a line bundle.
Choose $u_*\in U^\circ$. Cartan's theorem~A says that global sections
generate $\mathcal F_{u_*}$, so one global section has nonzero image in
$\mathcal F_{u_*}/\mathfrak m_{u_*}\mathcal F_{u_*}$.
Under
\[
 H^0(U,\pi_*\mathcal L)=H^0(Y,\mathcal L)
\]
it gives a section $s$ nonzero at the unique point above $u_*$.
Its zero divisor is effective Cartier and represents $\mathcal L$.
The section may vanish elsewhere; only its meromorphic inverse is
being used to obtain a meromorphic trivialization.

For the general case, fix $u\in U$ and $y\in\pi^{-1}(u)$, and choose
a $\pi$-ample line bundle $\mathcal H$. Relative Serre generation,
after shrinking to a Stein neighbourhood of $u$, gives an integer
$n$ for which both $\mathcal L\otimes\mathcal H^n$ and
$\mathcal H^n$ are relatively generated. Their proper direct images
are coherent. Cartan's theorem~A and the evaluation maps supply
sections $s$ and $t$ of these two bundles which are both nonzero at
$y$. The quotient $s/t$ is the required nonzero meromorphic section
of $\mathcal L$. At no point is $Y$ assumed Stein.
\end{proof}

\begin{proposition}[A global algebra for a detected flip]
\label{prog:detected-flip}
Let $(T,B+\mathbf M)$ be a gklt pair on a normal compact K\"ahler
space which is globally strongly $\Q$-factorial, and let $A\in H^{1,1}_{\BC}(T,\R)$
be its adjoint class. Let $f:T\to Z$ be a projective small contraction
onto a normal compact K\"ahler space. Suppose that the classes of all
$f$-vertical curves lie on one ray $R\subset\overline{\mathrm{NA}}(T)$, that $A\cdot R<0$,
and that a global line bundle $L$ satisfies $L\cdot R<0$.

Then $\bigoplus_{m\geq0}f_*L^m$ is locally finitely generated. The
relative Proj of a sufficiently divisible Veronese is a projective
small contraction $f^+:T^+\to Z$, where $T^+$ is compact K\"ahler and
globally strongly $\Q$-factorial. The reflexive transform $L^+$ is a
rational line bundle, with a positive power equal to the relatively
ample tautological bundle. The transformed boundary and the same
b-nef datum define a gklt pair on $T^+$, with relatively K\"ahler
adjoint $A^+$. This is the generalized flip, with the usual
strict discrepancy increase at centres in either exceptional locus.

Moreover, if $t=(A\cdot C)/(L\cdot C)>0$ for an $f$-vertical curve $C$,
then there is a unique $\eta\in H^{1,1}_{\BC}(Z,\R)$ such that
\[
 A=t\,c_1(L)+f^*\eta,\qquad
 A^+=t\,c_1(L^+)+(f^+)^*\eta.
\]
Every class on $T$ has the corresponding forward transport with an
actual exceptional-divisor correction. No assertion of surjectivity
on Bott--Chern groups is required.
\end{proposition}

\begin{proof}
Since $f$ has a global relatively ample bundle and all its vertical
curves lie on $R$, numerical invariance of ampleness on the projective
fibres shows that $-L$ is $f$-ample.

\smallskip\noindent
\emph{A local rational ordinary adjoint.}
Fix $z\in Z$ and a relatively compact Stein neighbourhood $U$.
Take a projective log resolution $p:V\to T_U$ carrying $\mathbf M$,
and put $\rho=f p$. Write its actual structure boundary as $B_V$, so
\[
 [K_V+B_V]+[\mathbf M_V]=p^*A.
\]
Every coefficient of $B_V$ is less than one. Apply
Lemma~\ref{prog:stein-representatives} to obtain a Cartier
representative $D_L$ of $L$ on $T_U$. The same birational direct-image
argument applies to the coherent rank-one reflexive canonical sheaf:
it supplies a canonical form nonzero at a chosen point of the common
smooth isomorphism locus. Use this meromorphic form on $T_U$ and its
pullback to $V$ to choose compatible canonical representatives, and set
\[
 N=t p^*D_L-K_V-B_V.
\]
For every $\rho$-vertical curve, its degree equals that of $\mathbf M_V$.
Thus $N$ is $\rho$-nef.

The relative bigness of $N$ uses the additional fact that
$\rho=f\circ p$ is \emph{bimeromorphic}: both the small contraction
$f$ and the resolution $p$ are bimeromorphic. It does not follow from
relative nefness alone. In fact every real Cartier divisor is relatively
big for this projective bimeromorphic map over the Stein neighbourhood.
Here is the required ample-plus-effective decomposition explicitly.
Write $N=\sum_j r_jN_j$, where $r_j>0$, $\sum_jr_j=1$, and the $N_j$
are rational Cartier divisors in its finite Cartier span. Choose a
$\rho$-ample line bundle and a Cartier representative
$P$ for it using Lemma~\ref{prog:stein-representatives}.
If $q_jN_j$ is integral, apply the birational case of that lemma to
$\mathcal O_V(q_jN_j-P)$. Its nonzero holomorphic section has an
effective Cartier zero divisor $E_j$, giving the actual relation
$q_jN_j\sim P+E_j$. This use depends on $\rho$ being bimeromorphic;
the section is allowed to vanish along exceptional fibres. Therefore
\[
 N\sim_{\R}
 \underbrace{\left(\sum_j\frac{r_j}{q_j}\right)P}_{H}
 +\underbrace{\sum_j\frac{r_j}{q_j}E_j}_{E}.
\]
Here $H$ is $\rho$-ample and $E\geq0$, proving $\rho$-bigness.
For a map with positive-dimensional general fibres the direct image
used above can be zero; no such bigness claim is made for that setting.
Make the finitely many section choices on a slightly larger Stein
neighbourhood, and now shrink $U$ so its closure is compact in that
neighbourhood. Properness makes the inverse image of this closure
compact. A log resolution of the finitely many resulting divisors
therefore gives a uniform positive bound for the following choice.
Choose $\epsilon>0$ small enough that $(V,B_V+\epsilon E)$ is sub-klt.
The divisor $(1-\epsilon)N+\epsilon H$ is relatively ample. Express it
as a positive combination of rational ample divisors and take general
members of sufficiently high multiples. Relative Bertini
\cite[Theorem~2.20]{DHP}, applied on a log resolution of $B_V+E$,
gives an effective $\Theta\sim_{\R}N$ such that $(V,B_V+\Theta)$ is
sub-klt. The high multiples make the added coefficients arbitrarily
small. All choices are made near the compact fibre; shrink $U$ once
to retain them.

Put $\Delta=p_*(B_V+\Theta)\geq0$. Write the actual linear equivalence
as $\Theta=N+\sum_j a_j\operatorname{div}_V(g_j)$.
The bimeromorphic map $p$ identifies meromorphic function fields, so
$g_j=p^*h_j$. Pushing down gives
$K_{T_U}+\Delta=tD_L+\sum_j a_j\operatorname{div}_{T_U}(h_j)$;
in particular this divisor is real Cartier. Pulling back this same
identity gives
\[
 K_{T_U}+\Delta\sim_{\R}tD_L,
 \qquad p^*(K_{T_U}+\Delta)=K_V+B_V+\Theta.
\]
Thus $(T_U,\Delta)$ is klt. Any finitely many base line bundles in a
relative linear equivalence can first be trivialized by shrinking $U$.
Record the resulting equality as
\begin{equation}\label{prog:detector-rationalization}
 K_{T_U}+\sum_i d_iD_i=tD_L+\sum_j b_j\operatorname{div}(g_j),
 \qquad d_i>0,
\end{equation}
using the positive support of $\Delta$. After a further relatively
compact shrinking, the displayed divisors have finitely many prime
components: their locally finite supports meet a compact inverse image.
Equality of coefficients in \eqref{prog:detector-rationalization} is a
finite rational affine system in $(d_i,t,b_j)$. The klt condition is
open within this system, as seen on one log resolution; the pullbacks
of its adjoints vary linearly by the displayed identity. A nearby
rational solution therefore gives a rational effective klt boundary
$\Delta_q$ and $t_q\in\Q_{>0}$ with
\begin{equation}\label{prog:detector-local-adjoint}
 K_{T_U}+\Delta_q\sim_{\Q}t_qD_L.
\end{equation}
The individual local primes $D_i$ need not be $\Q$-Cartier: the identity
ensures that the total adjoint is $\Q$-Cartier.

\smallskip\noindent
\emph{One global flip algebra.}
The adjoint in \eqref{prog:detector-local-adjoint} is $f_U$-antiample.
Its ordinary flip and local finite generation follow from
\cite[Theorems~1.14 and~1.18]{FujinoAnalyticBCHM}. Clearing the displayed
linear equivalence identifies a Veronese of its canonical algebra
with a Veronese of $\mathcal R:=\bigoplus_{m\geq0}f_*L^m$ over $U$.
Hence $\mathcal R$ is locally finitely generated by
\cite[Lemma~2.26]{FujinoAnalyticBCHM}. The graded pieces are coherent. Over a connected normal base open
the inverse image is irreducible; products of nonzero sections of line
bundles are nonzero there. Thus its graded section algebra is an
integral domain, without choosing any global meromorphic frame of $L$.
Compactness supplies one sufficiently divisible Veronese
$\mathcal R^{(r)}$ generated in degree one. Its relative Proj $T^+$
restricts to the ordinary flip over every such $U$. It is therefore
normal, locally klt, and small over $Z$, and has the global relatively
ample bundle
$\mathcal O_{T^+}(1)$. On the common big open this bundle is $L^r$.
Reflexivity gives
\[
 (L^+)^{[r]}\simeq\mathcal O_{T^+}(1).
\]
This proves that $L^+$ is a rational line bundle before any assertion
of factoriality. Relative positivity over the compact K\"ahler base
makes $T^+$ compact K\"ahler.

For any global rank-one reflexive sheaf $\mathcal F^+$ on $T^+$,
take its coherent reflexive transform $\mathcal F$ on $T$ using a
common resolution. Some $M=\mathcal F^{[m]}$ is a line bundle.
Choose $c\in\Q$ so that $(M+cL)\cdot C=0$, and clear its denominator
to obtain an integral bundle $J=n(M+cL)$ numerically trivial over $f$.
The local rational klt pairs \eqref{prog:detector-local-adjoint} have
antiample adjoint, so Lemma~\ref{prog:integral-descent} yields
one global line bundle $J_Z$ with $J=f^*J_Z$.
On $T^+$, pull back $J_Z$, undo the rational twist $ncL^+$, and compare
on the common big open. After clearing the already established index
of $L^+$, reflexivity gives an invertible positive reflexive power of
$\mathcal F^+$. This proves global strong factoriality.

\smallskip\noindent
\emph{The original generalized pair.}
The local log-Fano pairs above give $R^if_*\mathcal O_T=0$ for $i>0$
by \cite[Theorem~5.2]{FujinoAnalyticBCHM}. They make $T$ locally klt, hence rational.
Leray for a projective resolution then shows that $Z$ is rational.
Lemma~\ref{prog:bc-descent}(i) therefore gives
$f^*H^{1,1}_{\BC}(Z,\R)=R^\perp$, with injective pullback. Thus $A=t c_1(L)+f^*\eta$ for a unique $\eta$.
Define $A^+=t c_1(L^+)+(f^+)^*\eta$; it is relatively K\"ahler.

To construct its actual discrepancy divisor, let $\mathcal I$ be the
ideal defined by
\[
 \operatorname{im}(f^*f_*L^r\longrightarrow L^r)
   =\mathcal I\otimes L^r.
\]
Principalize $\mathcal I$ on a common resolution $p:W\to T$,
$q:W\to T^+$ carrying $\mathbf M$. If
$\mathcal I\mathcal O_W=\mathcal O_W(-F_L)$, the tautological quotient
gives the actual bundle identity
\[
 q^*\mathcal O_{T^+}(1)=p^*L^r\otimes\mathcal O_W(-F_L).
\]
Hence $E_L=F_L/r$ is effective and exceptional over both sides, and
$[E_L]=p^*c_1(L)-q^*c_1(L^+)$. If $B_W$ is the original structure
boundary, put $B_W^+=B_W-tE_L$. Then
\[
 [K_W+B_W^+]+[\mathbf M_W]=q^*A^+.
\]
Its pushforward is $B^+$, and its discrepancies do not decrease.
The local ordinary comparison in \eqref{prog:detector-local-adjoint}
has discrepancy divisor $t_qE_L$, so its strictness proves the stated
strictness for $tE_L$ as well. Thus the original pair remains gklt.

Finally, write any $\theta\in H^{1,1}_{\BC}(T,\R)$ uniquely as
$\theta=f^*\eta_\theta+s c_1(L)$ and set
$\theta^+=(f^+)^*\eta_\theta+s c_1(L^+)$. Its correction is the actual
divisor $sE_L$. For a rational line $M$, the number $s$ is rational. Applying
Lemma~\ref{prog:integral-descent} to a Cartier multiple of $M-sL$
gives $M=sL+f^*M_Z$ as actual rational lines. Reflexive extension
on the common big open gives $M^+=sL^++(f^+)^*M_Z$.
Consequently the cohomological transport agrees with the reflexive
transform on Chern classes.
\end{proof}


\subsubsection{Real boundaries and their linear transport}

The rational ordinary step constructed below also determines the step
for every sufficiently close real boundary. The relevant comparison is
an identity of actual rational lines over the contraction base.

\begin{proposition}[Ordinary real boundaries]
\label{prog:ordinary-replacement}
Let $T$ be normal compact K\"ahler and globally Weil $\Q$-factorial,
with canonical sheaf a rational line bundle. Let $B\geq0$ be a real
boundary with $(T,B)$ klt. Assume the rational ordinary-step result
of Proposition~\ref{prog:ordinary-step} for $T$ and its subsequent models. Then every
$(K_T+B)$-negative extremal analytic ray has an ordinary step with
projective contractions, compact K\"ahler models, and the expected
opposite relative ample signs. Global Weil $\Q$-factoriality is preserved;
global strong $\Q$-factoriality is preserved when imposed initially.
\end{proposition}
\begin{proof}
Put $D=K_T+B$ and fix a $D$-negative extremal ray $R$. On the finite
positive support of $B$, effectiveness, the klt condition and negativity
on $R$ persist under small coefficient changes. Klt openness is checked
on one log resolution of this support. Choose nearby effective rational
klt boundaries $B_0,\ldots,B_m$ and positive real numbers $u_j$ with
\[
 B=\sum_j u_jB_j,\qquad \sum_j u_j=1,
 \qquad D_j\cdot R<0,\quad D_j=K_T+B_j.
\]
For $B=0$ take just $B_0=0$. Every $D_j$ is a global rational line
bundle. The rational-step theorem for $(T,B_0)$ provides a projective
contraction $f:T\to Z$ onto a compact K\"ahler base, contracting
exactly $R$, with $-D_0$ relatively ample. The base is rational by
\cite[Lemma~8.8]{DHP}, so Lemma~\ref{prog:bc-descent} gives
$f^*H^{1,1}_{\BC}(Z,\R)=R^\perp$; use case~(ii) when $f$ is of
fiber type.
Let $C$ be a rational curve spanning $R$, and set
$a_j=(D_j\cdot C)/(D_0\cdot C)\in\Q_{>0}$. On every projective
fibre, $-D_j$ is numerically equivalent to the positive multiple
$-a_jD_0$. Numerical invariance of ampleness on that fibre, followed
by the relative-ampleness criterion, makes $-D_j$ relatively ample.
The positive combination $-D$ is relatively ample as a real line
bundle. This also handles a fibre-type contraction.

Suppose henceforth that $f$ is birational. A Cartier multiple
$M_j=n_j(D_j-a_jD_0)$ is numerically trivial over $Z$ and
$M_j-D_0$ is relatively ample. Lemma~\ref{prog:integral-descent},
applied to the rational klt pair $(T,B_0)$ gives a rational line bundle
$A_j$ on $Z$ with the actual identity
\begin{equation}\label{prog:real-line-descent}
 D_j=a_jD_0+f^*A_j.
\end{equation}
Here and below an identity of rational line bundles means an isomorphism
after a common integral multiple. The descent lemma is used only for
birational $f$.

For a small $f$, let $f^+:T^+\to Z$ be its rational $D_0$-flip.
The rational theorem supplies a global relatively ample adjoint
$D_0^+$ and preserves the stated factoriality and canonical-sheaf
conditions. The other transformed adjoints $D_j^+$ are therefore
rational line bundles. Transform \eqref{prog:real-line-descent} on the common
big open and extend by reflexivity to obtain
\[
 D_j^+=a_jD_0^++(f^+)^*A_j.
\]
All $D_j^+$ are relatively ample, as is
$D^+=\sum_j u_jD_j^+$. Thus this same small map is the required
ordinary real-boundary flip. The real ordinary discrepancy comparison
proves that $(T^+,B^+)$ is klt and gives the strict increases at
exceptional centres.

For a divisorial contraction, write
$D_0=f^*D_0'+e_0E$ with $e_0>0$. Transforming
\eqref{prog:real-line-descent} on the target and summing yields
\[
 D-f^*D'=\left(\sum_j u_ja_j\right)e_0E.
\]
This coefficient is positive. The same discrepancy comparison proves
the required assertions. The ambient category was already preserved
by the rational step in both cases.
\end{proof}


For later use, let \(D=K_T+B\) drive one of the birational steps just
constructed, and write \(f':T'\to Z\) for its positive morphism
(or the identity for a divisorial contraction). Lemma~\ref{prog:bc-descent}
gives the unique decomposition and forward transform
\begin{equation}
 \theta=f^*\eta+s\{D\},\qquad
 \theta'=(f')^*\eta+s\{D'\}.
 \label{prog:ordinary-forward-transport}
\end{equation}
On a common resolution \(p:W\to T\), \(q:W\to T'\), the ordinary
comparison \(p^*D-q^*D'=F\) is an actual effective real divisor,
exceptional over \(T'\) and over both sides for a flip. Hence
\(p^*\theta-q^*\theta'=s\{F\}\). On Chern classes this is the
reflexive transform, by the line identities in the proof above and
Lemma~\ref{prog:integral-descent}. This construction gives forward
transport without assuming surjectivity onto the next Bott--Chern group.

\subsubsection{Actual real lines and relative numerical spaces}

A second polarization argument will be used for maps whose fibers are
Moishezon. It applies to a real combination of global line bundles,
provided that its degrees on the fibers agree with a relatively Kähler
class.

\begin{lemma}[Polarization by an actual real line bundle]
\label{prog:actual-line-polarization}
Let $f:T\to Z$ be a proper Moishezon map, where $T$ is a compact
K\"ahler space. Suppose that $L\in\operatorname{Pic}(T)\otimes\R$
and a relatively K\"ahler class $\omega$ satisfy
\[
 L\cdot C=\omega\cdot C
 \quad\text{for every compact curve $C$ contracted by $f$.}
\]
Then $L$ is relatively ample and $f$ is projective.
\end{lemma}

\begin{proof}
Let $V$ be an irreducible positive-dimensional subspace of a reduced
fibre. It is Moishezon. Normalize $V$ and take a smooth projective
modification $\pi:\widetilde V\to V^{\mathrm n}$, chosen so that
an effective exceptional divisor $E$ has $-E$ relatively ample.
The pullback of the K\"ahler class from $V$ to its normalization is
K\"ahler. Consequently
$\pi^*\omega-\delta[E]$ is K\"ahler for sufficiently small
$\delta>0$. On the projective manifold $\widetilde V$ the actual
real line bundle $\pi^*L-\delta E$ has the same curve degrees as
that K\"ahler class, and is therefore ample. For example, subtract
a sufficiently small positive multiple of a fixed ample class from
the K\"ahler class; the corresponding real line bundle is nef by
the projective numerical criterion. Also $\pi^*L$ is nef, and the
decomposition
\[
 \pi^*L=(\pi^*L-\delta E)+\delta E
\]
shows that it is big. Thus
\[
 L^{\dim V}\cdot V=(\pi^*L)^{\dim V}>0.
\]
The real Nakai--Moishezon criterion for proper algebraic spaces
\cite[Theorem~1.6]{FujinoMiyamoto2020}, applied to the algebraizations
of the Moishezon fibres, proves that $L$ is ample on every reduced
fibre. The criterion is insensitive to nilpotents and reducible
components. In the coefficient space of the finitely many global
line bundles occurring in $L$, choose, for each fibre, a rational
simplex containing $L$ in its interior whose vertices restrict to
ample rational lines on that fibre. Openness of fibrewise ampleness
for each of these finitely many rational lines gives a base
neighbourhood on which the entire simplex is relatively ample.
Choose finitely many such neighborhoods covering the base and
intersect their coefficient neighborhoods of $L$. Every rational
point of the resulting neighborhood is relatively ample globally.
A rational simplex in this intersection also expresses $L$ as a
positive combination of relatively ample rational lines. Clearing
the denominator of one rational point proves projectivity.
\end{proof}

\begin{remark}
The proof uses only positivity of the top self-intersections of
$L$. Equality of curve degrees with $\omega$ does not assert
equality of their higher intersection numbers or of their
Bott--Chern classes.
\end{remark}




The projective morphism needed for the relative program can be chosen
on a smooth model of the rational quotient. Here the Hodge-theoretic
polarization has an elementary construction, independent of any
birational projectivity-descent assertion.

\begin{lemma}[Smooth Hodge-theoretic polarization]
\label{prog:smooth-polarization}
Let \(g:W\to S\) be a surjective morphism with connected fibers
between smooth compact Kähler manifolds. If
\(g^*:H^0(S,\Omega_S^2)\to H^0(W,\Omega_W^2)\) is an isomorphism,
then \(g\) is projective. This applies when its general fiber is
rationally connected.
\end{lemma}

\begin{proof}
We use the smooth argument of \cite[Theorem~3.1, Step~1]{CH}.
Let \(d=\dim W-\dim S\). Choose a rational class \(u\in H^2(W,\Q)\)
close to a Kähler class. Write its Hodge decomposition as
\(u=b+g^*v\), where \(b\) is Kähler and \(v\) has types \((2,0)\)
and \((0,2)\). Hodge type and the projection formula give
\[
 c=g_*(u^d)=g_*(b^d)>0,\qquad
 g_*(u^{d+1})=g_*(b^{d+1})+(d+1)c\,v.
\]
Since pushforward is rational, the class
\[
 \ell=u-g^*\frac{g_*(u^{d+1})}{(d+1)c}
      =b-g^*\frac{g_*(b^{d+1})}{(d+1)c}
\]
is rational and of type \((1,1)\). Lefschetz's theorem makes it the
Chern class of a rational line. The displayed equality supplies a
metric with positive curvature locally over \(S\), by adding a
potential for the base class. Thus this line is relatively ample.
For rationally connected general fibers, holomorphic two-forms are
pulled back from the smooth base, so the hypothesis holds.
\end{proof}

\begin{lemma}[A rational Hodge correction over a smooth base]
\label{prog:relative-hodge}
Let $f:T\to S$ be an already projective surjective morphism of normal
compact K\"ahler spaces. Assume that $S$ is smooth, $T$ has globally
strongly $\Q$-factorial klt singularities, and, for a projective resolution
$r:W\to T$, the morphism $g=f\circ r$ satisfies
\[
 g^*:H^0(S,\Omega_S^2)\xrightarrow{\ \simeq\ }
                         H^0(W,\Omega_W^2).
\]
Then every $\alpha\in H^{1,1}_{\BC}(T,\R)$ can be written
\[
                         \alpha=c_1(L)+f^*\gamma,
\]
where $L\in\operatorname{Pic}(T)\otimes_{\mathbb Z}\mathbb R$ is a
finite real combination of global line bundles and
$\gamma\in H^{1,1}(S,\mathbb R)$.
\end{lemma}

\begin{proof}
The composition $g$ is projective, since all the spaces are compact.
Choose a $g$-ample line bundle on $W$, with integral Chern class $H$.
Set $d=\dim W-\dim S$ and
\[
                            c=g_*(H^d)>0.
\]
Thus $c$ is the positive degree of $H^d$ on a general fibre. The
cohomological pushforward here is the usual pushforward between the
smooth compact manifolds $W$ and $S$, defined over $\mathbb Q$ and of
Hodge bidegree $(-d,-d)$.

Define a rational linear operator on degree-two cohomology by
\[
 \Pi(\eta)=\eta-g^*\!\left(\frac{g_*(\eta\smile H^d)}{c}\right).
\]
Every class of type $(2,0)$ on $W$ is pulled back from $S$, and the
same holds for type $(0,2)$. The projection formula therefore shows
that $\Pi$ kills both types. It preserves type $(1,1)$. Consequently
\[
        \Pi(H^2(W,\mathbb Q))\subset
                         H^2(W,\mathbb Q)\cap H^{1,1}(W).
\]
By the Lefschetz $(1,1)$ theorem the image consists of rational Chern
classes of global line bundles. Applying $\Pi$ to $r^*\alpha$ gives
\[
 r^*\alpha=c_1(L_W)+g^*\gamma,
 \qquad L_W\in\operatorname{Pic}(W)\otimes\mathbb R,
 \qquad
 \gamma=\frac{g_*(r^*\alpha\smile H^d)}{c}\in H^{1,1}(S,\mathbb R).
\]
For example, express $r^*\alpha$ in a rational basis of $H^2(W,\mathbb Q)$
and apply $\Pi$ to each basis vector. This gives the required finite
real combination $L_W$.

Write $L_W=\sum_j a_jL_j$ with actual line bundles $L_j$ on $W$.
For each $j$, let
\[
                   \mathcal F_j=(r_*L_j)^{**}.
\]
It is a rank-one reflexive coherent sheaf. Strong
$\mathbb Q$-factoriality gives an integer $m_j>0$ for which
$M_j=\mathcal F_j^{[m_j]}$ is a line bundle on $T$.
The canonical identification over the isomorphism locus of $r$ gives
\[
             L_j^{\otimes m_j}\simeq r^*M_j\otimes\mathcal O_W(E_j)
\]
for an integral $r$-exceptional divisor $E_j$.
One may obtain this comparison directly from the evaluation map for
$r_*L_j$: the two coherent rank-one sheaves agree away from the
exceptional locus, so their invertible transforms differ by an
exceptional divisor. No global meromorphic frame of $L_j$ or of the
canonical sheaf is required.

Set $L=\sum_j(a_j/m_j)M_j$ and
$E=\sum_j(a_j/m_j)E_j$. The preceding identities yield
\[
                r^*(\alpha-c_1(L)-f^*\gamma)=[E].
\]
The divisor $E$ is $r$-exceptional and is numerically trivial on every
$r$-contracted curve. The exceptional negativity lemma applied in both
signs gives $E=0$. Injectivity of pullback by a resolution, or pushforward
of the equality of currents, now gives
$\alpha=c_1(L)+f^*\gamma$.
\end{proof}


The hypothesis of Lemma~\ref{prog:relative-hodge} persists on each
projective globally strongly \(\Q\)-factorial klt model over this fixed
smooth base: on a common smooth resolution, holomorphic two-forms are
invariant under modifications. Thus the lemma applies separately on
every working model. It supplies real combinations of global lines;
Lemma~\ref{prog:stein-representatives} supplies their Cartier-divisor
representatives on sufficiently small Stein base neighborhoods.

\begin{lemma}[The projective relative cone in Bott--Chern cohomology]
\label{prog:projective-cone}
Let \(f:T\to S\) be projective between normal compact Kähler spaces
with rational singularities, and assume
\[
 H^{1,1}_{\BC}(T,\R)
   =c_1\bigl(\Pic(T)\otimes\R\bigr)
          +f^*H^{1,1}_{\BC}(S,\R).
\]
Let \(N_1^{\mathrm{line}}(T/S)\) denote the space of real combinations
of \(f\)-contracted curves modulo degrees of global line bundles,
and let \(\overline{\mathrm{NE}}^{\mathrm{line}}(T/S)\) be their
closed effective cone. Their analytic classes give an isomorphism
onto the span of the face
\[
 \mathcal F_f=
 \overline{\mathrm{NA}}(T)\cap(f^*\omega_S)^\perp,
\]
where \(\omega_S\) is any Kähler class on \(S\); this isomorphism
identifies \(\overline{\mathrm{NE}}^{\mathrm{line}}(T/S)\) with
\(\mathcal F_f\). In particular, relative extremal rays in the line
space are extremal rays of the full analytic cone.
\end{lemma}

\begin{proof}
The assumed decomposition makes analytic equivalence and line-degree
equivalence identical on combinations of vertical curves. Their map
into the analytic numerical space is therefore well-defined and
injective. We must show that its closed effective cone is the entire
face, including its classes represented by positive currents.

Fix \(z\in\mathcal F_f\). For every \(\gamma\in
H^{1,1}_{\BC}(S,\R)\), both \(a\omega_S+\gamma\) and
\(a\omega_S-\gamma\) are Kähler for large \(a\). Positivity then
shows that \(z\) annihilates \(f^*\gamma\). Fix a global relatively
ample line \(H\). If a real line \(D\) has nonnegative degree on
every vertical curve, the projective relative numerical criterion
makes \(D+\epsilon H\) relatively ample for every \(\epsilon>0\).
Its Chern class becomes Kähler after adding a sufficiently large
multiple of \(f^*\omega_S\). Therefore
\[
                 (c_1(D)+\epsilon c_1(H))\cdot z\geq0,
 \qquad c_1(D)\cdot z\geq0.
\]
Applying this to both signs of a relatively numerically trivial line
shows that pairing with \(z\) factors through the finite-dimensional
space of relative line degrees. It is nonnegative on its nef cone.
By duality for this projective relative space, it is represented by
an element of \(\overline{\mathrm{NE}}^{\mathrm{line}}(T/S)\).
The decomposition in the statement makes its analytic image equal
to \(z\). Conversely every vertical effective curve lies in
\(\mathcal F_f\), and finite-dimensional injectivity preserves
closedness. This proves the cone equality and the extremality claim.
\end{proof}

\subsection{Good models for already projective relative adjoints}

\begin{proposition}[Projective relative completion]
\label{prog:projective-good-model}
Let \(f:X\to S\) be an already projective surjective morphism of
normal compact Kähler spaces, with \(X\) smooth. Let
\((X,B+\mathbf M)\) be an effective generalized klt pair with adjoint
\(A\), carried by a projective resolution \(p:W\to X\) with smooth
\(W\), so that
\[
 p^*A=c_1(K_W)+[B_W]+\mathbf M_W.
\]
Suppose that there are an actual real line bundle
\(J\in\Pic(X)\otimes\R\) and
\(\beta\in H^{1,1}_{\BC}(S)\) such that
\[
 A=c_1(J)+f^*\beta,
 \qquad
 N:=p^*J-K_W-B_W
 \text{ is nef and big over }fp.
\]
Assume that \(J\) is pseudoeffective over \(S\).
Then there is a finite chosen \(A\)-negative projective program over
\(S\), with globally strongly \(\Q\)-factorial compact Kähler
working spaces, and a nonextracting endpoint \(X_m\) such that
\[
 J_m\sim_{\R}g^*H,
 \qquad
 A_m=g^*\bigl(c_1(H)+\rho^*\beta\bigr).
\]
Here \(J_m\) is the actual real-line transform of \(J\),
\(g:X_m\to Z\) is a projective connected-fibre morphism,
\(\rho:Z\to S\) is projective, and \(H\) is an actual real line
bundle ample over \(S\). Thus the class on \(Z\) in this formula is
relatively Kähler over \(S\). The generalized pair on \(X_m\) is
gklt. On a common resolution the comparison with \(A\) is effective
and exceptional over \(X_m\), with strictly positive coefficient
at the strict transform of every prime contracted from \(X\).
\end{proposition}

\begin{proof}
We reduce to fixed ordinary pairs, and construct one global program.
Choose finitely many Stein open sets \(U_\ell\subset S\) and
semianalytic Stein compact subsets \(W_\ell\Subset U_\ell\) whose
interiors cover \(S\). Such compacts are obtained by intersections
with closed polydiscs in local embeddings of \(S\). In particular
they satisfy condition~(P) of \cite{FujinoAnalyticBCHM}: the source
is normal, the base is Stein, the compact is Stein, and its
intersection with any analytic subset defined near it has finitely
many connected components. This remains true on every normal
projective model. Shrinking \(U_\ell\) around \(W_\ell\) preserves
the finite cover.

\smallskip\noindent
\emph{Actual ordinary replacements.}
On each chart represent the finitely many line bundles in the data
by Cartier divisors. For a projective map to a Stein space this can
be done by twisting a line bundle and its prospective meromorphic
frame by sufficiently positive relative lines and using nonzero
sections. The generic-rank-one direct-image argument is an
alternative only when the map is bimeromorphic.
Relative bigness gives
\(N\sim_{\R}P+E\), with \(P\) relatively ample and \(E\geq0\).
For sufficiently small \(\epsilon>0\),
\((W,B_W+\epsilon E)\) is sub-klt near the compact inverse image.
The line \((1-\epsilon)N+\epsilon P\) is relatively ample.
General divided relative sections give an effective representative
\(T\) such that
\[
 \Theta:=\epsilon E+T\sim_{\R}N,
 \qquad (W,B_W+\Theta)\text{ is sub-klt}.
\]
All supports and coefficient margins are fixed near that compact.
With compatible canonical representatives, push the finite
principal-divisor expression for \(\Theta-N\) down and pull it
back again. For \(\Delta_\ell=p_*(B_W+\Theta)\) this gives
\[
 \Delta_\ell\geq0,
 \qquad K_X+\Delta_\ell\sim_{\R}J|_{X_{U_\ell}},
 \qquad
 p^*(K_X+\Delta_\ell)=K_W+B_W+\Theta.
\]
In the last equality the chosen principal correction is included
in the representative of \(K_X+\Delta_\ell\).
Thus \((X_{U_\ell},\Delta_\ell)\) is ordinary klt and its boundary
is relatively big. Lemma~11.15 of \cite{FujinoAnalyticBCHM}, with
its big part equal to \(\Delta_\ell\) and its remaining part zero,
replaces it, in its actual real linear equivalence class, by a klt
boundary with an effective general relatively ample rational
summand. Its condition on non-klt centres is vacuous.
Choose the reserved ample summand to be \(\Q\)-linearly equivalent
to a small rational multiple of a fixed global \(f\)-ample line.
Its later strict transform is consequently \(\Q\)-Cartier by the
transport of that global line, without a local factoriality claim.

We make these choices for a finite rational polytope of actual
global line data. Start with a finite rational span containing
\(J\), and add relatively ample rational lines spanning the global
relative degree space. The above construction is valid in a
neighbourhood of \(J\): retain a small part of the ample summand
and use openness of ampleness and the strict klt inequalities.
It is also valid along a segment from \(J\) to a sufficiently
positive adjoint, by adding divided general ample divisors.
Use finitely many such representatives and one log resolution on
each chart. More explicitly, represent a finite simplex inside the
open ample neighbourhood used above; varying its positive
coefficients represents all nearby global line parameters on the
same finite support. Form the joint finite affine system consisting
of the divisor identities on all charts, with common global line
coordinates and separate boundary and principal-divisor
coordinates. This system is rational. Allow all boundary
coefficients to vary, including those originally contributed by
the real boundary \(B\); fix only the reserved rational ample part
and the zero-support constraints. Its projection to the global
line coordinates contains the just constructed neighbourhood, so
it has a rational affine section. Choose that section sufficiently
close to the constructed real affine lift, by rational approximation
in its affine space of sections, so that the strict klt and
effectiveness margins are retained along the compact segment.
In the joint solution space the
strict klt and effectiveness margins permit a rational polytope
containing the scaling segment. This supplies on every chart
a rational polytope of ordinary klt boundaries with a fixed
effective relatively ample rational summand, representing the
restrictions of the same global line parameters. These boundaries,
including the scaling representatives, are fixed before the
program starts. Their divisors need not agree on different charts.

\smallskip\noindent
\emph{One global scaling construction.}
Apply the global degree-space construction in the proof of
Proposition~\ref{prog:relative-klt} to \(J\) and a general ample
scaling direction in the chosen span. Here are the modifications
that allow real line data and a nonbirational map to \(S\).
The global relatively ample line still gives a compact normalized
slice of the curve cone. The fixed ordinary chart pairs give finite
negative-ray truncations. A general direction avoids equal wall
parameters for independent ray-degree functionals, exactly as in
that proof. A positive wall therefore selects a single ray in the
global degree space. A rational nef support is represented by a
global rational line; a large multiple minus the chartwise ordinary
adjoint is relatively ample. Ordinary relative base point freeness
contracts the ray using the evaluation of this one global line.

If the contraction is small, choose a global rational line having
negative ray degree and construct its flip by its one global
graded algebra. The local ordinary proof of
Proposition~\ref{prog:detected-flip} applies to this algebra:
all curves of the contraction have proportional degrees for global
lines, so the detector and the driving line are positively
proportional over the contraction; ordinary rational replacement
gives local finite generation. This part of that proof needs only
the already projective contraction and these line degrees.
It does not require a prior assertion about the full Bott--Chern
cone. The relative Proj is the next global model. In a divisorial
step the exceptional prime and
Lemma~\ref{prog:integral-descent} give the same continuation as in
Proposition~\ref{prog:relative-klt}. Killing the degree of any global
line by a rational multiple of the detector, then applying that
lemma, preserves global strong \(\Q\)-factoriality in both cases.
The global lines continue to span the relative degree spaces.
The rational support-simplex argument in that proposition makes
the wall a descended relatively ample real line on its contraction
base and gives a nonempty interval of ampleness immediately after
each step. This argument uses relative line data and projectivity,
not birationality of the map to \(S\).

We check explicitly the persistence of the ordinary pairs. Expand
the line parameters in finitely many global rational lines. Their
reflexive transforms are actual rational lines by the preceding
strong property. Pushing forward the fixed principal-divisor
identities gives
\[
 K_{X_i}+\Delta_{\ell,i}\sim_{\R}J_i|_{(X_i)_{U_\ell}}.
\]
Hence these ordinary adjoints are \(\R\)-Cartier; no local
factoriality of individual prime divisors is being assumed.
On a common resolution their comparison minus the appropriate
positive multiple of the detector comparison is exceptional and
relatively numerically trivial. Negativity in both signs makes
this difference zero. Every chosen step is consequently negative
for the fixed ordinary pairs. They remain effective and klt, and
their boundaries remain relatively big. The same equality of
comparison divisors, after adding the fixed base pullback,
preserves the generalized pair and its discrepancy inequalities.

For termination, every already constructed working model is a weak
log canonical model on each chart of a boundary in the fixed
polytope: choose a parameter in its nonempty ample scaling interval.
Indeed, actual descent at a preceding wall \(\lambda_j\) makes
the comparison for \(J+tH\) equal to
\((1-t/\lambda_j)\) times its \(J\)-comparison. For a parameter
in the current interval all preceding \(\lambda_j\) are at least
\(t\), so these comparisons are effective.
Theorem~E of \cite{FujinoAnalyticBCHM} gives finitely many marked
models on each chart. That theorem does not require locally
\(\Q\)-factorial targets. There are therefore finitely many possible
global marked working models. Indeed, isomorphisms between the
restrictions of two existing marked models agree on their common
dense open, and hence on overlaps. A repeated marked working model
would have the same boundary transform and discrepancies, contrary
to strict discrepancy increase at an intervening nontrivial step.
The global program is finite. Relative pseudoeffectivity excludes
a Mori fibre endpoint, so \(J_m\) is nef over \(S\).
This is a finite-cover argument for a single constructed program;
it does not glue independent minimal models.

\smallskip\noindent
\emph{Semiampleness and the global polarization.}
On the endpoint the fixed ordinary pair is klt, its effective
boundary is relatively big, and its adjoint is nef over
\(W_\ell\). Lemma~11.16 of \cite{FujinoAnalyticBCHM} applies to
this already existing weak model. More explicitly, Lemma~11.15
replaces the endpoint boundary by \(A_\ell+B_\ell\), where
\(A_\ell\geq0\) is relatively ample, \(B_\ell\geq0\), and the
pair is klt. Theorem~8.3 then gives semiampleness near \(W_\ell\):
its log canonical, klt-existence, real Cartier, ample-summand and
nef-over-compact hypotheses have all been checked.

The actual equivalences identify these local semiample adjoints
with restrictions of \(J_m\). Their ample models agree on overlaps
by \cite[Lemma~11.3]{FujinoAnalyticBCHM}. To retain a global
polarization, shrink the initial parameter polytope around \(J\)
after the finite program: all its finitely many strict step signs
persist, so the transformed ordinary boundaries remain klt and big.
On each chart the pushforward of its initial common ample rational
summand is effective and relatively big. Apply
\cite[Lemma~11.13]{FujinoAnalyticBCHM} to this common big rational
part and the small endpoint boundary polytope. All pairs are klt,
so the condition on non-klt centres is vacuous. The lemma gives a
common general relatively ample rational summand by one
\(\Q\)-linearly trivial translation. In particular the rational
parametrization is preserved; separate applications of
Lemma~11.15 to individual parameters would not suffice here.
In this finite rational span Theorem~11.17 of \cite{FujinoAnalyticBCHM}
makes the endpoint nef region a rational polytope. Intersect the
finitely many inverse images of these regions and express \(J_m\)
as a positive combination of rational points in its minimal face.
The corresponding global rational lines are semiample on every
chart by the same ordinary base point free argument. After clearing
denominators and taking common sufficiently divisible powers over
the finite cover, each is relatively generated globally: generation on
the charts is generation of its global coherent relative
evaluation map. The product of these finitely many relative
morphisms, followed by Stein factorization, gives a projective
\(g:X_m\to Z\), a projective \(\rho:Z\to S\), and descended
global lines whose positive combination \(H\) is relatively ample.
It realizes the glued ample model and gives
\(J_m\sim_{\R}g^*H\). Adding back \(f_m^*\beta\) proves the
asserted Bott--Chern identity. Composing the effective step
comparisons proves the final discrepancy statement.
\end{proof}

\begin{corollary}[Semiampleness on an existing projective model]
\label{prog:projective-semiample}
Let \(f:T\to S\) be a projective morphism of normal compact Kähler
spaces with \(T\) globally strongly \(\Q\)-factorial, and let
\(J\in\Pic(T)\otimes\R\) be nef over \(S\). Suppose that on a
finite Stein-compact cover satisfying condition~(P), there are
effective ordinary klt boundaries \(\Delta_\ell\), big over the
base, such that
\(K_T+\Delta_\ell\sim_{\R}J|_{T_{U_\ell}}\).
Then \(J\) has the projective relative semiample contraction and
actual descended ample real line of
Proposition~\ref{prog:projective-good-model}.
In particular this applies to any already existing weak model of
the fixed ordinary chart pairs in that proof; one need not run a
new program for the adjoint in question.
\end{corollary}

\begin{proof}
Lemma~11.15 and Theorem~8.3 of \cite{FujinoAnalyticBCHM} give
local semiampleness. For an existing weak model this is also
Lemma~11.16. To globalize the polarization, take a finite rational
span of the global lines defining \(J\). On each chart the ample
part of the normalized boundary allows small perturbations in
this span. The joint rational affine construction in the preceding
proof gives a rational family with a common rational ample part;
all boundary coefficients are allowed to vary. Theorem~11.17 then
gives a rational nef region on each chart. The finite-intersection,
rational-vertex and global-evaluation argument in that proof gives
one projective contraction and its actual descended ample real
line. Only its local ample models are identified by uniqueness.
\end{proof}

\begin{remark}[The two uses of projective completion]
\label{prog:projective-good-model-uses}
For a selective extraction over an already constructed weak model,
the carrier map is projective bimeromorphic. The inherited
exceptional splitting expresses the nef datum, modulo this base,
by actual real lines. Relative nefness follows from the carrier
datum, while relative bigness follows from birationality, using the
ample-plus-effective construction in the ordinary replacement.
The prepared adjoint is a base pullback plus an effective
exceptional divisor supported on exactly the unwanted primes.
Proposition~\ref{prog:projective-good-model} applies. On its endpoint
the remaining error is effective, exceptional over the fixed base,
and relatively nef; negativity makes it zero. Strict discrepancy
improvement prevents contraction of a requested zero-error prime.
The chosen extraction boundary must be effective and klt; in
particular its prescribed primes have the usual admissible
discrepancy range.

For the graph construction, suppose the already projective graph
resolution \(h:V\to Y_m\) has smooth source and
\[
 A_V=h^*A_{Y_m}+[G],\qquad G\geq0.
\]
Then the global relative real line is \(J=\OO_V(G)\), the base
class is \(A_{Y_m}\), and its carrier line is
\(N=G-K_V-B_V\). Its relative degrees equal those of the prepared
nef datum. That datum is nef and big on the prepared carrier, for
example a pullback of its prepared Kähler class; thus \(N\) is nef
and big over \(Y_m\). This bigness is a separate hypothesis check,
since \(h\) can have positive-dimensional general fibres.
Effectivity of \(G\) gives relative pseudoeffectivity. Apply the
proposition over \(Y_m\). The subsequent negativity argument must
still remove \(G_m\); only then is the endpoint adjoint the pullback
of \(A_{Y_m}\) and good over the original base.

For transport of all classes, verify separately that all
Bott--Chern classes on the graph carrier are real-line classes
modulo \(Y_m\), using the smooth rationally connected fibration
and the inherited exceptional splitting on the lower-dimensional
base model. Lemma~\ref{prog:projective-cone} then identifies
the relative degree rays with full Bott--Chern rays, so the
detected-step transport applies. The driving identity involving
\(G\) alone does not give this additional property.
\end{remark}

\begin{corollary}[Ordinary projective Mori output]
\label{prog:projective-mori}
Let \(f:X\to S\) be an already projective surjective morphism of
compact Kähler spaces, with \(X\) smooth. Let \((X,B)\) be an
effective ordinary real klt pair, and let \(D=K_X+B\). If \(D\)
is not pseudoeffective over \(S\), then its projective program over
\(S\), with scaling of a relatively ample actual real line \(N\),
terminates with a Mori fibre contraction. The initial scaling
parameter is chosen so that \(D+TN\) is nef over \(S\).
The working spaces are globally strongly \(\Q\)-factorial and
compact Kähler, and the actual line comparisons and integral
descent of Proposition~\ref{prog:projective-good-model} apply.
\end{corollary}

\begin{proof}
The line \(N\) is fixed; it need not be a general scaling direction.
In particular several extremal rays may have the same wall
parameter. Choose \(\epsilon>0\) such that \(D+2\epsilon N\) is
still not pseudoeffective over \(S\). On the finite Stein-compact
cover choose divided general effective representatives of \(N\)
so that all the resulting ordinary pairs representing \(D+tN\),
\(0\leq t\leq T\), are klt. For
\(\epsilon\leq t\leq T\) reserve a common effective relatively
ample rational summand. The joint rational-family construction in
Proposition~\ref{prog:projective-good-model} puts these fixed pairs
in the polytopes required by Theorem~E of
\cite{FujinoAnalyticBCHM}.

Suppose a finite prefix has reached \(X_i\), and put
\(\mu_i=\lambda_{i-1}\), with \(\mu_0=T\). Its transformed
\(D_i+\mu_iN_i\) is nef. Define
\[
 \lambda_i=\min\{t\in[0,\mu_i]:D_i+tN_i
                         \text{ is nef over }S\}.
\]
The feasible set is a nonempty closed interval. Inductively
\(\lambda_i>2\epsilon\): all previous steps have nonpositive
comparison for \(D+2\epsilon N\), so its transform is still not
pseudoeffective. A first threshold at most \(2\epsilon\) would
make this transform nef by convexity between the current and
preceding thresholds, a contradiction.

There is a \(D_i\)-negative extremal ray \(R_i\) with
\((D_i+\lambda_iN_i)\cdot R_i=0\). To see attainment without a
generality assumption, let \(t<\lambda_i\) tend to \(\lambda_i\).
A ray negative for \(D_i+tN_i\) has positive \(N_i\)-degree,
because \(D_i+\mu_iN_i\) is nef. When \(t>\lambda_i/2\), it is
also negative for \(D_i+(\lambda_i/2)N_i\). The latter has fixed
ordinary klt big-boundary representatives on the charts, since
\(\lambda_i/2>\epsilon\). Normalize each of these boundaries by
\cite[Lemma~11.15]{FujinoAnalyticBCHM} as \(A_\ell+B'_\ell\),
with \(A_\ell\geq0\) relatively ample and the new pair klt.
The \(A_\ell\)-truncation of the cone theorem for
\(K+B'_\ell\) gives finitely many rays negative for the entire
adjoint in question. The finite cover therefore gives finitely
many global negative rays. One of these rays attains the wall,
as required.

Choose one such ray, even if the wall vanishes on other rays.
The rational-support argument for an individual negative ray in
Proposition~\ref{prog:relative-klt} gives a global rational nef
support annihilating exactly \(R_i\), and ordinary relative base
point freeness contracts it. For the real ordinary boundary one
may first choose a nearby rational klt boundary that is still
negative on \(R_i\). The contraction therefore has a rational
ordinary antiample adjoint. Its flip, when small, is constructed
from one global detector algebra as above.

If the contraction is of fibre type, stop. Otherwise the wall
\(w_i=D_i+\lambda_iN_i\) has degree zero on \(R_i\).
In the finite rational span of its actual line factors, this
zero-degree subspace is rational: its defining functional has
rational values on rational lines. Thus \(w_i\) is a real
combination of rational lines of zero ray degree.
Lemma~\ref{prog:integral-descent} descends those lines to the
contraction base. Their combination is nef over \(S\), by lifting
curves through the projective contraction. Pulling it to the
positive side shows that the transformed wall stays nef. The next
threshold therefore satisfies \(\lambda_{i+1}\leq\lambda_i\);
equality is permitted. No interval of ampleness between successive
thresholds is required.

If \(E_i\geq0\) is the ordinary \(D_i\)-comparison on a common
resolution, this actual wall descent gives, for every real \(t\),
\[
 p_i^*(D_i+tN_i)
 \sim_{\R}q_i^*(D_{i+1}+tN_{i+1})
                 +(1-t/\lambda_i)E_i.
\]
For \(t\leq\lambda_i\) the error is effective. This proves all
the nonpositivity and fixed-pair persistence assertions used in
the induction, including the lower bound on thresholds.

At every working model use the \emph{wall parameter} \(t=\lambda_i\),
rather than an interior ample parameter. Since
\(\lambda_i\leq\lambda_j\) for every earlier step, the displayed
comparisons show that \(X_i\) is a weak log canonical model of
the initial fixed chart pairs representing \(D+\lambda_iN\).
Its trace is nef by definition. Theorem~E counts all these marked
weak log canonical models, including the ones at repeated wall
parameters. The finite chart cover gives a finite global marked
list. Repetition would contradict the strict ordinary
\(D\)-discrepancy increase at an intervening step. Thus this
program with the specified \(N\) is finite. Non-pseudoeffectivity
excludes a nef endpoint for \(D\), so the last contraction is a
Mori fibre contraction. The comparison at \(t=1=\lambda_i\) is
crepant: such a step supplies no strict discrepancy improvement
for \(D+N\), while it is still strictly negative for \(D\).
\end{proof}

\subsection{A restricted dimension induction}
\label{prog:dimension-induction}

All generalized nef data in this subsection are globally nef on a fixed
compact K\"ahler carrier. For the degree arguments we use weak NQC:
a positive combination of rational degree-two classes having
nonnegative degree on compact curves. Strong NQC implies this
condition, and the contraction assertion admits this weak form
\cite[Definition~2.43]{HX}. The nef adjoint itself remains nef in
Bott--Chern cohomology throughout. A modified-big boundary means the total trace
$B+\mathbf M_X$ is globally modified big. Relative assertions retain
these absolute hypotheses. A marked model records its bimeromorphic map
from the specified source; two markings agree only under an isomorphism
commuting with those maps. Weak models below may be arbitrary normal compact K\"ahler weak log
canonical models. Their generalized nef traces are understood as
closed currents pushed forward from the fixed nef carrier; only the
whole adjoint is required to be a Bott--Chern class. In particular no
factoriality, and no separate Bott--Chern class for either $K_Y+B_Y$
or the nef trace, is required of an auxiliary weak target. This
category contains both the working models and the adjunction strata
used in special termination. Chosen program outputs remain globally
strongly $\mathbb Q$-factorial.

For dimension at most $d$, write $\mathsf B_d$ for semiampleness of a
nef gklt adjoint with modified-big boundary, including a Moishezon
connected contraction onto a normal compact K\"ahler target. It makes
no assertion that every such contraction is projective.
Write $\mathsf C_d$ for the NQC non-klt contraction assertion of
\cite[Theorem~1.6]{HX}, and $\mathsf C_{d,\mathrm{big}}$ for its
big-adjoint case. The relative assertions are:
\begin{itemize}
\item $\mathsf M_d$: a relatively pseudo-effective gklt adjoint over
any proper map to a normal compact K\"ahler base $S$, with the globally
nef and modified-big data just specified, has a chosen nonextracting
good log terminal model over $S$. Its working space is globally
strongly $\mathbb Q$-factorial and compact K\"ahler; its connected
relative canonical morphism has normal compact K\"ahler target and
is Moishezon. From a smooth common carrier,
all Bott--Chern classes have forward traces, and the chosen model has
the rational degree-two and Bott--Chern exceptional splittings proved
below.
\item $\mathsf F_d$: a compact polytope of these data on one fixed
carrier has a polyhedral relative effective locus, finitely many
relative canonical chambers with chosen good terminal models, and
finitely many marked weak models, accounting for every weak model
in the category specified above.
\end{itemize}
Passing to the proper image and then to its normal Stein factor allows
all base maps to be taken surjective with connected fibres. These
operations preserve the compact K\"ahler base category.

The independent inputs are the analytic cone theorem, the projective
results of Proposition~\ref{prog:projective-good-model}, relative
vanishing, adjunction, relative-canonical positivity, the projective
canonical bundle formula, and the nef-and-big K\"ahler criterion.
The projective constructions use the actual-line lemmas of the preceding
subsection. No assertion of termination of an arbitrary sequence of
flips is used.

\subsubsection{A fixed integral grid and cohomology transport}

\begin{lemma}[Integral degree-two descent for a locally log-Fano map]
\label{prog:integral-h2}
Let $f:X\to Z$ be a projective surjective morphism with connected fibres
between normal complex analytic spaces. Assume that $f_*\mathcal O_X=
\mathcal O_Z$, that $R^if_*\mathcal O_X=0$ for $i>0$, and that every point
of $Z$ has a neighbourhood $U$ on which there is a klt pair
$(X_U,\Delta_U)$ with $-(K_{X_U}+\Delta_U)$ ample over $U$.
Then $f^*:H^2(Z,\mathbb Z)\to H^2(X,\mathbb Z)$ is injective, and its
image consists exactly of the integral classes having degree zero on
every compact curve contracted by $f$.
\end{lemma}

\begin{proof}
Properness, normality and connected fibres identify both
$f_*\mathcal O_X=\mathcal O_Z$ and
$f_*\mathcal O_X^*=\mathcal O_Z^*$. Apply $Rf_*$ to the analytic
exponential sequence. The local surjectivity of
$\mathcal O_Z\to\mathcal O_Z^*$ and the stated coherent vanishing give
\[
 R^1f_*\mathbb Z_X=0,
 \qquad
 R^1f_*\mathcal O_X^*\simeq R^2f_*\mathbb Z_X.
\]
Let $\xi\in H^2(X,\mathbb Z)$ have zero degrees on contracted curves.
A germ of its edge image in $R^2f_*\mathbb Z_X$ is, under this
isomorphism, represented after shrinking $U$ by an actual line bundle
$L$ on $X_U$. Equality of these germs means that, after further
shrinking, $c_1(L)$ and $\xi|_{X_U}$ have the same class. In particular
$L$ has degree zero on every curve in every fibre above this smaller
neighbourhood.

On a sufficiently small Stein base neighbourhood, twisting by a
relatively ample line and taking a quotient of two nonzero relative
sections gives a meromorphic Cartier representative for $L$. Thus
the Cartier-divisor form of base point freeness applies.
The line bundle $L$ is relatively nef, and
$aL-(K_{X_U}+\Delta_U)$ is relatively ample for every real $a$.
The projective analytic base-point-free theorem applies to the Cartier
line $L$. Its conclusion is generation for \emph{every} sufficiently
large integer power, not only for a divisible subsequence
\cite[Theorems~6.2 and~6.5]{FujinoAnalyticBCHM}. Thus, over a further relatively compact
neighbourhood, both $L^m$ and $L^{m+1}$ are relatively generated.
A generated line of degree zero on every curve of a projective connected
fibre defines a constant projective map on that fibre: the restriction
of the generated line is the pullback of $\mathcal O(1)$, and a
positive-dimensional projective image contains a curve of positive
degree. The corresponding relative image is therefore finite and
bimeromorphic over the normal base, and is the base itself. Consequently
\[
 L^m\simeq f^*M_m,\qquad L^{m+1}\simeq f^*M_{m+1}
\]
for line bundles on the neighbourhood in $Z$. Taking their quotient
proves $L\simeq f^*(M_{m+1}\otimes M_m^{-1})$. Its germ in
$R^1f_*\mathcal O_X^*$ is zero, so the edge image of $\xi$ is zero.

The integral Leray spectral sequence, using
$R^1f_*\mathbb Z_X=0$, gives
\[
 0\longrightarrow H^2(Z,\mathbb Z)
 \xrightarrow{f^*}H^2(X,\mathbb Z)
 \longrightarrow H^0(Z,R^2f_*\mathbb Z_X).
\]
Exactness at the middle term proves the assertion. The argument is
integral throughout and does not discard torsion. The converse follows
by restriction of a pulled-back class to a fibre.
\end{proof}

\begin{lemma}[A fixed integral grid along existing crepant steps]
\label{prog:fixed-integral-grid}
Let $X_0$ be a compact K\"ahler space, and suppose that
\[
 \alpha_0=\sum_{j=1}^r a_j\gamma_{0j},\qquad a_j>0,
 \qquad \gamma_{0j}\in H^2(X_0,\mathbb Q),
\]
where each $\gamma_{0j}$ has nonnegative degree on compact curves.
Fix integers $m_j>0$ such that $m_j\gamma_{0j}$ is the image of an
integral cohomology class. Consider an existing sequence of projective
birational steps
\[
 X_i\xrightarrow{f_i}Z_i\xleftarrow{f_i^+}X_{i+1},
\]
where $f_i^+$ is the identity for a divisorial step, $f_i$ satisfies
Lemma~\ref{prog:integral-h2}, and $\alpha_i$ is trivial on its
contracted ray. Assume every $f_i$-contracted curve has degree
proportional to that ray for the classes in the display. Then the
classes $\gamma_{ij}$ descend through $f_i$ and pull back through
$f_i^+$, the same integers $m_j$ remain integral multiples at every
stage, and the transported classes remain nonnegative on curves.
In particular, for every integral compact curve $C\subset X_i$,
\[
 \alpha_i\cdot C>0\quad\Longrightarrow\quad
 \alpha_i\cdot C\geq\delta,
 \qquad \delta:=\min_j\frac{a_j}{m_j}>0.
\]
\end{lemma}

\begin{proof}
Positivity of the coefficients and nonnegativity of the summands imply
that every $\gamma_{ij}$ annihilates the contracted ray. Apply the
integral descent lemma to the chosen integral lift of
$m_j\gamma_{ij}$ and pull its descended class back to $X_{i+1}$.
This keeps $m_j$ fixed.

To verify nonnegativity, take a curve on $X_{i+1}$. If it is contracted
by $f_i^+$, every descended summand has degree zero. Otherwise its
image is a curve $\Gamma\subset Z_i$. The projective map $f_i$ has a
compact curve mapping onto $\Gamma$ with positive degree: take a
component of its inverse image dominating $\Gamma$ and intersect with
sufficiently many relatively ample hyperplanes. Nonnegativity of the
old summand on this curve implies nonnegativity on $\Gamma$, hence on
the curve on $X_{i+1}$. This argument uses a positive-degree lift, not
a degree-one section. Finally $m_j\gamma_{ij}\cdot C$ is a
nonnegative integer, giving the stated lower bound.
\end{proof}

\begin{corollary}[Uniform crepancy parameter; no termination assertion]
\label{prog:uniform-canonical-parameter}
Suppose additionally that the working spaces are $d$-dimensional klt
spaces, that $\alpha_i$ are nef, and that the ordinary canonical
negative-ray length bound $0<-K_{X_i}\cdot C\leq2d$ is available.
For a fixed real number $t>2d/\delta$, every negative extremal ray
chosen in an existing $(K_{X_i}+t\alpha_i)$-program is
$\alpha_i$-trivial. If the steps are projective ordinary canonical
steps, Lemma~\ref{prog:fixed-integral-grid} therefore applies
inductively with this same value of $t$.

If $\alpha_0$ is not big and $\alpha_0-c_1(K_{X_0})$ is big, then
$K_{X_i}+t\alpha_i$ remains non-pseudo-effective along these crepant
steps. Consequently a terminating such program ends with a Mori
fibre contraction. This corollary does not assert that the requisite
contractions exist or that the program terminates.
\end{corollary}

\begin{proof}
A negative ray is $K_{X_i}$-negative because $\alpha_i$ is nef. Choose
its rational curve generator with the ordinary length bound. If its
$\alpha_i$-degree were positive, then
\[
 (K_{X_i}+t\alpha_i)\cdot C\geq-2d+t\delta>0,
\]
a contradiction. The step is thus an ordinary canonical step with
$\alpha_i$ pulled back from its base, and the preceding lemma applies.
On a common resolution the pullback of
$\alpha_{i+1}-c_1(K_{X_{i+1}})$ is the pullback of
$\alpha_i-c_1(K_{X_i})$ plus the effective canonical comparison.
Thus it stays big. Equality of the crepant pullbacks of $\alpha_i$
preserves its nonbigness. The
identity
\[
 (t+1)\alpha_i=(c_1(K_{X_i})+t\alpha_i)
                  +(\alpha_i-c_1(K_{X_i}))
\]
excludes pseudo-effectivity of its first summand. A nef endpoint is
therefore impossible.
\end{proof}


\begin{lemma}[Cohomology splitting along detected steps]
\label{prog:detected-cohomology}
Let $X_0$ be a smooth compact K\"ahler manifold. Consider a finite
sequence of divisorial contractions and small flips
\[
                         X_0\dashrightarrow X_1\dashrightarrow\cdots
                         \dashrightarrow X_r
\]
between normal globally strongly $\mathbb Q$-factorial compact
K\"ahler spaces with rational singularities.
Assume that every negative contraction $f_i:X_i\to Z_i$ is a
projective contraction of one ray of the analytic cone, that $Z_i$
has rational singularities, and that each small flip is the detected
flip of a global rational line $L_i$ with $L_i\cdot R_i<0$.
For a divisorial contraction assume the usual single exceptional
prime. Then, for every projective smooth resolution $p:U\to X_i$,
\begin{align*}
 H^2(U,\mathbb R)
   &=p^*H^2(X_i,\mathbb R)
       \oplus\bigoplus_{E\text{ exceptional for }p}\mathbb R[E],\\
 H^{1,1}(U,\mathbb R)
   &=p^*H^{1,1}_{\rm BC}(X_i,\mathbb R)
       \oplus\bigoplus_{E\text{ exceptional for }p}\mathbb R[E].
\end{align*}
The first equality also holds over $\mathbb Q$. There are compatible
forward linear transforms on $H^2(-,\mathbb R)$ and on Bott--Chern
cohomology. They are surjective, are isomorphisms for a small flip,
and preserve rational degree-two classes. On a common resolution the
pullback difference is an actual exceptional divisor class; for a
flip it is exceptional over both spaces.

The exceptional quotient also takes the Chern class of a holomorphic
line bundle on $U$ to a rational holomorphic line class on $X_i$.
\end{lemma}

\begin{proof}
For a modification of a smooth compact K\"ahler manifold the two
splittings are the degree-two modification formula. Suppose that they
hold for $X=X_i$. We use only the following birational descent fact:
for a proper bimeromorphic map between normal compact spaces in
Fujiki's class with rational singularities, pullback is injective on
$H^2(-,\mathbb R)$ and on Bott--Chern cohomology, and its image in
either group consists of the classes having degree zero on every
contracted curve. This is the bimeromorphic case of
\cite[Lemma~2.6(1)]{DH2024}.

First note that a degree-two class $\xi\in H^2(X,\mathbb R)$ has
the same degrees on curves as a Bott--Chern class. Indeed, pull it to
a smooth projective resolution $a:W\to X$ and take the $(1,1)$-part
of its Hodge decomposition. The induction hypothesis writes that
part as $a^*\beta+[F]$, with $F$ exceptional over $X$. On every
$a$-contracted curve, $a^*\xi$ and its $(2,0)$ and $(0,2)$ parts
have degree zero. Hence $F$ has degree zero on every such curve.
Exceptional negativity applied to both signs gives $F=0$.
Projective multisections lifting a curve of $X$ then show that
$\xi$ and $\beta$ have the same degrees on every curve of $X$.

Suppose first that $X\dashrightarrow X^+$ is a small detected flip.
Write $f:X\to Z$ and $f^+:X^+\to Z$, and fix an integral curve
$C$ on its negative ray. For $\xi\in H^2(X,\mathbb R)$ set
\[
                    s=\frac{\xi\cdot C}{c_1(L)\cdot C}.
\]
By the preceding paragraph, $\xi-s c_1(L)$ vanishes on every
$f$-contracted curve. Thus it is $f^*\eta$ for a unique
$\eta\in H^2(Z,\mathbb R)$. Define
\[
                    T\xi=(f^+)^*\eta+s c_1(L^+).
\]
For a Bott--Chern class, the same descent takes place in
Bott--Chern cohomology, so the two transforms agree. The detected
flip construction gives an effective rational divisor $E_L$ on a
common projective resolution $a:W\to X$, $b:W\to X^+$ such that
\[
             a^*c_1(L)-b^*c_1(L^+)=[E_L].
\]
Consequently
\[
                         a^*\xi=b^*T\xi+s[E_L].
\]
If $\xi$ is rational, then $s$ is rational. Since the image of
$f^*$ on real degree-two cohomology is the realification of a
rational subspace, its unique preimage $\eta$ is rational as well.
This proves rationality of $T\xi$.

The sets of prime divisors exceptional for $a$ and $b$ are identical:
the two working spaces are small bimeromorphic. The old splitting
and the displayed pullback identity therefore span
\[
              H^2(W,\mathbb R)
                =b^*H^2(X^+,\mathbb R)
                   +\operatorname{span}_{\mathbb R}\{[E]:E\text{ is }b
                   \text{-exceptional}\},
\]
and similarly in bidegree $(1,1)$. The sum is direct. If a pullback
from $X^+$ equals an exceptional divisor class, that divisor has
degree zero on all $b$-contracted curves, so both signs of negativity
make it zero; pullback injectivity then makes the original class
zero. In particular the dimensions on the two sides agree, $T$ is
injective by the old direct sum, and $T$ is surjective. The rational
version follows either from the same argument or from realification.

For a divisorial contraction $f:X\to X^+$, let $D$ be its
exceptional prime. Its degree on the contracted ray is negative:
otherwise $D$ would be $f$-nef, contrary to exceptional negativity.
For any $\xi$, subtract
\[
                    \frac{\xi\cdot C}{D\cdot C}[D]
\]
and descend the result by the same birational cohomology lemma.
This defines $T\xi$. On a common resolution its pullback difference
is a multiple of the total transform of $D$. The target's exceptional
prime set is the old one together with the strict transform of $D$.
The same spanning and directness argument gives the target splitting
and surjectivity; its kernel is $\mathbb R[D]$. Rationality is again
immediate from the quotient of rational intersection numbers.

These arguments initially prove the splitting on a chosen common
resolution. To obtain it on any projective smooth resolution
$U\to X^+$, dominate that resolution and the chosen one by a smooth
common projective resolution. Apply the smooth degree-two
modification formula upstairs and push down to $U$. An exceptional
divisor pushes either to zero or to a divisor exceptional over $X^+$.
This proves spanning on $U$; negativity gives directness there.

Finally let $\mathcal H$ be a holomorphic line on $U$ and put
$\mathcal Q=(p_*\mathcal H)^{**}$. Global strong
$\mathbb Q$-factoriality gives an $m>0$ for which
$\mathcal M=\mathcal Q^{[m]}$ is invertible. The coherent evaluation
comparison gives
\[
                  \mathcal H^{\otimes m}
                  \simeq p^*\mathcal M\otimes\mathcal O_U(F)
\]
for an integral $p$-exceptional divisor $F$. This comparison is
obtained from the coherent direct image and evaluation map; it does
not assume that $\mathcal H$ has a global meromorphic section.
Dividing first Chern classes by $m$ proves the line assertion.
\end{proof}


\begin{corollary}[Stability under projective extraction]
\label{prog:extraction-cohomology}
Let $r:Y\to X$ be a projective bimeromorphic morphism of normal compact
K\"ahler spaces with rational singularities. Suppose that $Y$ is globally
strongly $\mathbb Q$-factorial and that $X$ has the splitting property in
the preceding lemma. Then $Y$ has that property as well. More precisely,
if $E_1,\ldots,E_s$ are the prime divisors exceptional for $r$, then
\[
 H^2(Y,\mathbb R)
   =r^*H^2(X,\mathbb R)\oplus
      \bigoplus_{j=1}^{s}\mathbb R[E_j],
\]
and the analogous statements hold over $\mathbb Q$ and in
Bott--Chern cohomology.
\end{corollary}

\begin{proof}
Choose any projective smooth resolution $q:W\to Y$. The composite
$p=rq$ is a projective resolution of $X$. Its exceptional primes are
exactly the $q$-exceptional primes together with the strict transforms
$\widetilde E_j$ of the $E_j$. Since every $E_j$ is $\mathbb Q$-Cartier,
\[
 [\widetilde E_j]=q^*[E_j]+[F_j]
\]
for a rational $q$-exceptional divisor $F_j$. Substitute these identities
in the splitting for $p$. The result spans $H^2(W,\mathbb R)$ by
$q^*H^2(Y,\mathbb R)$ and the $q$-exceptional divisor classes. Their sum
is direct: an exceptional divisor whose class is a pullback has degree
zero on every $q$-contracted curve, and applying exceptional negativity
to both signs makes the divisor zero. Pullback injectivity then applies.
The same argument in bidegree $(1,1)$ and over $\mathbb Q$ proves the
splitting property for $Y$.

Now expand $q^*\xi$, for $\xi\in H^2(Y,\mathbb R)$, using the original
$p$-splitting and replace the $[\widetilde E_j]$ as above. The difference
between $q^*\xi$ and a class in
$q^*(r^*H^2(X,\mathbb R)+\sum_j\mathbb R[E_j])$ is $q$-exceptional,
and is therefore zero by the direct sum just proved. Thus the displayed
formula for $H^2(Y,\mathbb R)$ spans. Its directness follows from
exceptional negativity for $r$ and pullback injectivity. The rational
and Bott--Chern statements follow by the identical argument. The line
bundle assertion of the preceding lemma uses only global strong
$\mathbb Q$-factoriality and hence applies to $Y$ as well.
\end{proof}



\paragraph{Support in the special-termination comparison.}
\label{prog:special-support}
The following support comparison applies to the detected steps in
the dimension induction. Let $\psi:X\dashrightarrow X^+$ be a detected
small negative step with maps $f:X\to Z$ and $f^+:X^+\to Z$,
and let $L$ be its negative detector. For a sufficiently
divisible $r>0$, put
\[
 \operatorname{im}(f^*f_*L^r\longrightarrow L^r)
      =\mathcal I\otimes L^r.
\]
The zero set of $\mathcal I$ is exactly the exceptional locus of $f$.
Outside that locus $f$ is an isomorphism, so the evaluation map is
surjective. On a positive-dimensional projective fibre, every section
of $L^r$ vanishes: its restriction to each integral curve has negative
degree, and those curves cover every positive-dimensional fibre
component. Connectedness supplies the whole nontrivial fibre.

On a resolution principalizing $\mathcal I$, the detector comparison
is a positive multiple of the effective divisor defined by
$\mathcal I\mathcal O_W$. Consequently its support is the full
inverse image of the exceptional locus. The driving adjoint differs
from a positive real multiple of $c_1(L)$ by a class pulled back from
the contraction base. Its comparison divisor is therefore the same
positive multiple of the detector comparison: their class difference
is exceptional and zero, so both-sign negativity makes it zero as a
divisor. The comparison for a longer
negative program is at least this first-step divisor, by discrepancy
monotonicity on a common resolution.

Suppose the beginning and end restrictions to a surviving normalized
lc stratum $S$ are isomorphic and their strict adjunction data agree.
The general point of this stratum avoids every intervening surgery:
a zero-discrepancy place remains such a place, whereas strict
comparison would increase its discrepancy if its center were contained
in the exceptional locus. Thus its strict transform $S_W$ on a common
resolution is not contained in the comparison support. Iterated strict
adjunction identifies the restriction of the effective comparison with
the difference of the identified adjoints, so its class is zero.
A nonzero effective divisor on a positive-dimensional compact K\"ahler
stratum has positive K\"ahler mass. The restricted divisor is therefore
zero. If the first step meeting the stratum were nontrivial there,
the full inverse-image support and surjectivity $S_W\to S$ would make
this restriction nonzero, a contradiction. For a zero-dimensional
stratum, generic-point avoidance already excludes an intersection.
The whole program is consequently an isomorphism near the stratum.
This is the same strict-adjunction support argument used in the proof
of Theorem~\ref{prog:special-termination}; it does not use its
lower-dimensional abundance assumption.

\begin{proposition}[A uniform relative NQC bound]
\label{prog:relative-nqc-grid}
Let $f_0:U\to Z$ be an already projective surjective morphism with
connected fibres between smooth compact
K\"ahler manifolds, and suppose that
\[
 f_0^*:H^0(Z,\Omega_Z^2)\longrightarrow H^0(U,\Omega_U^2)
\]
is an isomorphism. Let $\alpha_0$ be a nef class with an expression
\[
       \alpha_0=\sum_{j=1}^k r_j\eta_{j,0},\qquad
       r_j>0,\quad \eta_{j,0}\in H^2(U,\mathbb Q),
\]
where $\eta_{j,0}$ has nonnegative degree on every curve vertical
over $Z$. Let $D_0=K_U+\Delta_U+\mathbf M_U$ be a generalized klt
adjoint. Consider its projective relative program for
$D_0+a\alpha_0$, under the projective local ordinary reduction described in the
proof of Lemma~\ref{prog:base-point-free-induction}.

If $\alpha_0=0$, triviality is immediate. Otherwise there are
positive integers $m_j$, chosen once on $U$, and a number
\[
                         \delta=\min_j\frac{r_j}{m_j}>0
\]
such that, for $a\delta>2\dim U$, the entire relative program is
$\alpha_0$-trivial. The same number $a$ works at every step.
\end{proposition}

\begin{proof}
\emph{Rational classes modulo the fixed base.}
Choose an $f_0$-ample line bundle with integral Chern class $H$, and
put $e=\dim U-\dim Z$ and $c=f_{0*}(H^e)>0$. The rational operator
\[
 \Pi(\eta)=\eta-
 f_0^*\!\left(\frac{f_{0*}(\eta\smile H^e)}{c}\right)
\]
preserves type $(1,1)$ and kills types $(2,0)$ and $(0,2)$.
Indeed these latter classes are pulled back from $Z$, and the
projection formula applies. Thus
$\Pi(H^2(U,\mathbb Q))\subset H^{1,1}(U)\cap H^2(U,\mathbb Q)$.
By the Lefschetz $(1,1)$ theorem, choose an actual line bundle
$L_{j,0}$ and an integer $m_j>0$ with
\[
 \eta_{j,0}=\frac1{m_j}c_1(L_{j,0})+f_0^*\gamma_j,
 \qquad \gamma_j\in H^2(Z,\mathbb Q).
\]
The original summands $\eta_{j,0}$ need not have type $(1,1)$.
Their base components account for this. Their weighted sum gives
\begin{equation}\label{prog:eq:relative-nqc}
 \alpha_0=\sum_j\frac{r_j}{m_j}c_1(L_{j,0})+f_0^*\Gamma,
 \qquad \Gamma=\sum_jr_j\gamma_j\in H^{1,1}(Z,\mathbb R).
\end{equation}
The final assertion follows because $f_0^*$ is an injective Hodge
map and the other terms have type $(1,1)$. The identity is therefore
also one of Bott--Chern classes.
Each $L_{j,0}$ has nonnegative integral degree on every
$Z$-vertical curve.

\smallskip\noindent
\emph{Induction across an actual contraction.}
Suppose at a finite stage $f_i:U_i\to Z$ we have actual line bundles
$L_{j,i}$, nef over $Z$, and
\[
 \alpha_i=\sum_j\frac{r_j}{m_j}c_1(L_{j,i})+f_i^*\Gamma.
\]
Assume also that $\alpha_i$ is nef and that the transformed
$D_i$ is generalized klt. These assertions hold initially.
Let $R$ be a $(D_i+a\alpha_i)$-negative extremal ray over $Z$.
It is $D_i$-negative because $\alpha_i$ is nef. Lemma~\ref{prog:relative-hodge} and Lemma~\ref{prog:projective-cone}
identify the relative curve cone with the corresponding face of the
full analytic cone, at the initial model and after every step. The
analytic cone theorem therefore supplies a rational generator $C$ with
\[
                     0<-D_i\cdot C\leq2\dim U.
\]
All $L_{j,i}\cdot C$ are nonnegative integers. If
$\alpha_i\cdot C>0$, one of them is at least one, and hence
$\alpha_i\cdot C\geq\delta$. Therefore
\[
 (D_i+a\alpha_i)\cdot C
             \geq-2\dim U+a\delta>0,
\]
a contradiction. We have $\alpha_i\cdot C=0$ and
$L_{j,i}\cdot C=0$ for every $j$. Every curve contracted by the
projective extremal contraction $h_i:U_i\to Y_i$ has class on
$R$ when tested by global line bundles, so the same vanishing
holds for all of them.

On each fixed Stein chart of $Z$, choose the effective real
ordinary representative of the initial relatively ample nef
datum once, before running the program. Its actual real linear
equivalence to the fixed real line-bundle representative persists
under pushforward. The Bott--Chern comparison with the generalized
adjoint modulo the fixed base class persists separately by
exceptional negativity, as detailed below.
Thus every $h_i$ is, locally on $Y_i$, a contraction with an
ordinary real klt log-Fano boundary. On a relatively compact base
chart, rational approximation in the finite affine system of Cartier
adjoint identities gives a rational effective klt boundary with the
same antiample sign, as in Proposition~\ref{prog:detected-flip}.
Lemma~\ref{prog:integral-descent} therefore applies.
Consequently
\[
 M_{j,i}:=(h_i)_*L_{j,i}\ \text{is a line bundle},\qquad
 h_i^*M_{j,i}=L_{j,i}.
\]
For a divisorial contraction set $L_{j,i+1}=M_{j,i}$.
For a flip $h_i^+:U_{i+1}\to Y_i$, set
$L_{j,i+1}=(h_i^+)^*M_{j,i}$.
These are actual line bundles with the \emph{same} integers $m_j$.
The descended Bott--Chern class
\[
 \alpha_{Y_i}
     =\sum_j\frac{r_j}{m_j}c_1(M_{j,i})+f_{Y_i}^*\Gamma
\]
satisfies $h_i^*\alpha_{Y_i}=\alpha_i$.
It is nef by descent of nefness through the projective surjection
$h_i$; its pullback to the next model is $\alpha_{i+1}$.
This proves that the step is crepant for $\alpha_i$. Hence it is
also a $D_i$-negative step, and the original generalized pair
remains gklt.

Finally each $M_{j,i}$ is nef over $Z$. For a $Z$-vertical curve
$B\subset Y_i$, projectivity supplies a curve $B'\subset U_i$
mapping onto it with some positive degree $d_B$. Then
\[
 d_B(M_{j,i}\cdot B)=L_{j,i}\cdot B'\geq0.
\]
No assertion $d_B=1$ is required. Pullback preserves this relative
nefness, so the induction applies on $U_{i+1}$. Its curve degrees
are again integers because the transported objects are line
bundles, not because any curves were lifted with degree one.
This proves the fixed bound through the whole program.
\end{proof}

\begin{remark}[Scope of the degree grid]
The transported rational classes
\[
 \eta_{j,i}=\frac1{m_j}c_1(L_{j,i})+f_i^*\gamma_j
\]
have degrees in $m_j^{-1}\mathbb Z_{\geq0}$ on curves vertical over
the fixed smooth base $Z$. This relative assertion is exactly what
the initial projective program requires. It does not assert
nonnegativity of these individual classes on all curves of every
birational model, and uses neither rational connectedness of
resolution fibres nor Graber--Harris--Starr.
\end{remark}

\subsubsection{The negative-part comparison for a nef adjoint}
We write $N_\sigma(\xi)$ for the divisorial negative part of a
pseudo-effective class. We use its homogeneity, subadditivity,
continuity after adding a vanishing K\"ahler perturbation, and the
identities
\[
 N_\sigma(r^*\xi+[E])=N_\sigma(r^*\xi)+E,
 \qquad r_*N_\sigma(r^*\xi)=N_\sigma(\xi)
\]
for an effective $r$-exceptional divisor $E$. A nef class has zero
negative part. These are the divisorial Zariski-decomposition
properties of \cite[Section~3]{Boucksom2004} and
\cite[Appendix~A]{DHY}. No identity
$N_\sigma(r^*\xi)=r^*N_\sigma(\xi)$ for arbitrary $\xi$ is assumed.

\begin{lemma}[The absolute negative part after a relative program]
\label{prog:negative-part-program}
Let $\mu:U\to X$ be a projective resolution, let $\alpha$ be nef on $X$,
let $c>0$, and let $F_U\geq0$ be $\mu$-exceptional. Put
\[
 D_U=c\mu^*\alpha+[F_U].
\]
Suppose that $\phi:U\dasharrow V$ is a finite $D_U$-negative program,
possibly relative to another base, and denote its transformed class and
divisor by $D_V$ and $F_V=\phi_*F_U$. Then
\[
                         N_\sigma(D_V)=F_V.
\]
\end{lemma}

\begin{proof}
Take a common projective resolution $p:W\to U$, $q:W\to V$.
The negativity comparison for the finite program is
\[
                p^*D_U=q^*D_V+[E],
 \qquad E\geq0,\quad E\text{ is }q\text{-exceptional}.
\]
This comparison is valid for a relative negative program: its proof
uses the negativity of the contracted rays and the positive adjoint on
each flipped side, not absolute nefness of the final adjoint.
Since $p^*F_U$ is effective and exceptional over $X$, the exceptional
identity and nefness of $(\mu p)^*\alpha$ give
\[
 N_\sigma(p^*D_U)=p^*F_U
                 =N_\sigma(q^*D_V)+E.
\]
Pushing forward by $q$ proves the assertion. In particular, the
negative part in this statement is absolute, even when the program
that produced $V$ was relative.
\end{proof}

\subsubsection{Descent of the effective vertical divisor}

\begin{lemma}[Vertical numerical triviality]
\label{prog:vertical-divisor}
Let $g:V\to S$ be a projective surjective morphism with connected fibres.
Assume that $S$ is globally Weil $\mathbb{Q}$-factorial. Let $F\geq0$ be
an effective real Cartier divisor on $V$, vertical over $S$, such that
\[
                    F\cdot C=0
       \quad\text{for every curve }C\text{ contracted by }g.
\]
Then there is an effective real Cartier divisor $F_S$ on $S$ such that,
as actual divisors,
\[
                              F=g^*F_S.
\]
\end{lemma}

\begin{proof}
For a prime divisor $P\subset S$, write $m_Q$ for the multiplicity of
$Q$ in $g^*P$, where $Q$ ranges over the prime divisors dominating $P$.
These multiplicities may be computed over the smooth generic locus of
$P$, where $P$ is Cartier. Put
\[
 b_P=\min_{g(Q)=P}\frac{\operatorname{coeff}_Q F}{m_Q},
 \qquad F_S=\sum_P b_PP.
\]
Only finitely many $b_P$ are nonzero, since any such $P$ is the image of
a component of $F$. Thus $F_S$ is a genuine effective Weil real divisor.
Global Weil $\mathbb{Q}$-factoriality makes this finite divisor real
Cartier.

We first check equality over codimension one in $S$. Work near a
general point of $P$ and restrict to a transverse disk. Resolving the
source, and cutting by general relative ample hypersurfaces if the
relative dimension exceeds one, reduces to a projective surface over
that disk with connected fibres. These operations can be performed
over a relatively compact Stein neighbourhood; they do not require
global sections on $S$. The intersection matrix of the components of
the special fibre is negative semidefinite, with kernel generated by
the full fibre with its multiplicities. The restriction of $F$ has
zero intersection with every fibre component. Its coefficients are
therefore proportional to those multiplicities. Equivalently, all the
ratios in the definition of $b_P$ are equal. When $g$ is birational this
assertion over the generic point of $P$ is immediate.

Consequently
\[
                              G=F-g^*F_S
\]
is a signed real Cartier divisor supported over a subset of codimension
at least two in $S$. Moreover, $G$ is numerically trivial over $S$.
We check that such a signed exceptional divisor is zero. The assertion
is local on $S$. Over a relatively compact Stein neighbourhood, take
$d=\dim V-\dim S$ general relative ample hypersurfaces and resolve their
intersection. They give a projective generically finite morphism
$H\to S$. The cuts may be chosen to meet any prescribed component of
$G$ in a nonzero divisorial trace. After normalization and Stein
factorization, $H\to S$ factors as a projective birational morphism
$H\to S'$ followed by a finite morphism $S'\to S$. The restricted
divisor $G|_H$ is exceptional for $H\to S'$ and numerically trivial
over $S'$. Applying the ordinary exceptional negativity lemma to both
$G|_H$ and $-G|_H$ gives $G|_H=0$. The choice of the cuts therefore
excludes every nonzero component of $G$. Hence $G=0$.

The surface argument above is also the usual proof of the
degenerate-divisor negativity statement: after subtracting the minimum
multiple of the full fibre, an effective residual fibre divisor omits
a component and cannot be numerically trivial. Notice that projectivity
of $g$ was a hypothesis; it was not deduced from factoriality.
\end{proof}

\subsubsection{The negative part on the lower-dimensional base}

\begin{lemma}[Detecting the negative part by pullback currents]
\label{prog:negative-part-base}
Let $g:V\to S$ be a surjective morphism. Suppose that $\alpha_S$ is nef,
$F_S\geq0$ is real Cartier, $c>0$, and
\[
 D_S=c\alpha_S+[F_S],\qquad D_V=g^*D_S,\qquad
 N_\sigma(D_V)=g^*F_S.
\]
Then $N_\sigma(D_S)=F_S$.
\end{lemma}

\begin{proof}
Nefness gives $N_\sigma(D_S)\leq F_S$. Fix a prime divisor $P\subset S$
and a prime divisor $Q\subset V$ dominating it. Write
$m=\operatorname{mult}_Q(g^*P)>0$ and $b=\operatorname{coeff}_P F_S$.
At their generic smooth points, pullback of positive currents satisfies
\[
                         \nu(g^*T,Q)=m\nu(T,P).
\]
Indeed, the divisorial term $\nu(T,P)[P]$ pulls back with multiplicity
$m$, and the residual current has zero generic Lelong number there.

Choose K\"ahler forms $\omega_S,\omega_V$ and a constant $C>0$ for which
$C\omega_V-g^*\omega_S$ is K\"ahler. For $\varepsilon>0$, let $T_\varepsilon$
be any positive current in the big class $D_S+\varepsilon[\omega_S]$.
Minimality of the multiplicity and monotonicity under adding a nef
class give
\begin{align*}
 m\nu(T_\varepsilon,P)
 &=\nu(g^*T_\varepsilon,Q)\\
 &\geq\nu(D_V+\varepsilon g^*[\omega_S],Q)\\
 &\geq\nu(D_V+C\varepsilon[\omega_V],Q).
\end{align*}
Taking the infimum over $T_\varepsilon$, and then letting
$\varepsilon\downarrow0$, gives
\[
 m\nu(D_S,P)\geq\nu(D_V,Q)
                 =\operatorname{coeff}_Q(g^*F_S)=mb.
\]
This proves the reverse inequality for every $P$.
All multiplicities are unchanged by resolving away from the generic
points in question, so the same argument applies to the normal spaces
under consideration.

The $\varepsilon$ perturbation is necessary: at a pseudoeffective
boundary class, the infimum over its exact positive currents need not
equal its minimal multiplicity. The argument uses that equality only
for the big perturbed classes.
\end{proof}

\begin{lemma}[Pullback when the positive part is nef]
\label{prog:negative-part-pullback}
Suppose
\[
 \xi=P+[F],\qquad P\text{ nef},\quad F\geq0\text{ real Cartier},
 \qquad N_\sigma(\xi)=F.
\]
For every projective resolution $p:W\to S$,
\[
                         N_\sigma(p^*\xi)=p^*F.
\]
\end{lemma}

\begin{proof}
Write $N'=N_\sigma(p^*\xi)$. Since $p^*P$ is nef, $N'\leq p^*F$.
Birational invariance gives $p_*N'=F$, so
$E=p^*F-N'\geq0$ is $p$-exceptional. The positive part of $p^*\xi$ is
\[
                         p^*P+[E],
\]
and is modified nef. If $E\neq0$, exceptional negativity in its
covering-curve form provides a component of $E$ covered by curves $C_t$
contracted by $p$, with $E\cdot C_t<0$
\cite[Appendix~A, Lemma~A.3]{DHY}. A modified nef class has nonnegative
intersection with a general member of a family of curves covering a
prime divisor: use currents with arbitrarily small negative part and
zero generic divisorial Lelong number, restrict to a general member,
and let the negative bound tend to zero. But here
\[
                       (p^*P+E)\cdot C_t=E\cdot C_t<0,
\]
a contradiction. Thus $E=0$.
\end{proof}

\begin{proposition}[Any lower-dimensional good model suffices]
\label{prog:nef-good-comparison}
Suppose the initial relative program has the data in
Lemma~\ref{prog:negative-part-program}. Assume in addition that
its class $\alpha_V$ is nef and agrees with $\mu^*\alpha$ on a common
resolution, and that there is an already projective connected-fibre
morphism $g:V\to S$ with
\[
\alpha_V=g^*\alpha_S,\qquad D_V=g^*D_S.
\]
Suppose there is a projective small modification
$s:S^{\mathrm q}\to S$ such that $S^{\mathrm q}$ is globally Weil
$\mathbb Q$-factorial; the identity is allowed. Assume that the
lower-dimensional adjoint $s^*D_S$ has a good model:
there are a normal compact K\"ahler space $S_m$, a common projective
resolution $p:W\to S^{\mathrm q}$, $q:W\to S_m$, and a nef semiample
class $D_m$ such that
\[
                  p^*s^*D_S=q^*D_m+[E],
 \qquad E\geq0,\quad E\text{ is }q\text{-exceptional}.
\]
Then $\alpha$ is semiample on $X$. If the contraction defining
semiampleness of $D_m$ is Moishezon, the resulting contraction of
$\alpha$ is Moishezon as well. In particular, it is unnecessary to
lift a program on $S$ to $V$, or to require that a chosen program
producing $S_m$ preserve $\alpha_S$ at every intermediate step.
\end{proposition}

\begin{proof}
Nefness descends under the surjective morphism $g$, so $\alpha_S$ is nef.
The class of $F_V$ is pulled back from $S$, so $F_V$ is numerically
trivial over $S$. It is therefore vertical: an effective horizontal
component would restrict to a nonzero effective divisor on a general
projective fibre, with positive intersection against a suitable power
of an ample class, contradicting numerical triviality. For a
birational $g$, verticality is automatic by dimension.
First take a projective resolution $\pi:\widetilde V\to V$ of the main
component of $V\times_S S^{\mathrm q}$. The induced map
$\widetilde g:\widetilde V\to S^{\mathrm q}$ is projective and has
connected fibres. Pull back $\alpha_V,D_V,F_V$ along $\pi$, and pull
back $\alpha_S,D_S$ along $s$. All displayed class identities are
preserved, and $\pi^*F_V$ remains an effective vertical divisor.
Lemma~\ref{prog:negative-part-program} gives
$N_\sigma(D_V)=F_V$ before this modification. Because
$D_V=c\alpha_V+[F_V]$ with $\alpha_V$ nef,
Lemma~\ref{prog:negative-part-pullback} then gives
\[
              N_\sigma(\pi^*D_V)=\pi^*F_V.
\]
We may therefore replace $(V,S,g)$ by
$(\widetilde V,S^{\mathrm q},\widetilde g)$ and suppress the new
superscripts in the rest of the proof. The good-model comparison now
reads $p^*D_S=q^*D_m+[E]$. No assertion that the crepant subboundary on
$\widetilde V$ is effective is needed: this space is used only for
classes, currents and the effective divisor $\pi^*F_V$. The
lower-dimensional generalized pair is the crepant pullback to the
small model $S^{\mathrm q}$.

The class identity $[F_V]=g^*(D_S-c\alpha_S)$ shows that $F_V$ is
numerically trivial over $S$. Lemma~\ref{prog:vertical-divisor}
gives an actual effective real Cartier divisor $F_S$ with
$F_V=g^*F_S$. Injectivity of pullback yields
\[
                         D_S=c\alpha_S+[F_S].
\]
The established identity $N_\sigma(D_V)=F_V$ and
Lemma~\ref{prog:negative-part-base} therefore give
$N_\sigma(D_S)=F_S$, and
Lemma~\ref{prog:negative-part-pullback} gives
\[
                         N_\sigma(p^*D_S)=p^*F_S.
\]
On the other hand $q^*D_m$ is nef, so the exceptional-divisor identity
applied to the good-model comparison gives
\[
                         N_\sigma(p^*D_S)=E.
\]
Thus $E=p^*F_S$ as actual divisors. Subtracting them from the comparison
proves the exact class equality
\[
                         q^*D_m=c\,p^*\alpha_S.
\]
This proves the needed preservation of $\alpha_S$ from the final good
model, without an assumption about its intermediate models.

Choose a contraction $h:S_m\to T$ to a normal compact K\"ahler space
and a K\"ahler class $\kappa$ on $T$ with $D_m=h^*\kappa$. Taking a
resolution of the main component of $V\times_S W$, and then a common
projective resolution with $U$, produces a projective modification
$r:R\to X$ and a holomorphic map $\ell:R\to T$ satisfying
\[
                         r^*\alpha=\frac1c\ell^*\kappa.
\]
For every curve $C$ in a fibre of $r$ this equality implies
$\ell(C)$ is a point, since a K\"ahler class has positive degree on
every nonconstant image curve. The fibres of the projective modification
$r$ are connected projective complex spaces, hence are connected by
chains of curves. Therefore $\ell$ is
constant on each fibre of $r$. The factorization theorem for a proper
map onto a normal space gives a holomorphic map $f:X\to T$ with
$\ell=f\circ r$. Pushing forward the class equality by $r$ gives
\[
                              \alpha=f^*(\kappa/c).
\]
The maps from the main fibre-product component to $S_m$ have connected
general fibres; their Stein factorizations are finite birational over
the normal space $S_m$, hence have connected fibres everywhere.
Together with the connected fibres of $h$, this shows that $\ell$,
and therefore $f$, has connected fibres. This is the required
semiampleness of $\alpha$.

Finally, $R\to S_m$ is projective: it is obtained from the projective
map $g$ by base change, followed by projective resolutions and the
projective map $q$. If $h$ is Moishezon, choose a projective surjection
$H\to S_m$ such that $H\to T$ is projective. A component of
$R\times_{S_m}H$ dominating $R$ is projective and surjective over $X$,
and its map to $T$ is projective. This is a Moishezon witness for $f$.
\end{proof}


\begin{lemma}[The contraction step in the dimension induction]
\label{prog:base-point-free-induction}
Assume $\mathsf C_{d-1}$, $\mathsf B_{d-1}$ and $\mathsf M_{d-1}$.
Then $\mathsf C_{d,\mathrm{big}}$ and $\mathsf B_d$ hold.
\end{lemma}
\begin{proof}
\smallskip\noindent
\emph{The big non-klt contraction.}
Apply the construction of \cite[Section~3]{HX}. On its smooth
modification $\nu:U\to X$ it writes
\[
 \nu^*\alpha=[D_U]+\eta_U,
 \qquad D_U\geq0,\quad \eta_U\text{ K\"ahler}.
\]
The induction over the jumping coefficients of the multiplier
ideal reduces the new reduced stratum to dimension at most $d-1$.
The given contraction on the old non-klt subspace and
$\mathsf C_{d-1}$ provide the Moishezon maps of those strata.
Every map that is promoted to a projective map in that construction
contracts exactly the curves on which $\nu^*\alpha$ vanishes.
On each such curve the actual real line bundle $-D_U$ has degree
$\eta_U\cdot C$. On a further projective common resolution
$q:U''\to U$, choose an effective exceptional divisor $E$ with
$-E$ relatively ample and choose $\varepsilon>0$ sufficiently small.
Use the decomposition
\[
 q^*\nu^*\alpha=[D'']+\eta'',\qquad
 D''=q^*D_U+\varepsilon E,\qquad
 \eta''=q^*\eta_U-\varepsilon[E].
\]
Here $\eta''$ is K\"ahler, and the actual real line bundle $-D''$
has degree $\eta''\cdot C$ on every contracted curve. Restrict this
decomposition to the smooth reduced components in the gluing
construction. Lemma~\ref{prog:actual-line-polarization} then supplies
precisely the relative ampleness needed in place of
\cite[Lemma~2.42]{HX}, including the common-resolution step in its
Claim~3.5. The uncorrected pullback $q^*\eta_U$ is not being treated
as a K\"ahler class.

The exact sequences of multiplier ideals, relative vanishing,
finite pushouts, and extension from sufficiently thickened
$\operatorname{Supp}D_U$ in that proof then apply unchanged.
They give the birational Moishezon contraction in
$\mathsf C_{d,\mathrm{big}}$. In particular the gluing step uses
neither generalized termination nor a projectivity criterion for
an undetected ray.

\smallskip\noindent
\emph{Preparation for both cases of $\mathsf B_d$.}
The modified-big perturbation
\cite[Lemma~2.23]{HX} supplies a nef datum that dominates a
K\"ahler class on its carrier. Subtracting a sufficiently small
multiple of the pullback of a K\"ahler class $\omega_X$ still leaves
a nef datum on that carrier. Thus
$\alpha$ is a gklt adjoint plus $\varepsilon\omega_X$, and the cone
argument of \cite[Lemma~2.45]{HX} makes $\alpha$ NQC.
This preparation applies whether or not $\alpha$ is big.

\smallskip\noindent
\emph{The nef-and-big case of $\mathsf B_d$.}
For a gklt pair the multiplier ideal is the unit ideal, so the
preceding big non-klt assertion now applies and gives a birational
contraction $h:X\to Y$.
On a projective resolution $p:W\to X$ chosen projective over $Y$,
relative Kawamata--Viehweg vanishing for the effective exceptional
divisor $-\lfloor B_W\rfloor$ gives
$R^ih_*\mathcal O_X=0$ for $i>0$. Thus $Y$ has rational
singularities. Proper birational Bott--Chern descent
\cite[Lemma~8.7]{DHP} gives $\alpha=h^*\gamma$.
The descended adjoint is gklt, nef and big, and has no trivial
curve. The criterion \cite[Theorem~4.3]{HLX} makes $\gamma$
K\"ahler.

For completeness, the projectivity input in that criterion is
restricted to a map $S\to Z'$ between smooth compact K\"ahler
manifolds with rationally connected general fibre: $S$ is the
chosen smooth divisor on a resolution and $Z'$ is the resolution
of the null-locus component. The smooth-source, smooth-base
argument in \cite[Theorem~3.1, Step~1]{CH} applies. Its subsequent
relative program starts with this projective map. Thus this application of the criterion uses only the
smooth-source, smooth-base projectivity argument and an
already-projective relative program.

\smallskip\noindent
\emph{The non-big case of $\mathsf B_d$.}
Use a projective small factorialization and the modified-big
perturbation of the pair. The non-pseudo-effectivity of $K_X$
gives an MRC fibration by \cite{Ou2025}. Choose a smooth K\"ahler
MRC base $Z$ and a smooth modification $\mu:U\to X$ such that
$U\to Z$ is projective. Only the smooth case of
\cite[Theorem~3.1, Step~1]{CH} is needed: pullback identifies
the holomorphic two-forms since the general fibre is rationally
connected. Lemma~\ref{prog:relative-hodge} identifies any $(1,1)$-class on $U$,
modulo a class pulled back from $Z$, with an actual real line bundle.
Lemma~\ref{prog:projective-cone} identifies its relative rays
with rays of the analytic cone; both statements persist along the
projective relative construction.

Write, as in \cite[Theorem~4.1, Claim~4.1]{HX},
\[
 D_U(a)=K_U+\Delta_U+\mathbf M_U+a\alpha_U
     =(a+1)\alpha_U+F_U,
 \quad \alpha_U=\mu^*\alpha,\quad
 \Delta_U,F_U\geq0,
\]
where $\Delta_U$ is klt and $F_U$ is $\mu$-exceptional.
The modified-big replacement is chosen, as in the cited Claim~4.1,
so that the nef datum $\mathbf M_U$ on this smooth carrier is
K\"ahler. If a further projective resolution is needed, subtract a
sufficiently small effective exceptional divisor from its nef part
and add that divisor to the subboundary; this preserves the
adjoint and the klt inequalities.
Choose $a_0$ using the uniform bound of
Proposition~\ref{prog:relative-nqc-grid}. It makes the relative
$D_U(a_0)$-program $\alpha_U$-trivial and preserves that choice
through every step. This is a
\emph{projective} relative program from its first step: the nef
data are represented by real line bundles modulo $Z$.
Lemma~\ref{prog:relative-hodge} applies on every
intermediate model over this fixed smooth base.
More precisely, put $r:U\to Z$ and choose one actual global real
line bundle $L$ with
\[
 \mathbf M_U+a_0\alpha_U=c_1(L)+r^*\gamma.
\]
The bundle $L$ is relatively ample, because $\mathbf M_U$ is
K\"ahler and $\alpha_U$ is nef. On each of a fixed finite collection
of relatively compact Stein charts of $Z$, choose an effective
real divisor $\Theta$ representing $L$, with $(U,\Delta_U+\Theta)$
klt. The representatives can retain a positive relatively ample
part in their boundaries. Here the equality
$K_U+\Delta_U+\Theta\sim_{\R}K_U+\Delta_U+L$ is an actual
real linear equivalence of line bundles; the equality with
$D_U(a_0)$ modulo $r^*\gamma$ is separately an equality of
Bott--Chern classes.

Fix these ordinary pairs on the initial charts. Their actual
linear equivalences and their Bott--Chern comparisons with $D_U(a_0)$
modulo the fixed base class persist separately throughout a global
relative program. Proposition~\ref{prog:projective-good-model} applies:
the source is smooth, its adjoint is represented by an actual global
real line modulo $Z$, its carrier nef part is relatively ample, and
$D_U(a_0)=(a_0+1)\mu^*\alpha+[F_U]$ is pseudo-effective.
It gives a finite program $U\dasharrow V$ and a good relative endpoint.
This is one global program controlled by fixed ordinary chart pairs.
Its construction uses global line algebras, not a choice of unrelated
local minimal models. Proposition~\ref{prog:relative-nqc-grid} makes
every step $\alpha_U$-trivial with the same value $a_0$.

Here is an explicit ample-model comparison that gives the required
descent of $\alpha_V$. Both $D_V(a_0)$ and $\alpha_V$ are nef over
$Z$. Put $a_1=a_0+1$ and $a_2=a_0+2$. Before the program starts,
choose fixed ordinary representatives on the finite chart cover for
all three relatively ample nef parts $\mathbf M_U+a_j\alpha_U$,
$j=0,1,2$, on one simultaneous log resolution. The
$\alpha_U$-crepancy of the constructed program makes its endpoint
$V$ a weak model for each of these three fixed ordinary adjoints.
Their traces $D_V(a_j)=D_V(a_0)+(a_j-a_0)\alpha_V$ are nef.
Corollary~\ref{prog:projective-semiample} therefore makes
\[
 D_V(a_j)=D_V(a_0)+(a_j-a_0)\alpha_V
\]
semiample over $Z$. Their zero curves are exactly the common zero
curves of $D_V(a_0)$ and $\alpha_V$. Their projective ample-model
contractions therefore have the same fibres: these fibres are
connected by curves, and each contraction is constant on the
fibres of the other. Identify the two normal targets, writing the
common contraction as $g:V\to S$. If
$D_V(a_j)=g^*\lambda_j$, with $\lambda_j$ relatively K\"ahler over
$Z$, subtraction gives
\[
 \alpha_V=g^*(\lambda_2-\lambda_1)=g^*\alpha_S.
\]
This is an equality of Bott--Chern classes, not merely of curve
degrees. Since the initial program is $\alpha_U$-trivial, it is
also $D_U(a_1)$-negative. From now on put $a=a_1$ and suppress
the argument in $D_U(a)$ and $D_V(a)$.

The maps $V\to Z$ and $S\to Z$ are projective. Moreover $g$ is
projective: a relatively ample bundle for $V\to Z$ restricts to
an ample bundle on every fibre of $g$, and hence is $g$-ample.
Each $D_U(a_j)$ is not big globally: pushing a big such class to
$X$ would make $(a_j+1)\alpha$ big. If it were big over the smooth
MRC base $Z$, relative-canonical positivity
\cite[Theorem~2.48]{HX}, together with the pseudo-effectivity of
$K_Z$, would make it big globally. It is therefore not big over
$Z$. This property persists to its relative semiample model, so
its relative canonical morphism has $\dim S<d$.
The projective canonical bundle formula
\cite[Theorem~2.53]{HX} provides
\[
 D_S=K_S+B_S+\mathbf N_S+a\alpha_S,
 \qquad D_V=g^*D_S,
\]
with a gklt pair on $S$ and modified-big nef data.
The class $\alpha_S$ is nef by descent of nefness through the
surjective projective map $g$.
Thus adding $a\overline{\alpha_S}$ to the nef datum preserves the
gklt condition and modified bigness: on a common carrier this
addition is a nef pullback and preserves the existing big class.

Take a projective small factorialization $s:S^{\mathrm q}\to S$
and resolve the main component of $V\times_S S^{\mathrm q}$.
The resulting maps $\pi:\widetilde V\to V$ and
$\widetilde g:\widetilde V\to S^{\mathrm q}$ are projective,
and $\widetilde V$ is compact K\"ahler. Pull back $D_V$,
$\alpha_V$ and $F_V$ to $\widetilde V$, and $D_S$ and
$\alpha_S$ to $S^{\mathrm q}$. The pair on $S^{\mathrm q}$ is
gklt and its boundary-plus-nef class remains modified big.
Lemma~\ref{prog:negative-part-program} first gives
$N_\sigma(D_V)=F_V$. Lemma~\ref{prog:negative-part-pullback} then gives
$N_\sigma(\pi^*D_V)=\pi^*F_V$. The new total space is used only as a
carrier for this equality and for pullbacks of positive currents;
no effective boundary on it is needed. Rename these spaces $V$
and $S$. No additional relative program is needed for this modification.

The vertical-divisor descent and negative-part lemmas above now
give actual effective divisors $F_S,F_V$ with
\[
 F_V=g^*F_S,\qquad
 D_S=(a+1)\alpha_S+[F_S],\qquad N_\sigma(D_S)=F_S.
\]
Apply $\mathsf M_{d-1}$ to the generalized pair with adjoint $D_S$.
On a common resolution $p:W\to S$, $q:W\to S_m$ of the resulting
good model, write
\[
 p^*D_S=q^*D_m+[E],\qquad E\geq0
 \text{ exceptional over }S_m.
\]
The negative-part comparison proved above gives
$E=p^*F_S$, and hence
\[
 q^*D_m=(a+1)p^*\alpha_S.
\]
Since $D_m$ is semiample, $\alpha_S$ is semiample after descent
through $p$. Pulling back by $g$ and using the
$\alpha$-trivial initial program gives semiampleness of
$\alpha$ on $X$. The target is compact K\"ahler, and the
resulting map is Moishezon, as follows by taking common projective
modifications of the maps just constructed. The comparison uses only the chosen lower good model; no
lower-dimensional program is lifted to $V$.

\end{proof}

\begin{proposition}[A negative ray with an actual detector]
\label{prog:ordinary-step}
Assume $\mathsf B_d$. Let $A$ be a gklt adjoint in dimension at most
$d$ on a globally strongly $\mathbb Q$-factorial compact K\"ahler
space, with globally nef carrier data. Let $R$ be an $A$-negative
extremal ray of the full analytic cone. If an actual global rational
line $L$ has $L\cdot R<0$, the ray has a projective contraction;
a small contraction has its detected flip. The working models retain
the strong factoriality and compact K\"ahler properties.
In particular this applies to an ordinary rational adjoint
$A=c_1(K_X+B)$ with detector $K_X+B$, and to ordinary real boundaries
by Proposition~\ref{prog:ordinary-replacement}.
\end{proposition}
\begin{proof}
The local polyhedrality of the negative cone makes $R$ an exposed ray.
Choose a nef support $\lambda$ whose null analytic face is $R$.
Normalize the analytic cone by a K\"ahler class. On a neighbourhood
of the point representing $R$, the class $-A$ is positive. On the
remaining compact part of this slice, $\lambda$ has a positive
minimum. Consequently, for a sufficiently large $b$, the class
$\omega_R=b\lambda-A$ is K\"ahler. Thus
$A+\omega_R=b\lambda$ is nef with precisely the null face $R$.
The gklt presentation with nef carrier datum increased by the pullback
of $\omega_R$ has modified-big total boundary. Apply $\mathsf B_d$
to construct its contraction. Lemma~\ref{prog:detected-projectivity}
makes the same global line $-L$ relatively ample, and
Proposition~\ref{prog:detected-flip} constructs the flip when needed.
For a divisorial contraction the exceptional-prime negativity and
integral line descent give strong factoriality on the target: pull
back a rank-one reflexive sheaf, remove its degree by a rational
multiple of the exceptional divisor, descend the resulting line,
and compare on the common big open. Local ordinary log-Fano
replacement gives rational singularities by relative vanishing.
The supporting target and projective flip are compact K\"ahler.
\end{proof}

\subsubsection{Relative good models and finite geography}

\begin{lemma}[Removing the preparation error]
\label{prog:remove-preparation}
Let $\mu:W\to X$ be a projective smooth preparation over $S$ with
$A_W=\mu^*A_X+[F]$, where $F\geq0$ is $\mu$-exceptional and has
positive coefficient on every added prime. If $W\dashrightarrow Y$
is a chosen good log terminal model of $A_W$ over $S$, then it induces
a good log terminal model of $A_X$ over $S$.
\end{lemma}
\begin{proof}
On a common projective resolution $p:V\to W$, $q:V\to Y$, write
$p^*A_W=q^*A_Y+[E]$ with $E\geq0$ exceptional over $Y$.
Put $r=\mu p$ and $D=E-p^*F$. Then
\[
 [D]=r^*A_X-q^*A_Y,\qquad r_*D=r_*E\geq0.
\]
The divisor $-D$ is $r$-nef because $A_Y$ is nef over $S$.
Negativity gives $D\geq0$, hence $E\geq p^*F$. Every added prime
is therefore exceptional over $Y$. This proves nonextraction from
$X$ and gives its effective exceptional comparison with $A_Y$.
Strictness at any contracted prime of $X$ follows from strictness
for the chosen log terminal model of $W$, since $F$ has coefficient
zero there. The same semiample endpoint makes the model good.
\end{proof}

\begin{lemma}[Recovering the uncontracted original primes]
\label{prog:weak-to-terminal}
Let a gklt pair on $X$ have a nonextracting weak good model
$X\dashrightarrow Y$ over $S$. Assume $Y$ is compact K\"ahler,
globally strongly $\mathbb Q$-factorial and has the exceptional
splitting of Lemma~\ref{prog:detected-cohomology}. Then a projective
crepant extraction $Y'\to Y$ gives a good log terminal model of the
original pair. The exceptional splitting persists on $Y'$.
\end{lemma}
\begin{proof}
On a common resolution write the actual comparison
$p^*A_X=q^*A_Y+[E]$, where $E\geq0$ is $q$-exceptional.
There are finitely many primes of $X$ contracted by its marking.
Let $\mathcal P$ consist of those with coefficient zero in $E$.
Take a projective log resolution $r:W\to Y$ which contains all
these valuations and dominates the fixed nef carrier. In the crepant
structure boundary $B_W$, the coefficient of each member of
$\mathcal P$ is its original effective boundary coefficient, hence
lies in $[0,1)$. Keep these coefficients unchanged. Increase each
other $r$-exceptional coefficient, if necessary from a negative
value, to a number in $(\max\{0,b_Q\},1)$. Keep all nonexceptional
coefficients unchanged. The resulting effective klt boundary
$\Gamma_W$ satisfies
\[
 [K_W+\Gamma_W]+\mathbf M_W=r^*A_Y+[F],
\]
where $F\geq0$ has precisely the unwanted exceptional primes as its
support. The inherited splitting gives actual real-line representatives
for every class modulo $Y$. Since $r$ is birational, the nef carrier
datum is relatively big; the adjoint is relatively represented by
$F$ and is pseudo-effective. Apply
Proposition~\ref{prog:projective-good-model} over $Y$.
Its projective endpoint $r':Y'\to Y$ has an effective relatively nef
exceptional trace $F'$, so negativity gives $F'=0$. A prime outside
$\operatorname{Supp}F$ cannot be contracted by an $F$-negative
birational program: on a general curve through its generic point the
effective divisor has nonnegative degree. Thus all members of
$\mathcal P$ survive, and all unwanted exceptional primes disappear.
No new prime is extracted by the program. It follows that $r'$ is
crepant and that the marking from $X$ contracts exactly primes with
strictly positive original comparison coefficient. It is nonextracting
and is therefore a log terminal model. The pullback of the semiample
adjoint on $Y$ makes it good. Finally apply
Corollary~\ref{prog:extraction-cohomology}.
\end{proof}

\begin{lemma}[A detector throughout the big-adjoint construction]
\label{prog:big-good-model}
Assume $\mathsf B_d$ and $\mathsf F_{d-1}$. The big-adjoint
construction of \cite[Section~5.1]{HX} gives a chosen good model in
dimension $d$ using only detected steps. It also gives the case of
$\mathsf M_d$ in which $A+c f^*\omega_S$ is big for some $c>0$,
with the map to $S$ retained throughout.
\end{lemma}

\begin{proof}
After the smooth preparation in \cite[Proposition~5.12]{HX}, write the
adjoint as
\[
 A=D+\omega,\qquad D\geq0,\quad \omega\text{ K\"ahler},
 \qquad \operatorname{Supp}D=N_+(A).
\]
Here $N_+(A)$ denotes the reduced divisor with support
$\bigcap_{0<\epsilon\ll1}\operatorname{Supp}N_\sigma(A-\epsilon\omega)$;
this support is constant for sufficiently small positive $\epsilon$
and is independent of the chosen Kähler class $\omega$
\cite[Lemma~5.3]{HX}. At a nonterminal scaling step,
\[
 A_i\cdot R_i<0,\qquad
 (A_i+t_i\omega_i)\cdot R_i=0,\qquad t_i>0.
\]
Consequently $\omega_i\cdot R_i>0$ and $D_i\cdot R_i<0$.
A component of the actual effective divisor $D_i$ supplies a global
rational Cartier detector. The supporting contraction follows from
$\mathsf B_d$, and its projectivity and flip follow from
Proposition~\ref{prog:detected-flip}. This argument does not require
$\omega_i$ to remain nef.

To apply $\mathsf B_d$ to a gdlt presentation, lower its finitely many
floor coefficients on the initial smooth carrier and add the same small
class to the K\"ahler nef datum. The adjoint does not change, and the
resulting gklt presentation persists through the same negative steps.
The original gdlt presentation is retained for adjunction to the floor.
Thus the special-termination proof of \cite[Lemma~5.8]{HX} applies with
$\mathsf F_{d-1}$ on its smooth lower-dimensional carrier. Its other
inputs are adjunction, discrepancy monotonicity, and the local
Cartier-index calculation. The last neighbourhood-isomorphism step
uses the evaluation-ideal support argument in
the \hyperref[prog:special-support]{support comparison above}: the comparison of
the first detected flip has support equal to the full inverse image
of its exceptional locus, and later comparisons dominate it.
One does not infer disjointness of images merely from a vanishing
restriction of an arbitrary exceptional divisor. The theta induction of
\cite[Proposition~5.11]{HX} now gives a weak model. Applying
$\mathsf B_d$ to its nef adjoint makes it good, and
Lemma~\ref{prog:weak-to-terminal} gives the log terminal model.
All its steps are detected or belong to the already-projective
relative extraction in that lemma.

For completeness, the negative-part identity needed in this argument
is only
\[
 N(p^*\xi+E)=N(p^*\xi)+E\quad(E\geq0\text{ $p$-exceptional}),
\]
not the stronger formula $N(p^*\xi)=p^*N(\xi)$ for an arbitrary
pseudo-effective class. The comparison with a nef endpoint identifies
its exceptional error with the negative part. The formula for
$N_+$ follows by writing a K\"ahler form upstairs as
$p^*\omega-F$, with $F$ positive on every exceptional prime, and applying
the displayed identity to
$p^*(\xi-\epsilon\omega)+E+\epsilon F$. Pushforward identifies the
nonexceptional coefficients; $\epsilon F$ supplies all exceptional
primes. These are exactly the uses in the theta argument.

Now put $P=f^*\omega_S$ and suppose $A+cP$ is big for some $c$.
Fix $\delta>0$ smaller than the $\omega_S$-degree of every compact
integral curve on $S$ and enlarge $c$ so $c>4d/\delta$.
Use the following direct preparation, which keeps the original
unshifted globally nef datum. On a smooth projective carrier
$\nu:W\to X$ write
\[
 \nu^*(A+cP)=[G]+\omega,\qquad G\geq0,\quad\omega\text{ K\"ahler},
 \qquad [K_W+B_W]+M_W=\nu^*A,
\]
with the non-K\"ahler support property used in the big-class
preparation. Here $\operatorname{Ex}(\nu)$ below denotes the reduced
divisor formed by all $\nu$-exceptional primes. Write
$B_W=B_W^+-B_W^-$. For small $\epsilon>0$ put
\[
 \begin{split}
 B'&=B_W^++\epsilon B_W^-+
                   \epsilon(G+\operatorname{Ex}(\nu)),\\
 M'_W&=M_W+c\nu^*P+\epsilon\omega,\\
 b&=B_W^++\frac{\epsilon}{1+\epsilon}\operatorname{Ex}(\nu).
 \end{split}
\]
The new boundary is klt and $M'_W$ is K\"ahler. Direct addition gives
\[
 A'=[K_W+B']+M'_W
   =(1+\epsilon)\bigl([K_W+b]+M_W+c\nu^*P\bigr).
\]
Moreover $B'\geq b$. In every theta branch the cumulative added
boundary satisfies $0\leq C\leq1-B'$, so
\[
 b+\frac{C}{1+\epsilon}
 \leq b+\frac{1-b}{1+\epsilon}<1.
\]
Thus the unshifted presentation
$[K_W+b+C/(1+\epsilon)]+M_W$ is genuinely gklt with the
\emph{original} globally nef datum. This is the presentation used
for the horizontal length bound. No base form is subtracted from
an arbitrary newly prepared nef datum.

A horizontal negative ray would have $P$-degree at least $\delta$;
the bound $0<-A_{\rm unshifted}\cdot C\leq4d$ would contradict
$c\delta>4d$. Every constructed step is therefore over $S$.
Its projective connected fibres are contracted by $f$, so $f$ factors
through its base and through the flip. The final finite extraction is
also over $S$.

The shifted nef endpoint is semiample by $\mathsf B_d$. Its zero face
contains no horizontal curve: decompose such a curve by the unshifted
adjoint's cone theorem; nefness of the shifted class forces all
summands into the zero face, and a summand of positive $P$-degree again
contradicts the length bound. The shifted class is big, so the
connected semiampleness map is bimeromorphic. Its connected fibres are Moishezon and are connected
by chains of compact curves. The preceding zero-face argument makes
the map to $S$ constant on each such curve, hence on every fibre.
The analytic factorization lemma therefore factors the map to $S$
through the semiampleness map. Subtracting $c\omega_S$ from the
normalized class on its target gives the required relative K\"ahler class for the unshifted adjoint. The effective
resolution error from the preparation is removed by the usual
negativity comparison, giving a model of the original data.

The forward traces and exceptional splittings persist by
Lemma~\ref{prog:detected-cohomology}. In the last extraction,
the newly extracted prime classes move from the old exceptional span
into the pullback span; they are actual rational Cartier classes.
Both-sign negativity proves directness as before.
\end{proof}

\begin{lemma}[Changing the base of a relatively trivial adjoint]
\label{prog:graph-good-model}
Suppose $g:X\to Y$ is projective, its general fibre is rationally
connected, and a gklt adjoint satisfies $A_X=g^*A_Y$. Assume the
prepared nef datum is K\"ahler on a smooth carrier. Let
$Y\dashrightarrow Y_m$ be the chosen nonextracting good model
supplied by $\mathsf M_{d-1}$ over $S$, with its exceptional
cohomology splitting, and let $\dim Y<d$. Then a single projective relative construction gives a weak good model of $(X,A_X)$ over $S$ with that splitting.
Lemma~\ref{prog:weak-to-terminal} makes it log terminal. No step
of the base program is lifted.
\end{lemma}

\begin{proof}
Take a common projective smooth resolution $T\to Y$, $T\to Y_m$
and resolve the main component of $X\times_Y T$, further dominating
the prepared nef carrier; denote the resulting smooth space by $V$.
Its morphism $h:V\to Y_m$ is projective. Before constructing a
program, verify algebraicity for all classes. The projective RC map
$V\to T$ and Lemma~\ref{prog:relative-hodge} express every class
modulo $T$ by actual real lines. The splitting of the chosen $Y_m$
expresses classes from $T$, modulo $Y_m$, by exceptional divisor
classes. Hence
\[
 H^{1,1}_{\rm BC}(V)=c_1(\operatorname{Pic}(V)\otimes\mathbb R)
                         +h^*H^{1,1}_{\rm BC}(Y_m).
\]
Lemma~\ref{prog:projective-cone} identifies the relative cone
with the full analytic face. These properties and the exceptional
splitting persist along the projective relative steps by
Lemma~\ref{prog:detected-cohomology}.

Choose the effective gklt preparation on this final carrier. Its
adjoint has the identity
\[
 A_V=h^*A_{Y_m}+[G],\qquad G=G_{\rm base}+F\geq0,
\]
where $G_{\rm base}$ is pulled back from the actual effective base
comparison and $F$ is exceptional over $X$, with positive coefficient
on every added prime. Its nef datum is nef and big on the prepared
carrier. Modulo $Y_m$ it is the actual real-line class
$[G]-c_1(K_V)-[B_V]$. Thus every hypothesis of
Proposition~\ref{prog:projective-good-model} holds, with the effective
relative adjoint $G$. Denote its global relative endpoint by $V_m$
and the strict transform of $G$ by $G_m$.

On a common resolution $p:W\to V$, $q:W\to V_m$, the actual adjoint
comparison has effective $q$-exceptional error $E$. The difference
$p^*G-q^*G_m-E$ has zero Bott--Chern class and is $q$-exceptional,
because $q_*p^*G=G_m$. Negativity in both signs makes it zero:
\[
 p^*G=q^*G_m+E.                                      \tag{*}
\]
Choose a big common-isomorphism open $U\subset Y_m$ for
$Y\dashrightarrow Y_m$. Such an open exists because that map is
nonextracting. On $U$ we have $G_{\rm base}=0$. Let
$r:W_U\to X_U$ be the induced projective birational map. The divisor
$q^*G_m$ is $r$-nef, since the endpoint adjoint is nef over $U$ and
the base class has zero degree on $r$-fibres. Equation $(*)$ gives
\[
 r_*q^*G_m=-r_*E\leq0,
\]
since $F$ is exceptional over $X$. Negativity yields $q^*G_m\leq0$.
Effectivity yields $G_m|_U=0$. This step removes the possible
\emph{horizontal} components of $F$; they were not assumed very
exceptional at the start.

Now $G_m$ is supported over codimension at least two in $Y_m$.
It is nef over $Y_m$, so very-exceptional negativity gives $G_m=0$.
The endpoint adjoint is the pullback of $A_{Y_m}$, hence is nef and
good over $S$. Its relative canonical morphism is Moishezon: it is
the connected factor of the composite of the projective morphism to
$Y_m$ and the lower-dimensional Moishezon canonical morphism.
Lemma~\ref{prog:remove-preparation} removes the added
source-exceptional errors and proves nonextraction. All these operations are over $S$.

The all-class invariant was established before the relative program;
its preservation follows at each step from the same cone identification
and forward cohomology splitting. If the construction has contracted
any original prime crepantly, apply
Lemma~\ref{prog:weak-to-terminal} to retain it.
\end{proof}

\begin{lemma}[The global rational quotient over a fixed base]
\label{prog:relative-model}
Assume $\mathsf M_{d-1}$ and the case of $\mathsf M_d$ in which a
base shift makes the adjoint big. Then $\mathsf M_d$ holds when no
class $A+c f^*\omega_S$, $c\geq0$, is big.
\end{lemma}

\begin{proof}
First take the usual smooth carrier with effective gklt boundary
and adjoint $\mu^*A+[F_0]$, where $F_0$ is effective exceptional
with full exceptional support. Its boundary-plus-nef trace is big and its adjoint
is not big, so its canonical class is not pseudo-effective. Take
its global MRC and resolve the MRC graph and the fixed carrier,
obtaining a projective map to a smooth compact K\"ahler $Z$ with
$K_Z$ pseudo-effective.
Perform the K\"ahler-nef-part preparation \emph{after} these graph
resolutions, preserving the maps to $Z$ and $S$. This order ensures
that the final scaling datum is K\"ahler, rather than merely the nef
pullback of an earlier K\"ahler datum. Write the resulting adjoint as
\[
 A_0=K_U+B_U+M_U=\mu^*A+[F],
 \qquad F\geq0\text{ exceptional},\qquad M_U\text{ K\"ahler}.
\]
The global modified-big hypothesis makes a preliminary base shift
unnecessary. Relative pseudo-effectivity is preserved.
No base shift of $A_0$ is big, since pushing a big such class down
would make a base shift of $A$ big. The fixed morphism $h:U\to Z$
is projective and has RC general fibres. All these preparations
precede the next, fresh base shift.

Put $P=f_U^*\omega_S$ and choose $c>2d/\delta$, where every
integral curve on $S$ has $\omega_S$-degree greater than
$\delta>0$. Set
\[
 A_c=A_0+cP=D+N,\qquad D=K_U+B_U,\quad N=M_U+cP.
\]
Since $A_c$ is not big globally, relative-canonical positivity
\cite[Theorem~2.48(2)]{HX}, together with pseudo-effectivity of $K_Z$,
shows that $A_c$ is not big over $Z$, whether or not it is globally
pseudo-effective. The class $N$ is big over $Z$. Consequently $D$
cannot be pseudo-effective over $Z$: otherwise $D+N=A_c$ would be
big over $Z$. The smooth projective RC Hodge projection represents
every class modulo $Z$ by an actual real line. The ordinary global klt adjoint $D$ is therefore relatively
non-pseudo-effective. The projective ordinary Mori construction of
Corollary~\ref{prog:projective-mori}, with scaling of the actual
relatively ample real line represented by $N$, terminates in a Mori
fibre space. Lemma~\ref{prog:relative-hodge} and
Lemma~\ref{prog:projective-cone} identify every relative ray
with a full analytic ray. Its nef b-data retain the fixed globally
nef carrier $U$.

First suppose $A_c$ is not pseudo-effective globally. The retained
relative-canonical positivity theorem on $h$, together with
pseudo-effectivity of $K_Z$, then says that $A_c$ is not
pseudo-effective over $Z$. The terminating relative program ends
in a Mori contraction at a threshold $\tau>1$. At every step its
scaled zero ray has $A_c$-degree negative. Inductively the map to
$S$ is retained: if such a ray were horizontal, its $P$-degree
would exceed $\delta$, while the gklt cone theorem for $A_0$ gives
a rational generator of $-A_0$-degree at most $2d$, contradicting
$A_c\cdot R\leq0$. Once $P\cdot R=0$, the step is negative or
crepant for $A_0$, so its gklt presentation persists. The same
argument applies to the final Mori ray. This would be an
$A_c$-negative Mori contraction over $S$, contradicting the
preserved pseudo-effectivity over $S$. Hence $A_c$ is
pseudo-effective globally.

It is not big globally. Relative-canonical positivity now implies
that it is not big over $Z$. Its pseudo-effective threshold for
$D+tN$ over $Z$ is exactly one: for $t>1$ the class is big, while
pseudo-effectivity for $t<1$ would make $A_c$ big because $N$ is
big over $Z$. The same finite relative program therefore ends in
a Mori contraction $g:U'\to Y$ at threshold one.
All birational steps, including any steps at threshold one, and
the final zero ray are over $S$ by the preceding length argument
with $A_c\cdot R\leq0$. For a zero ray with positive $P$-degree,
$A_0$ would again have degree less than $-c\delta$.
Thus $Y$ carries a map to $S$ and $\dim Y<d$.
The class $A_c$ has zero degree on every $g$-contracted curve.
This is an ordinary projective log-Fano contraction: rationalize its
effective real ordinary boundary, preserving klt and relative
antiampleness. Relative vanishing gives rational singularities on
$Y$, which is compact K\"ahler because it is projective over $Z$.
Lemma~\ref{prog:bc-descent} therefore gives an actual equality
$A_c=g^*A_Y$ in Bott--Chern cohomology. This descent precedes the
projective canonical bundle formula. That formula now provides gklt
data for $A_Y$, globally nef on a carrier, whose total boundary is
globally modified big. Relative pseudo-effectivity descends through
the projective surjection $g$, so $\mathsf M_{d-1}$ applies over $S$.
Keep the base shift in this application. Use its chosen output in
Lemma~\ref{prog:graph-good-model}.

Steps at threshold one can be $A_c$-crepant divisorial contractions.
Consequently the composite constructed so far need only be a weak
good model of the original prepared adjoint. Apply
Lemma~\ref{prog:weak-to-terminal} to extract the finitely many original
primes with zero comparison coefficient. This gives the strict log
terminal model and preserves the cohomology invariant. The shift
changes neither relative comparisons nor relative nefness.
Lemma~\ref{prog:remove-preparation} removes the full-support
preparation error and proves the claimed $\mathsf M_d$.
\end{proof}

\subsubsection{Finite geography and closing the induction}

\begin{lemma}[Canonical models first, weak models second]
\label{prog:finite-geography}
Assume $\mathsf B_d$ and $\mathsf M_d$. Then $\mathsf F_d$ holds for compact polytopes
of globally nef b-data on a fixed carrier and globally modified-big
boundary, over the fixed base.
\end{lemma}

\begin{proof}
First prove, by induction on the parameter-polytope dimension,
polyhedrality of the effective locus, finite canonical-model chambers,
and finitely many \emph{chosen} terminal models on those chambers,
following \cite[Theorem~5.15]{HX}.
The zero-dimensional case uses $\mathsf M_d$ and uniqueness of the
canonical model; it does not yet assert finiteness of all weak models.

At a central point choose $X\dashrightarrow X_m\to X_c$ over
$S$. The transport invariant in $\mathsf M_d$ supplies the traces
of the whole nearby polytope on $X_m$. Shrink so the finitely many
strict exceptional inequalities remain strict. Add a common pullback
from $S$ so that the central canonical class on $X_c$ is K\"ahler
absolutely. On the boundary of the parameter polytope, apply the
smaller-parameter induction \emph{over $X_c$}; its existence input
is $\mathsf M_d(X_c)$, just proved independently. No projectivity of
$X_m\to X_c$ or Hodge projection over the possibly singular $X_c$
is used.

The central adjoint is zero over $X_c$, so the relative effective
locus is the radial cone on its boundary effective locus. Apply the
cone-length estimate of \cite[Claim~5.1]{HX} on the finitely chosen
\emph{gklt terminal models} $h_i:Y_i\to X_c$, all of dimension $d$.
Write $A_{u,i}$ for the boundary-parameter adjoint, nef over $X_c$,
and $\psi_i:Y_i\to Z_i$ for its relative canonical morphism.
The cone theorem is not applied to the possibly lower-dimensional
canonical targets $Z_i$.

Let $H_i=h_i^*\omega_c$ and choose $\delta>0$ below the
$\omega_c$-degree of every compact integral curve in $X_c$. Fix $\bar\lambda<\delta/(\delta+2d)$ and for
$0<\lambda\leq\bar\lambda$ put
\[
 A_{\lambda,i}=(1-\lambda)H_i+\lambda A_{u,i}.
\]
The same cone estimate shows more than nefness: for a sufficiently
small $\epsilon>0$, uniform on the finite list and the chosen radial
interval, $A_{\lambda,i}-\epsilon H_i$ is nef. Indeed on a
horizontal $A_{u,i}$-negative rational ray its degree is greater than
$(1-\lambda-\epsilon)\delta-2d\lambda>0$; on the nonnegative
part of the cone and on the vertical cone it is nonnegative.
It follows for every class $\xi\in\overline{\mathrm{NA}}(Y_i)$ that
\[
 A_{\lambda,i}\cdot\xi=0
 \quad\Longleftrightarrow\quad
 H_i\cdot\xi=A_{u,i}\cdot\xi=0.
\]
In particular its null curves are exactly those contracted by
$\psi_i$. The central data are crepant through these maps over $X_c$;
thus the convex data defining $A_{\lambda,i}$ remain gklt with
globally modified-big boundary. Apply $\mathsf B_d$ to $A_{\lambda,i}$. Its connected
semiampleness map and $\psi_i$ both have Moishezon fibres, and each
connected such fibre is connected by chains of compact curves.
Equality of their contracted curves makes each map constant on the
fibres of the other. The analytic factorization lemma identifies
the two maps. Thus the selected relative canonical model is also
the canonical model over $S$ on this radial interval.
This proves local polyhedrality and the finite local list; compactness
gives the global finite list of canonical models and chosen terminal
models.

\paragraph{Uniform preparation preserving every weak model.}
\label{prog:uniform-preparation}
Write $A_u=[K_X+B_u+\mathbf M_{u,X}]$ for the adjoints in the
parameter polytope. We first record why a weak model of these gklt data
extracts no divisor. On a common resolution $p:U\to X$, $q:U\to Y$
put $D=p^*A_u-q^*A_{u,Y}$, using the actual structure-boundary
comparison. The weak-model discrepancy inequalities give $p_*D\geq0$,
and $-D$ is $p$-nef because $A_{u,Y}$ is nef over the fixed base.
Negativity therefore gives $D\geq0$. An extracted prime on $Y$ has
coefficient one in the weak-model boundary, whereas its coefficient in
the structure boundary of the gklt source is less than one; it would
give a negative coefficient of $D$. Hence there is no such prime.
Consequently $D$ is $q$-exceptional and the target is itself gklt.

Fix a parameter $u_0$. The boundary plus nef part is globally
modified big by hypothesis. The modified-big preparation
\cite[Lemma~2.23]{HX} gives, on a fixed projective smooth carrier
$a:W\to X$, an effective klt boundary $\Gamma$, a K\"ahler class
$\gamma_0$, and an effective $a$-exceptional real divisor $F$ such that
\begin{equation}\label{prog:eq:uniform-preparation}
             [K_W+\Gamma]+\gamma_0=a^*A_{u_0}+[F].
\end{equation}
Here is the effectivity detail in this use of the preparation. Its new
nef datum dominates a K\"ahler form on a carrier. After resolving the
new structure boundary, choose an effective exceptional divisor whose
negative is relatively ample and subtract a sufficiently small multiple
from the pulled-back datum to make it K\"ahler. Add the same multiple
to the structure boundary,
then replace its negative coefficients by zero and, if needed, add
small positive coefficients on the remaining exceptional primes.
All coefficients stay below one. Only exceptional coefficients have
been increased; their increase is the divisor $F\geq0$ in
\eqref{prog:eq:uniform-preparation}.

On a sufficiently small closed polyhedral neighbourhood of $u_0$ put
\[
                \gamma_u=\gamma_0+a^*(A_u-A_{u_0}).
\]
These classes remain K\"ahler by openness, and the same $\Gamma$ and
the same $F$ satisfy
\[
                  [K_W+\Gamma]+\gamma_u=a^*A_u+[F].
\]
Thus this is one affine family of gklt pairs on one smooth carrier,
with a fixed effective boundary and K\"ahler nef data.

Let $\psi:X\dashrightarrow Y$ be \emph{any} weak model at a parameter
$u$ in this neighbourhood, and resolve the induced map $W\dashrightarrow
Y$ by $r:T\to W$, $q:T\to Y$. Its original comparison is
\[
              (ar)^*A_u=q^*A_{u,Y}+E_u,
              \qquad E_u\geq0\text{ is }q\text{-exceptional}.
\]
Since $\psi$ is nonextracting, every $a$-exceptional prime is exceptional
over $Y$. Hence $r^*F$ is also $q$-exceptional, and
\begin{equation}\label{prog:eq:prepared-weak}
 r^*\bigl([K_W+\Gamma]+\gamma_u\bigr)
       =q^*A_{u,Y}+E_u+r^*F.
\end{equation}
This proves that $W\dashrightarrow Y$ is a weak model of the prepared
pair. In detail, subtract the effective exceptional error on the right
from its structure boundary on $T$ and push down to $Y$. The resulting
boundary is precisely the transform of $\Gamma$, it is effective, the
adjoint is the already existing Bott--Chern class $A_{u,Y}$, and all
discrepancies are at least those of the prepared gklt pair. This uses generalized data at the level of currents. On $T$ keep
the nef datum $r^*\gamma_u$, and subtract the displayed effective
$q$-exceptional error from the prepared structure boundary. Pushing
forward by $q$ gives the boundary $\Gamma_Y$ and the trace of this
fixed nef b-class. Their sum with the canonical current represents
the already existing locally exact adjoint $A_{u,Y}$; the resolution
identity defines its structure boundary. The generalized discrepancies
have not decreased. This argument neither requires $K_Y+\Gamma_Y$ to
be real Cartier nor asserts that the nef trace is separately a
Bott--Chern class. It therefore applies to every K\"ahler weak target,
including the adjunction strata used in special termination.

\paragraph{Every weak model occurs in the finite canonical list.}
Choose a small box in $H^{1,1}(W,\mathbb R)$ which spans that vector
space and whose addition to the compact family $\{\gamma_u\}$ stays
inside the K\"ahler cone. For any $Y$ above choose a K\"ahler class
$\eta_Y$ and define $\eta_W=r_*q^*\eta_Y$. Smooth-source cohomology
puts $\eta_W$ in $H^{1,1}(W,\mathbb R)$ and gives an actual
$r$-exceptional real divisor $D_\eta$ with
\[
                    r^*\eta_W-q^*\eta_Y=[D_\eta].
\]
Its negative is $r$-nef, so $D_\eta\geq0$ by exceptional negativity.
As $W\dashrightarrow Y$ extracts no divisors, $D_\eta$ is also
$q$-exceptional. Scale $\eta_Y$ by a sufficiently small positive number
so that $\eta_W$ lies in the fixed box. Adding this identity to
\eqref{prog:eq:prepared-weak} yields an effective
$q$-exceptional comparison with target class $A_{u,Y}+\eta_Y$, which
is K\"ahler over the original base. Thus $W\dashrightarrow Y$ is the
canonical model of a member of this one enlarged prepared polytope.
Its canonical-model list is finite by the first part of the geography
proof. The induced list of maps from $X$ is therefore finite as well.
A finite cover of the original compact parameter polytope completes
the proof for all weak models.

\end{proof}


\begin{theorem}[The restricted package in every dimension]
\label{prog:restricted-package}
The assertions $\mathsf B_d$, $\mathsf C_d$, $\mathsf M_d$ and
$\mathsf F_d$ hold in every finite dimension, for the globally nef
carrier data and globally modified-big boundary traces specified above.
All birational steps used to prove them are detected steps or belong
to an already-projective relative construction with the stated
relative-line algebraicity. A general semiample morphism is asserted
to be Moishezon; projectivity is asserted only when the construction
provides an actual relatively ample line.
\end{theorem}
\begin{proof}
The assertions in dimensions zero and one follow from degrees and
factorization of maps of compact curves. Suppose the package holds
in dimensions less than $d$. Lemma~\ref{prog:base-point-free-induction}
first proves $\mathsf C_{d,\mathrm{big}}$ and then $\mathsf B_d$.
The former uses $\mathsf C_{d-1}$; the nonbig part of the latter
uses only a chosen $\mathsf M_{d-1}$ model and the explicit
negative-part comparison. In particular, no lower program is lifted.

With $\mathsf B_d$ available, Proposition~\ref{prog:ordinary-step}
constructs detected negative steps in dimension $d$.
Lemma~\ref{prog:big-good-model}, using the already known
$\mathsf F_{d-1}$, gives the big-shift case of $\mathsf M_d$.
Lemma~\ref{prog:relative-model} gives its remaining case, using
$\mathsf M_{d-1}$ and the projective graph construction. This proves
$\mathsf M_d$ over every proper compact K\"ahler base, with its
chosen-model invariant. Lemma~\ref{prog:finite-geography}, whose
hypotheses are \emph{both} $\mathsf B_d$ and $\mathsf M_d$, now
proves $\mathsf F_d$.

It remains to prove $\mathsf C_d$ when its nef NQC adjoint $\alpha$
is not big. Take the projective dlt preparation in
\cite[Section~6]{HX}. Its underlying space is klt and globally
strongly $\mathbb Q$-factorial, and
$\alpha-c_1(K_X)$ is a big Bott--Chern class. Choose an NQC
expression with positive coefficients and integral multiples of its
rational degree-two summands. Corollary~\ref{prog:uniform-canonical-parameter}
provides a single number $t$ such that each negative ray of
$K_X+t\alpha$ is $\alpha$-trivial. Choose a K\"ahler scaling class
$\omega$ with $K_X+t\alpha+\omega$ nef. Every selected negative
ray is ordinary $K_X$-negative; the actual rational canonical line
is its detector. Proposition~\ref{prog:ordinary-step} supplies its
projective contraction or flip, and the fixed integral grid persists.
The descended $\alpha$ remains nef by projective descent of nefness.

The driving class is non-pseudo-effective throughout, by the
corollary. Its initial pseudo-effective scaling threshold is positive;
choose a cutoff $\varepsilon>0$ below that threshold. On the fixed
initial carrier the compact segment of globally nef data
\[
                    t\alpha+[\varepsilon,1]\omega
\]
has modified-big total boundary. At its scaling parameter each working
model is a marked weak model of a member of this segment. Each selected
step is negative for the fixed driving adjoint, so its discrepancies
do not decrease and increase at a valuation affected by that step.
An infinite sequence would repeat a marked model by $\mathsf F_d$;
the cumulative nonzero effective comparison for the fixed adjoint
would then contradict equality of that marking. Hence this program
terminates. Non-pseudo-effectivity excludes a nef endpoint, leaving
an $\alpha$-trivial ordinary Mori contraction $f:X'\to Z$.

For clarity, the remaining contraction argument uses only this
already-projective map. Relative vanishing transports the multiplier
ideal defining the non-klt closed subspace. Lemma~\ref{prog:bc-descent}
descends $\alpha$ in Bott--Chern cohomology, and
Lemma~\ref{prog:integral-h2} descends its fixed rational NQC summands;
projective multisections preserve their nonnegative degrees.
Make the adjoint-preserving modified-big replacement of
\cite[Lemma~2.23 and Section~6]{HX}, with a K\"ahler summand on
the source, and restrict the given contraction to the resulting
non-klt closed subspace, taking its Stein factor. If this non-klt
subspace dominates $Z$, relative Nadel vanishing gives
$\mathcal O_Z=f_*\mathcal O_{\operatorname{Nklt}}$, and its given
Moishezon contraction factors to the required contraction on $Z$.
If it does not dominate, apply the projective canonical bundle formula. The multiplier-ideal calculation identifies the
induced non-klt subspace on its lower-dimensional base and its
Moishezon contraction. Thus $\mathsf C_{d-1}$ applies there.
Pullback and factorization through the projective modifications give
the required Moishezon contraction on the original space. These are
relative vanishing, canonical-bundle-formula and closed-subspace
factorization operations; they require no further global program.
This proves $\mathsf C_d$ and completes the induction.
\end{proof}

\begin{corollary}[Finite positive parts of detected scalings]
\label{prog:detected-positive-truncation}
Let a gklt adjoint with globally nef carrier data be run with scaling
of the trace of a K\"ahler class on a fixed smooth carrier. Suppose
that every chosen negative analytic ray has an actual global rational
line detector and that the total boundary trace on each positive
scaling segment is globally modified big. Every segment whose scaling
parameters lie in a fixed interval $[\varepsilon,T]$, with
$\varepsilon>0$, is finite. If the driving adjoint is not
pseudo-effective, the scaling terminates with a Mori fibre space.
\end{corollary}
\begin{proof}
Proposition~\ref{prog:ordinary-step} constructs each step and
Lemma~\ref{prog:detected-cohomology} supplies the forward traces from
the fixed carrier. Each reached model at a scaling parameter is a
marked weak model of that parameter's adjoint. The compact parameter
segment satisfies $\mathsf F_d$, so there are finitely many such
markings. Strict increase for the fixed driving adjoint rules out
repetition. If that adjoint is not pseudo-effective, the positive
pseudo-effective scaling threshold gives a fixed positive cutoff
containing every parameter of a continuing program. The finite
program cannot end nef, hence ends with the claimed Mori map.
\end{proof}

